\chapter{Rings, fields, polynomials and characters}\label{ch:rings}

The paper is written for readers who handle rings, finite fields and quadratic characters without
thinking about them. This chapter makes that background explicit, with complete proofs, and organises it
around the seven places where the paper actually uses it: the $\ell$-adic valuations of rational numbers
in Definition def:terms and the table of Section 1; the quadratic characters $\chi_{-3}$, $\chi_{-4}$,
$\chi_{q^*}$ and $\chi_S$ as Dirichlet characters; quadratic reciprocity, which makes $\chi_{q^*}$ the
quadratic character of conductor $q$; the $\ell$-th powers in $\F_q^\times$ and the class $c$ of Lemma
lem:tamfactor; the ring theory behind Lemma lem:crit ($\T/(I+J)$ is cyclic, and its order is divisible by
the residue characteristic of every maximal ideal containing $I+J$); the congruence $a^2\equiv1\pmod{\ell^k}$
behind Lemma lem:h0; and the foundations (ideals, unique factorisation, Gauss's lemma, Eisenstein's
criterion, finite fields, cyclotomic polynomials, Gauss sums) that Chapters 3, 5, 6 and 7 build on. The
group theory it needs (cyclic groups, finite abelian groups, finite subgroups of $K^\times$) is in
\Cref{ch:groups} and is used, not reproved.

\section*{Core results at a glance}

Rings are commutative with $1$. The letters $p$, $q$, $\ell$ denote primes, except in
\Cref{sec:rings:finite-fields,sec:rings:powers}, where $q=p^r$ is the size of a finite field; in the paper
$q$ is a prime, so $\F_q=\Z/q\Z$ there. Quadratic characters are named by their discriminant:
$\chi_D=\bigl(\frac D\cdot\bigr)$ for a fundamental discriminant $D$ (for other discriminants the Kronecker symbol
$\bigl(\frac D\cdot\bigr)$ is written out), so $\chi_{q^*}$ (never $\chi_q$) is the character of conductor $q$, and
$\chi_\ell$ is the $\ell$-adic cyclotomic character.

\begin{enumerate}[label=\textbf{\thechapter.\arabic*}]
\item \textbf{Conventions.} Rings, homomorphisms and subrings are unital; $p,q,\ell$ are primes ($q$ a prime power in
Sections 6 and 7); $v_p(0)=+\infty$, $\abs{x}_p=p^{-v_p(x)}$, and for $c\in\Z/\ell^m$, $v_\ell(c)=\min(v_\ell(\tilde c),m)$; the Kronecker
symbol at $2$ and at $-1$; Dirichlet characters and the names $\chi_D$ ($D$ fundamental), $\chi_{q^*}$, $\omega_\ell$, $\chi_0$, $\chi^{(p)}$;
$\zeta_n=e^{2\pi i/n}$ and
$\sqrt{-3}=1+2\zeta_3$. \Cref{conv:rings:rings,conv:rings:primes,conv:rings:valuations,conv:rings:symbols,conv:rings:characters,conv:rings:roots}.
\item \textbf{Rings} (definition: ring, unit group $R^\times$, zero divisor, nilpotent, field, integral domain, subring,
homomorphism, kernel, image). \Cref{def:rings:ring}.
\item \textbf{Characteristic.} For every ring $R$ there is a unique homomorphism $\Z\to R$; its kernel is $n\Z$ for a unique
$n\ge0$, the characteristic. A domain has characteristic $0$ or a prime, and a finite domain is a field.
\emph{(proved here)} \Cref{prop:rings:characteristic}.
\item \textbf{Ideals} (definition: ideal, the ideal generated by a set, $I+J$, $IJ$, $I\cap J$, the annihilator $\Ann_R(S)$,
comaximal ideals, the quotient ring $R/I$). \Cref{def:rings:ideal}.
\item \textbf{Ideals and kernels.} Kernels of homomorphisms are exactly the ideals; $IJ\subseteq I\cap J$; every ideal of $\Z$ is
$n\Z$ for a unique $n\ge0$; $m\Z+n\Z=\gcd(m,n)\Z$ and $m\Z\cap n\Z=\lcm(m,n)\Z$. \emph{(proved here)} \Cref{prop:rings:ideals}.
\item \textbf{Isomorphism theorems.} (i) $R/\ker\phi\simeq\phi(R)$. (ii) Ideals of $R/I$ correspond to ideals $J\supseteq I$ of $R$,
with $(R/I)/(J/I)\simeq R/J$. (iii) If $\phi\colon R\to S$ is surjective with kernel $I$ and $J$ is an ideal of $R$, then $\phi(J)$ is
an ideal and $R/(I+J)\simeq S/\phi(J)$. (iv) Every quotient ring of $\Z$ is $\Z/r\Z$, cyclic as a group. \emph{(proved here)}
\Cref{thm:rings:isomorphism}.
\item \textbf{Two ideals.} The sequence $0\to R/(I\cap J)\to R/I\times R/J\to R/(I+J)\to0$, $x\mapsto(x,x)$, $(y,z)\mapsto y-z$, is
exact. If $I+J=R$ then $I\cap J=IJ$ and $R/IJ\simeq R/I\times R/J$; if $I_1,\dots,I_k$ are pairwise comaximal,
$R/I_1\cdots I_k\simeq\prod R/I_i$. \emph{(proved here)} \Cref{prop:rings:two-ideals}.
\item \textbf{Idempotents.} If $e^2=e$ then $R\simeq Re\times R(1-e)$; $R$ is a product of two nonzero rings iff it has an
idempotent $e\ne0,1$. The ideals of $R_1\times R_2$ are the $I_1\times I_2$, its prime ideals the $P_1\times R_2$ and $R_1\times P_2$.
\emph{(proved here)} \Cref{prop:rings:idempotents}.
\item \textbf{Frobenius in characteristic $p$.} $(x+y)^{p^k}=x^{p^k}+y^{p^k}$, so $x\mapsto x^p$ is a ring endomorphism; it is
injective on a domain and bijective on a finite field; $a^p\equiv a\pmod p$. \emph{(proved here)} \Cref{prop:rings:freshman}.
\item \textbf{Prime and maximal ideals.} $R/P$ is a domain iff $P$ is prime, $R/M$ a field iff $M$ is maximal; maximal ideals are
prime; preimages of primes are prime; the kernel of a map onto a field (or to a field with finite image) is maximal; if $R/P$ is
finite, $P$ is prime iff maximal; the primes of $\Z$ are $0$ and $p\Z$; every non-unit lies in a maximal ideal (Zorn).
\emph{(proved here)} \Cref{thm:rings:prime-maximal}.
\item \textbf{Residue characteristic.} If $\mathfrak m$ is maximal with $R/\mathfrak m$ of characteristic $\ell>0$ and $\mathfrak a\subseteq\mathfrak m$ is an ideal with
$R/\mathfrak a$ finite, then $\ell\mid\#R/\mathfrak a$; a homomorphism $\Z/r\to k$ to a field of characteristic $\ell$ exists iff $\ell\mid r$.
\emph{(proved here)} \Cref{prop:rings:residue-char}.
\item \textbf{Field of fractions.} A domain $R$ embeds in a field $\Frac(R)$, and every injective homomorphism from $R$ to a field
extends uniquely to $\Frac(R)$. \emph{(proved here)} \Cref{thm:rings:fractions}.
\item \textbf{The ring $\Z_{(p)}$.} $\Z_{(p)}=\{a/b:p\nmid b\}$ has units $\{a/b:p\nmid ab\}$, nonzero ideals $p^k\Z_{(p)}$, a unique maximal
ideal $p\Z_{(p)}$, and a surjective reduction map $\Z_{(p)}\to\F_p$ with kernel $p\Z_{(p)}$; $\Z_{(p)}=\{x:v_p(x)\ge0\}$ and
$\Z=\bigcap_p\Z_{(p)}$. \emph{(proved here)} \Cref{prop:rings:zp-local}.
\item \textbf{Divisibility} (definition: associates, irreducible and prime elements, Euclidean domain, PID, UFD, gcd).
\Cref{def:rings:factorisation}.
\item \textbf{Primes and irreducibles.} In a domain a prime element is irreducible, and $(a)$ is a nonzero prime ideal iff $a$
is a prime element; a factorisation into prime elements is unique up to order and units. \emph{(proved here)}
\Cref{lem:rings:prime-irreducible}.
\item \textbf{Euclid, PID, UFD.} A Euclidean domain is a PID; a PID is Noetherian and a UFD, with Bezout's identity and
nonzero primes maximal; in a UFD irreducibles are prime. $\Z$ and $F[x]$ are Euclidean; $\Z[x]$ and $F[x,y]$ are UFDs, not PIDs.
\emph{(proved here)} \Cref{thm:rings:euclid-pid-ufd}.
\item \textbf{Arithmetic of $\Z$.} Bezout, Euclid's lemma, $n=\pm\prod p^{v_p(n)}$, $\gcd$ and $\lcm$ by $\min$ and $\max$ of exponents.
\emph{(proved here)} \Cref{prop:rings:integers}.
\item \textbf{$\Z[i]$ and $\Z[\zeta_3]$.} Both are Euclidean for $N(z)=\abs z^2$ (remainders of norm at most $\frac12N(\beta)$, resp.
$\frac13N(\beta)$), hence PIDs and UFDs; $\Z[i]^\times=\mu_4$, $\Z[\zeta_3]^\times=\mu_6$; elements of prime norm are prime.
\emph{(proved here)} \Cref{thm:rings:gauss-eisenstein-integers}.
\item \textbf{Polynomials.} Over a domain $\deg fg=\deg f+\deg g$ and $R[x]^\times=R^\times$; division by a monic polynomial over any
ring; $F[x]$ is Euclidean, its maximal ideals are the $(f)$, $f$ irreducible; factor theorem; at most $d$ roots in a domain.
\emph{(proved here)} \Cref{prop:rings:poly-division}.
\item \textbf{Repeated roots.} $a$ is a multiple root of $f$ iff $f(a)=f'(a)=0$; $f$ has no repeated root in any extension iff
$\gcd(f,f')=1$; $x^{p^r}-x$ and, when $\operatorname{char}F\nmid n$, $x^n-1$ are separable. \emph{(proved here)} \Cref{prop:rings:derivative}.
\item \textbf{Gauss's lemma.} Over a UFD $R$ with fraction field $K$: contents multiply; a factorisation in $K[x]$ descends to
$R[x]$; a primitive polynomial is irreducible in $R[x]$ iff in $K[x]$; $R[x]$ is a UFD; rational root test.
\emph{(proved here)} \Cref{thm:rings:gauss-lemma}.
\item \textbf{Eisenstein's criterion.} If $a_n\notin P$, $a_i\in P$ for $i<n$, $a_0\notin P^2$ ($P$ prime), then $f$ is not a product of two
polynomials of positive degree; over a UFD a primitive such $f$ is irreducible in $R[x]$ and $K[x]$. \emph{(proved here)}
\Cref{thm:rings:eisenstein}.
\item \textbf{Reduction modulo $p$.} If $p$ does not divide the leading coefficient of a primitive $f\in\Z[x]$ and $\bar f$ is
irreducible in $\F_p[x]$, then $f$ is irreducible in $\Q[x]$; the degrees of a rational factorisation are sums of degrees of
the factors of $\bar f$. \emph{(proved here)} \Cref{prop:rings:reduction-mod-p}.
\item \textbf{Kronecker's construction.} For $f\in F[x]$ irreducible of degree $d$, $F[x]/(f)$ is a field of degree $d$ in which $\bar x$
is a root of $f$; $F[x]/(f)\simeq F[\alpha]$ for any root $\alpha$; every polynomial splits in a finite extension. \emph{(proved here)}
\Cref{thm:rings:adjoin-root}.
\item \textbf{The ring $\Z/n\Z$.} $a$ is a unit iff $\gcd(a,n)=1$; $\Z/n$ is a field iff $n$ is prime; $\#(\Z/n)^\times=\varphi(n)$,
$a^{\varphi(n)}\equiv1$, $\varphi(p^k)=p^{k-1}(p-1)$, $\sum_{d\mid n}\varphi(d)=n$; $(\Z/n)^\times=\Aut(\Z/n)$. \emph{(proved here)} \Cref{prop:rings:zn}.
\item \textbf{Chinese remainder theorem.} $\Z/mn\simeq\Z/m\times\Z/n$ iff $\gcd(m,n)=1$; $\Z/n\simeq\prod\Z/p_i^{k_i}$,
$(\Z/n)^\times\simeq\prod(\Z/p_i^{k_i})^\times$, $\varphi$ is multiplicative. \emph{(proved here)} \Cref{thm:rings:crt}.
\item \textbf{Lifting the exponent.} For odd $\ell$, $\ell\mid a-b$, $\ell\nmid ab$: $v_\ell(a^n-b^n)=v_\ell(a-b)+v_\ell(n)$; for odd $a$ and even
$n$, $v_2(a^n-1)=v_2(a-1)+v_2(a+1)+v_2(n)-1$; $1+\ell$ has order $\ell^{k-1}$ modulo $\ell^k$; $v_\ell(q^f-1)$ for $q\ne\ell$;
$v_\ell(q^2-1)=v_\ell(q-1)+v_\ell(q+1)$. \emph{(proved here)} \Cref{lem:rings:lte}.
\item \textbf{Primitive roots.} $(\Z/p^k)^\times$ is cyclic for odd $p$; $(\Z/2^k)^\times=\{\pm1\}\times\langle5\rangle$ for $k\ge3$;
$(\Z/n)^\times$ is cyclic iff $n\in\{1,2,4,p^k,2p^k\}$, $p$ odd. \emph{(proved here)} \Cref{thm:rings:primitive-roots}.
\item \textbf{Squares congruent to one.} For odd $\ell$ and $k\ge1$, $a^2\equiv1\pmod{\ell^k}$ iff $a\equiv\pm1$; every unit squares to $1$ modulo $\ell^k$
iff $\ell^k=3$; a homomorphism onto $(\Z/\ell^k)^\times$ (or one whose image contains $2$) with $\chi^2\equiv1$ forces
$\ell^k=3$. \emph{(proved here)} \Cref{cor:rings:squares-one}.
\item \textbf{Finite fields.} A finite field has prime characteristic $p$ and $p^r$ elements; $F^\times$ is cyclic;
$x^{p^r}-x=\prod_{a\in F}(x-a)$; $F=\F_p[\alpha]$ for a generator $\alpha$ of $F^\times$. \emph{(proved here)}
\Cref{prop:rings:finite-field-basic}.
\item \textbf{Existence and uniqueness.} For every prime power $q=p^r$ there is a field with $q$ elements, unique up to
isomorphism, isomorphic to $\F_p[x]/(f)$ for every irreducible $f$ of degree $r$. \emph{(proved here)}
\Cref{thm:rings:finite-field-exists}.
\item \textbf{Frobenius and subfields.} $\Frob_p\colon x\mapsto x^p$ is an automorphism of $\F_{p^r}$ of order $r$ with fixed field $\F_p$;
$\F_{p^r}$ has a subfield with $p^d$ elements iff $d\mid r$, and then exactly one; the fixed field of $\Frob_p^s$ is $\F_{p^{\gcd(s,r)}}$.
\emph{(proved here)} \Cref{prop:rings:frobenius}.
\item \textbf{Irreducible polynomials over $\F_q$.} $x^{q^r}-x$ is the product of the monic irreducibles of degree dividing $r$;
$q^r=\sum_{d\mid r}dN_d$ and $N_r=\frac1r\sum_{d\mid r}\mu(d)q^{r/d}>0$. \emph{(proved here)} \Cref{prop:rings:irreducible-count}.
\item \textbf{Roots of unity in $\F_q$.} $\mu_n(\F_q)$ is cyclic of order $\gcd(n,q-1)$; the least $f$ with $\mu_n\subset\F_{q^f}$ is the order of $q$
modulo $n$; no nontrivial $p$-power roots of unity; $-1$ is a square iff $q\equiv1\pmod4$, $-3$ (for $p\ne2,3$) iff $q\equiv1\pmod3$; for
$\sigma$ of order $3$ (resp. $4$) and $m$ prime to $3$ (resp. odd), $\sigma^m=\sigma^{\pm1}$ according to $m$ modulo $3$ (resp. $4$).
\emph{(proved here)} \Cref{prop:rings:mu-in-fq}.
\item \textbf{Powers in $\F_q^\times$.} For $d=\gcd(m,q-1)$, $(\F_q^\times)^m=(\F_q^\times)^d$ has index $d$; $u$ is an $m$-th power iff
$u^{(q-1)/d}=1$; $[\F_q^\times:(\F_q^\times)^{\ell^k}]=\ell^{\min(k,v_\ell(q-1))}$. \emph{(proved here)} \Cref{thm:rings:power-index}.
\item \textbf{The class $c(u)$.} For $\ell\nmid q$ and $v=v_\ell(q-1)$, $\F_q^\times/(\F_q^\times)^{\ell^v}$ is cyclic of order $\ell^v$; a generator identifies it
with $\Z/\ell^v$, and $v_\ell(c(u))$ does not depend on the generator; for $k\le v$, $u$ is an $\ell^k$-th power iff $v_\ell(c(u))\ge k$;
$\min(v_\ell(n),v,v_\ell(c(u)))\ge1$ iff $\ell\mid n$, $q\equiv1\pmod\ell$ and $u\in(\F_q^\times)^\ell$. \emph{(proved here)} \Cref{prop:rings:class-c}.
\item \textbf{Roots of unity} (definition: $\mu_n(K)$, $\mu(K)$, primitive $n$-th root of unity). \Cref{def:rings:roots-of-unity}.
\item \textbf{Roots of unity in fields.} $\mu_n(K)$ is cyclic of order dividing $n$, and of order $n$ only if $\operatorname{char}K\nmid n$;
$\mu_n(\C)=\langle\zeta_n\rangle$; $\mu(\Q)=\{\pm1\}$, $\mu(\Q(i))=\mu_4$, $\mu(\Q(\zeta_3))=\mu_6$. \emph{(proved here)} \Cref{prop:rings:roots-of-unity}.
\item \textbf{Cyclotomic polynomials.} $x^n-1=\prod_{d\mid n}\Phi_d$; $\Phi_n\in\Z[x]$ is monic of degree $\varphi(n)$;
$\Phi_n=\prod_{d\mid n}(x^d-1)^{\mu(n/d)}$; formulas for $\Phi_{p^k}$, $\Phi_{2m}$, $\Phi_n(1)$, $\Phi_n(0)$; $\Phi_{p^k}$ is irreducible over $\Q$.
\emph{(proved here)} \Cref{thm:rings:cyclotomic}.
\item \textbf{Cyclotomic polynomials modulo $p$.} For $p\nmid n$ the roots of $\Phi_n$ in characteristic $p$ are the elements of order $n$;
$p\mid\Phi_n(a)$ iff $a$ has order $n$ modulo $p$, so $p\equiv1\pmod n$; there are infinitely many primes $\equiv1\pmod n$; $\Phi_n$
splits modulo $p$ iff $p\equiv1\pmod n$. \emph{(proved here)} \Cref{prop:rings:cyclotomic-mod-p}.
\item \textbf{Legendre symbol} (definition). \Cref{def:rings:legendre}.
\item \textbf{Euler's criterion.} $\bigl(\frac ap\bigr)\equiv a^{(p-1)/2}\pmod p$; the symbol is the unique character of order $2$ of $(\Z/p)^\times$;
$\bigl(\frac{-1}p\bigr)=(-1)^{(p-1)/2}$; $\bigl(\frac a3\bigr)\equiv a\pmod3$. \emph{(proved here)} \Cref{prop:rings:legendre}.
\item \textbf{The quadratic Gauss sum.} In a field of characteristic $\ne q$ with a primitive $q$-th root of unity $\zeta$,
$g=\sum\bigl(\frac aq\bigr)\zeta^a$ satisfies $g^2=q^*=(-1)^{(q-1)/2}q$ and $\sum\bigl(\frac aq\bigr)\zeta^{ab}=\bigl(\frac bq\bigr)g$; so $\sqrt{q^*}\in\Q(\zeta_q)$.
\emph{(proved here; the sign of $g$ over $\C$ is stated in \Cref{rem:rings:gauss-sign})} \Cref{prop:rings:quadratic-gauss-sum}.
\item \textbf{Quadratic reciprocity.} For odd primes $p\ne q$: $\bigl(\frac pq\bigr)\bigl(\frac qp\bigr)=(-1)^{\frac{p-1}2\frac{q-1}2}$, i.e.
$\bigl(\frac{q^*}p\bigr)=\bigl(\frac pq\bigr)$; $\bigl(\frac{-1}p\bigr)=(-1)^{(p-1)/2}$; $\bigl(\frac2p\bigr)=(-1)^{(p^2-1)/8}$. \emph{(proved here)} \Cref{thm:rings:qr}.
\item \textbf{Jacobi symbol} (definition). \Cref{def:rings:jacobi}.
\item \textbf{Jacobi reciprocity.} The Jacobi symbol is multiplicative in both entries, periodic in the top one, and satisfies
reciprocity and both supplements; $\bigl(\frac an\bigr)=1$ does not make $a$ a square modulo $n$. \emph{(proved here)} \Cref{prop:rings:jacobi}.
\item \textbf{Kronecker symbol, discriminants} (definition: the Kronecker symbol $\bigl(\frac an\bigr)$ for all $a,n\in\Z$; discriminant;
fundamental and prime discriminants; $\chi_D=\bigl(\frac D\cdot\bigr)$ for fundamental $D$). \Cref{def:rings:kronecker}.
\item \textbf{Kronecker characters.} For a discriminant $D$, $n\mapsto\bigl(\frac Dn\bigr)$ is completely multiplicative, vanishes iff
$\gcd(D,n)>1$, and is a real Dirichlet character modulo $\abs D$ with value $\sgn D$ at $-1$; values at primes;
$\bigl(\frac{DD'}\cdot\bigr)=\bigl(\frac D\cdot\bigr)\bigl(\frac{D'}\cdot\bigr)$ for discriminants $D,D'$, and $\bigl(\frac{Df^2}n\bigr)=\bigl(\frac Dn\bigr)$ for $\gcd(n,f)=1$.
\emph{(proved here)} \Cref{prop:rings:kronecker}.
\item \textbf{The character $\chi_{q^*}$.} $q^*\equiv1\pmod4$, and $\bigl(\frac{q^*}n\bigr)=\bigl(\frac nq\bigr)$ for every $n\in\Z$; $\chi_{q^*}$ is the unique
character of order $2$ and conductor $q$; the same for $\ell^*$; $3^*=-3$, $5^*=5$, $7^*=-7$, $11^*=-11$, $13^*=13$. \emph{(proved here)}
\Cref{prop:rings:qstar}.
\item \textbf{Dirichlet characters} (definition: character modulo $n$, principal character, parity, the group $X(n)$,
characters with values in another field, $\omega_\ell$). \Cref{def:rings:dirichlet}.
\item \textbf{Basic properties.} Complete multiplicativity, periodicity, values in $\mu_{\varphi(n)}$, $\bar\chi=\chi^{-1}$, $\chi(-1)=\pm1$;
$X(n)\simeq(\Z/n)^\times$ non-canonically. \emph{(proved here)} \Cref{prop:rings:dirichlet}.
\item \textbf{Orthogonality.} $\sum_a\chi(a)=\varphi(n)[\chi=\chi_0]$, $\sum_\chi\chi(a)=\varphi(n)[a\equiv1]$, where the Iverson bracket $[P]$ is $1$ if the statement $P$ holds and $0$ otherwise; the characters are an orthonormal basis
of the functions on $(\Z/n)^\times$. \emph{(proved here)} \Cref{thm:rings:orthogonality}.
\item \textbf{Conductor} (definition: $\chi$ defined modulo $d$, conductor $f_\chi$, primitive character). \Cref{def:rings:conductor}.
\item \textbf{Conductors.} If $\chi$ is defined modulo $d_1$ and $d_2$, it is defined modulo $\gcd(d_1,d_2)$; the conductor divides every
such modulus and $\chi$ is induced by a unique primitive $\chi^{\mathrm{prim}}$; $f_\chi=\prod_pf_{\chi^{(p)}}$ for the $p$-components $\chi^{(p)}$; products of primitive characters
of coprime conductors are primitive; $\chi_{-4}$, $\chi_{\pm8}$ have conductor $4$, $8$, and $\chi_S$ has conductor $\prod_{q\in S}q$.
\emph{(proved here)} \Cref{prop:rings:conductor}.
\item \textbf{Fundamental discriminants.} A fundamental discriminant $D\ne1$ is uniquely a product of coprime prime discriminants
$-4,8,-8,q^*$; $\chi_D$ is primitive of conductor $\abs D$. \emph{(proved here)} \Cref{cor:rings:fundamental-primitive}.
\item \textbf{Gauss sums.} $\tau(\chi,b)=\bar\chi(b)\tau(\chi)$ for $b$ prime to $n$, and for all $b$ if $\chi$ is primitive; then $\abs{\tau(\chi)}^2=n$
and $\tau(\chi)\tau(\bar\chi)=\chi(-1)n$; $\tau(\chi_{-4})=2i$, $\tau(\chi_{-3})=\sqrt{-3}$; $\tau(\chi)=\mu(n/f)\chi^{\mathrm{prim}}(n/f)\tau(\chi^{\mathrm{prim}})$ when
$\gcd(n/f,f)=1$. \emph{(proved here)} \Cref{thm:rings:gauss-sums}.
\item \textbf{The paper's characters.} The table of $\chi_{-4}$, $\chi_{-3}$, $\chi_{q^*}$, $\chi_S$, $\chi_{\ell^*}$, $\omega_\ell$ (Kronecker symbol,
conductor, parity, values at primes); $p\equiv\chi_{-3}(p)\pmod3$ for $p\ne3$ and $p\equiv\chi_{-4}(p)\pmod4$ for odd $p$;
$\bigl(\frac\cdot\ell\bigr)=\omega_\ell^{(\ell-1)/2}$, equal to $\omega_\ell$ iff $\ell=3$; $3\mid q-\chi_{-4}(q)$ iff $q\equiv\pm1\pmod{12}$. \emph{(proved here)}
\Cref{prop:rings:paper-characters}.
\item \textbf{Valuations} (definition: $v_p$ on $\Q$, $v_p(0)=+\infty$). \Cref{def:rings:valuation}.
\item \textbf{Properties of $v_p$.} $v_p(xy)=v_p(x)+v_p(y)$; $v_p(x+y)\ge\min$, with equality if the valuations differ; $v_p\ge0$ on
$\Z_{(p)}$ exactly; $p^k\mid n$ iff $v_p(n)\ge k$; Legendre's formula for $v_p(n!)$. \emph{(proved here)} \Cref{prop:rings:valuation}.
\item \textbf{The $p$-adic absolute value} (definition: $\abs x_p=p^{-v_p(x)}$). \Cref{def:rings:padic-abs}.
\item \textbf{Absolute values on $\Q$.} $\abs\cdot_p$ is an ultrametric absolute value; the product formula $\prod_{p\le\infty}\abs x_p=1$;
$\Q$ is not complete for $\abs\cdot_p$. \emph{(proved here)} \Cref{prop:rings:padic-abs}.
\item \textbf{Ostrowski's theorem.} Every nontrivial absolute value on $\Q$ is equivalent to $\abs\cdot_\infty$ or to exactly one
$\abs\cdot_p$. \emph{(stated; reference)} \Cref{thm:rings:ostrowski}.
\item \textbf{$\ell$-integrality.} For primes $q\ne\ell$, $q$ is an $\ell$-adic unit, $P(q^{-k})\in\Z_{(\ell)}$ for $P\in\Z[X]$, and
$v_\ell(\frac{q-1}q)=v_\ell(q-1)$, $v_\ell(1\mp q^{-1})=v_\ell(q\mp1)$ for either sign, $v_\ell(1-q^{-2})=v_\ell(q-1)+v_\ell(q+1)$. \emph{(proved here)}
\Cref{lem:rings:ell-integral}.
\item \textbf{The $\kod{I}_0^*$ identity.} If $\alpha\beta=q$, $\alpha+\beta=a$, then
$(1-\frac\alpha{\beta q})(1-\frac1q)(1-\frac\beta{\alpha q})=\frac{q-1}q\cdot\frac{(q+1)^2-a^2}{q^2}$, and $(q+1)^2-a^2\ge(q-1)^2$ if $a^2\le4q$.
\emph{(proved here)} \Cref{prop:rings:i0star-identity}.
\item \textbf{The paper's local terms.} For primes $q\ge5$, odd $\ell\ne q$ and $a\in\Z$ with $a^2\le4q$: (i) $v_\ell(q-\chi_{-3}(q))\ge1$ always at $\ell=3$ and
iff $q\equiv\chi_{-3}(q)\pmod\ell$ otherwise; (ii) $v_\ell(q-\chi_{-4}(q))\ge1$ iff $q\equiv\pm1\pmod{12}$ at $\ell=3$, iff $q\equiv\chi_{-4}(q)\pmod\ell$ otherwise;
(iii) $v_\ell(q^2-1)\ge1$ always at $\ell=3$, iff $q\equiv\pm1\pmod\ell$ otherwise; (iv) $v_\ell(q-1)+v_\ell((q+1)^2-a^2)\ge1$ iff $q\equiv1\pmod\ell$ or
$a\equiv\pm(q+1)\pmod\ell$; (v) for $n\ge1$ and $\bar u\in\F_q^\times$, $v_\ell(n)+\min(v_\ell(n),v_\ell(q-1),v_\ell(c(\bar u)))>v_\ell(n)$ iff $\ell\mid n$, $q\equiv1\pmod\ell$
and $\bar u\in(\F_q^\times)^\ell$. \emph{(proved here)} \Cref{prop:rings:local-terms}.
\end{enumerate}

\section*{Where the paper uses this chapter}

\begin{itemize}
\item \textbf{Definition def:terms and the table of Section 1.} The terms $e_\ell(E,q)$ and $t_\ell(E,q)$ are $\ell$-adic valuations
of integers and of rationals with denominator a power of $q$: \Cref{prop:rings:valuation,lem:rings:ell-integral,prop:rings:i0star-identity,prop:rings:local-terms}
and \Cref{ex:rings:paper-values}; the same results give Remark rem:afterquant(1) ($e_3>0$ at every prime $q\ge5$ of type $\kod{I}_n^*$). The
characters $\chi_{-3}$, $\chi_{-4}$ of the table: \Cref{prop:rings:paper-characters}.
\item \textbf{Lemma lem:euler (the omitted factors at tame primes) and Lemma lem:IV.} "Frobenius centralises $\sigma$ if
$q\equiv1\pmod3$ and inverts it if $q\equiv2$": \Cref{prop:rings:mu-in-fq}(v) and \Cref{prop:rings:paper-characters}(ii); the
identity for type $\kod{I}_0^*$: \Cref{prop:rings:i0star-identity}; $e_\ell=v_\ell(1-q^{-2})=v_\ell(q^2-1)$: \Cref{lem:rings:ell-integral}.
\item \textbf{Proposition prop:forms.} "$e_3(E,q)=v_3(q-\chi_{-3}(q))\ge1$ since $q\equiv\chi_{-3}(q)\pmod3$": \Cref{prop:rings:paper-characters}(i),
\Cref{prop:rings:local-terms}(i).
\item \textbf{Lemma lem:twist, Theorems thm:twist, thm:quant, thm:qual and Example ex:qual.} The character $\chi_{q^*}$ of
conductor $q$, the field $\Q(\sqrt{q^*})$, the character $\chi_S$ of conductor $D=\prod_{q\in S}q$ and the vanishing of
$\chi_S(n)$ for $\gcd(n,D)>1$: \Cref{thm:rings:qr,prop:rings:qstar,prop:rings:conductor,prop:rings:paper-characters}.
\item \textbf{Lemma lem:tamfactor, Remark rem:ko(1) and Section sec:tam.} $\ell$-th powers in $\F_q^\times$, the class $c$ of the
reduction of the unit part $u$, and the reading of $v_\ell(c)$: \Cref{thm:rings:power-index,prop:rings:class-c} and
\Cref{ex:rings:powers}; reduction of units: \Cref{prop:rings:zp-local}(iii).
\item \textbf{Lemma lem:crit (congruence criterion), Lemma lem:hecke and Step 1 of the proof of Theorem thm:quant.}
$\T/(I+J)\simeq\Z/\phi_f(J)$ is cyclic, the cokernel of $\T\subset\T/I\times\T/J$ is $\T/(I+J)$, the decomposition
$\T\otimes\Q=\Q\times\T'$, and "$\ell$ divides $\#\T/(I+J)=r$": \Cref{thm:rings:isomorphism,prop:rings:two-ideals,prop:rings:idempotents,thm:rings:prime-maximal,prop:rings:residue-char}.
\item \textbf{Lemma lem:fpr.} "$L_p(A_f,1)^{-1}$ is a nonzero rational number which is an $\ell$-adic integer":
\Cref{prop:rings:zp-local,lem:rings:ell-integral}.
\item \textbf{Remark rem:wild.} $P_2(X)=1+2X$ gives $e_3=v_3(\frac32)=1$; the cubic factors at $2$: \Cref{ex:rings:paper-values}(e).
\item \textbf{Section sec:anomaly.} The type $\mathrm{III}$ primes with $e_3>0$ are those with $q\equiv\pm1\pmod{12}$:
\Cref{prop:rings:paper-characters}(iv).
\item \textbf{Lemma lem:h0 and Proposition prop:ns.} $\chi_\ell^2\equiv1\pmod{\ell^k}$ on a group mapping onto $\Z_\ell^\times$ forces $k=0$ or $\ell^k=3$:
\Cref{cor:rings:squares-one}; $\bar\chi_3^{\mathrm{cyc}}=\chi_{-3}$ on the Dirichlet side: \Cref{prop:rings:paper-characters}(iii); $\F_9=\F_3[i]$:
\Cref{thm:rings:adjoin-root}; the $3$-division polynomial test: \Cref{prop:rings:reduction-mod-p} and \Cref{ex:rings:psi3}.
\end{itemize}

\section{Rings, ideals, quotients and products}\label{sec:rings:rings}

This section fixes the language of rings and ideals and proves the isomorphism theorems, the exact sequence of two
ideals and the decomposition of a ring by idempotents. In the paper these are exactly the steps of Lemma lem:crit:
the Hecke algebra $\T$ is a ring, $I=\ker\phi_f$ and $J=\Ann_\T(f^\perp)$ are ideals with $I\cap J=0$, the quotient
$\T/(I+J)$ is identified with $\Z/\phi_f(J)$, and $\T\otimes\Q$ splits as $\Q\times\T'$. We begin with the conventions of
the whole chapter.

\begin{convention}[rings]\label{conv:rings:rings}
Rings are commutative with an identity $1$; ring homomorphisms send $1$ to $1$; subrings contain $1$. The zero ring
(in which $1=0$) is allowed, but fields and integral domains are nonzero by definition. $R^\times$ is the group of units
of $R$. The paper's Hecke algebra $\T$ is such a ring. Matrix rings appear only as examples and are then flagged as
non-commutative.
\end{convention}

\begin{convention}[primes and the letter $q$]\label{conv:rings:primes}
The letters $p$, $q$, $\ell$ denote primes unless stated otherwise. In \Cref{sec:rings:finite-fields,sec:rings:powers},
and wherever $\F_q$ is written with $q$ not declared prime, $q=p^r$ is a prime power and $\F_q$ is a field with $q$
elements. In the paper $q$ is always a prime (a bad prime of $E$), so $\F_q=\Z/q\Z$ there and primitive roots modulo
$q$ exist. The coefficient prime $\ell$ is odd in every use of the paper. In one place, \Cref{pl:rings:characters}, $q$ is also
the variable $q=e^{2\pi iz}$ of $q$-expansions (Appendix B, B-C1), and this is said there. No two of these three meanings of $q$
occur in one display.
\end{convention}

\begin{convention}[valuations]\label{conv:rings:valuations}
For a prime $p$, $v_p\colon\Q\to\Z\cup\{+\infty\}$ is the $p$-adic valuation (\Cref{def:rings:valuation}), with
$v_p(0)=+\infty$ and $v_p(p)=1$, and $\abs x_p=p^{-v_p(x)}$, $\abs0_p=0$. For a nonzero ideal $r\Z$, $v_p(r\Z):=v_p(r)$. For a class
$c\in\Z/\ell^m\Z$ put
\[
v_\ell(c):=\min\bigl(v_\ell(\tilde c),m\bigr)\qquad\text{for any lift }\tilde c\in\Z\text{ of }c.
\]
This does not depend on the lift: two lifts differ by a multiple of $\ell^m$, and $\min(v_\ell(x),m)=\min(v_\ell(x+\ell^my),m)$ by
\Cref{prop:rings:valuation}(ii). So $v_\ell(0)=m$, and for $0\le k\le m$, $v_\ell(c)\ge k$ if and only if $c\in\ell^k\Z/\ell^m\Z$. With this
reading the paper's $\min(v_\ell(n),v_\ell(q-1),v_\ell(c))$ in Lemma lem:tamfactor makes sense also when $c=0$, and also when
$v_\ell(q-1)=0$ (then $c$ lies in the zero group $\Z/1$ and the minimum is $0$).
\end{convention}

\begin{convention}[symbols]\label{conv:rings:symbols}
$\bigl(\frac ap\bigr)$ is the Legendre symbol for an odd prime $p$ (\Cref{def:rings:legendre}), $\bigl(\frac an\bigr)$ the Jacobi symbol for an odd
$n>0$ (\Cref{def:rings:jacobi}), and $\bigl(\frac an\bigr)$ the Kronecker symbol for arbitrary $a,n\in\Z$ (\Cref{def:rings:kronecker}). The three agree
where more than one is defined. The Kronecker symbol has $\bigl(\frac a2\bigr)=0$ for even $a$, $1$ for $a\equiv\pm1\pmod8$ and $-1$ for
$a\equiv\pm3\pmod8$; $\bigl(\frac a{-1}\bigr)=1$ for $a\ge0$ and $-1$ for $a<0$; $\bigl(\frac a1\bigr)=1$; and $\bigl(\frac a0\bigr)=1$ if $a=\pm1$, $0$ otherwise.
\end{convention}

\begin{convention}[characters and their names]\label{conv:rings:characters}
A Dirichlet character modulo $n$ is a function $\chi\colon\Z\to\C$ coming from a homomorphism $(\Z/n)^\times\to\C^\times$, with
$\chi(a)=0$ when $\gcd(a,n)>1$ (\Cref{def:rings:dirichlet}). Quadratic characters are named by a fundamental discriminant
(\Cref{def:rings:kronecker}): for such $D$, $\chi_D(a)=\bigl(\frac Da\bigr)$ is (for $D\ne1$) the character of $\Q(\sqrt D)$ (Appendix B, B-C6(b)), primitive of conductor
$\abs D$ (\Cref{cor:rings:fundamental-primitive}); so $\chi_{-3}=\bigl(\frac{-3}\cdot\bigr)$ and $\chi_{-4}=\bigl(\frac{-4}\cdot\bigr)$. For a discriminant $D$ that is not
fundamental the Kronecker symbol $\bigl(\frac D\cdot\bigr)$ is always written out and never called $\chi_D$: for instance $\bigl(\frac{-12}\cdot\bigr)$ is an imprimitive
character modulo $12$ that agrees with $\chi_{-3}$ at odd integers, whereas under B-C6(b) the symbol $\chi_{-12}$ would mean $\chi_{-3}$ itself, of conductor
$3$; the chapter does not use the symbol $\chi_{-12}$. The letter $D$ denotes a discriminant; as in the paper, $D=\prod_{q\in S}q$ is also the conductor of $\chi_S$
(the absolute value of the fundamental discriminant $\prod_{q\in S}q^*$). For an odd prime $q$
put $q^*=(-1)^{(q-1)/2}q$; the quadratic character of conductor $q$ is $\chi_{q^*}$, never $\chi_q$. Thus $\chi_{-7}$, $\chi_{-11}$ are the characters of conductor $7$, $11$, while $\chi_5$, $\chi_{13}$ have conductor $5$, $13$.
A prime-letter subscript ($\chi_\ell$, $\chi_p$) is reserved for the cyclotomic character ($\chi_D$, $\chi_S$, $\chi_{q^*}$ name quadratic
characters): $\chi_\ell$ is the $\ell$-adic cyclotomic character of
\Cref{ch:galois}, $\bar\chi_\ell$ its reduction, written $\chi_3^{\mathrm{cyc}}$, $\bar\chi_3^{\mathrm{cyc}}$ at the prime $3$; and
$\omega_\ell\colon a\mapsto a\bmod\ell\in\F_\ell^\times$ is the Dirichlet character that corresponds to $\bar\chi_\ell$ (\Cref{ch:galois,ch:ant}
make the correspondence). Two further names are not of this kind: $\chi_0$ is the principal character modulo $n$, the standard
name (the numeral $0$ is not a discriminant, so it does not collide with B-C6(b)); and the $p$-component of a character $\chi$, a character
modulo a power of $p$, is written $\chi^{(p)}$ with a superscript (\Cref{prop:rings:conductor}(v)), never $\chi_p$, and members of a family of characters are
written $\chi^{(i)}$ in the same way. The paper (revised version)
follows the same rules.
\end{convention}

\begin{convention}[roots of unity]\label{conv:rings:roots}
$\mu_n(K)=\{x\in K:x^n=1\}$; $\zeta_n=e^{2\pi i/n}\in\C$; $\zeta_3=(-1+\sqrt{-3})/2$ with $\sqrt{-3}:=1+2\zeta_3=i\sqrt3$. A primitive cube
root of unity is always $\zeta_3$, never $\omega$ (the letter $\omega$ is reserved for differentials and for the
characters $\omega_\ell$). In this chapter $\pi$ is the real number $3.14159\ldots$; the element $1-\zeta_3$ is not given a name.
\end{convention}

\begin{definition}\label{def:rings:ring}
Let $R$ be a ring. An element $a\in R$ is a \emph{unit} if $ab=1$ for some $b\in R$; the units form the group $R^\times$. A
\emph{zero divisor} is an $a\ne0$ with $ab=0$ for some $b\ne0$; $a$ is \emph{nilpotent} if $a^n=0$ for some $n\ge1$. $R$ is a
\emph{field} if $R\ne0$ and $R^\times=R\setminus\{0\}$, and an \emph{integral domain} (or \emph{domain}) if $R\ne0$ and $R$ has no zero divisors.
A \emph{subring} is a subset containing $0$ and $1$ and closed under $+$, $-$, $\times$. A \emph{homomorphism} $\phi\colon R\to S$ is a
map with $\phi(a+b)=\phi(a)+\phi(b)$, $\phi(ab)=\phi(a)\phi(b)$, $\phi(1)=1$; its \emph{kernel} is $\phi^{-1}(0)$ and its \emph{image}
$\phi(R)$, a subring of $S$.
\end{definition}

\begin{example}\label{ex:rings:rings}
$\Z\subset\Q\subset\R\subset\C$; $\Z/n\Z$; the polynomial ring $R[x]$ (\Cref{sec:rings:polynomials}); the product $R_1\times R_2$ with
componentwise operations, which has the zero divisors $(1,0)$ and $(0,1)$; the subring $\Z[i]=\{a+bi:a,b\in\Z\}$ of $\C$. If $A$
is an abelian group, $\End_\Z(A)$ is a ring under addition and composition, in general non-commutative; a set of
commuting endomorphisms generates a commutative subring. This is the shape of the Hecke algebra: the paper's $\T$ is
"the subring of $\End_\C S_2(\Gamma_0(N))$ generated over $\Z$ by all $T_n$" (\Cref{ch:hecke}). The element $x$ of $\Z[x]/(x^2)$ is
nilpotent; a domain has no nonzero nilpotents.
\end{example}

\begin{proposition}\label{prop:rings:characteristic}
Let $R$ be a ring.
\begin{enumerate}[label=(\roman*)]
\item There is exactly one homomorphism $\Z\to R$, namely $n\mapsto n\cdot1$. Its kernel is $n\Z$ for a unique $n\ge0$, called the
\emph{characteristic} $\operatorname{char}R$; its image is isomorphic to $\Z/n\Z$.
\item If $R$ is a domain, $\operatorname{char}R$ is $0$ or a prime.
\item A finite domain is a field.
\end{enumerate}
\end{proposition}

\begin{proof}
(i) A homomorphism $\psi\colon\Z\to R$ has $\psi(1)=1$, hence $\psi(n)=n\cdot1$ for $n\ge0$ by additivity and $\psi(-n)=-\psi(n)$;
conversely $n\mapsto n\cdot1$ is a homomorphism by the distributive law. The kernel is an ideal of $\Z$, hence of the form $n\Z$
with a unique $n\ge0$ (\Cref{prop:rings:ideals}(iii) below, whose proof uses only division with remainder), and the image is
isomorphic to $\Z/n\Z$ by \Cref{thm:rings:isomorphism}(i).
(ii) Let $n=\operatorname{char}R>0$. Since $R\ne0$, $n\ne1$. If $n=ab$ with $1<a,b<n$, then $(a\cdot1)(b\cdot1)=0$ with both factors
nonzero (by minimality of $n$), which is impossible in a domain. So $n$ is prime.
(iii) Let $R$ be a finite domain and $a\ne0$. The map $x\mapsto ax$ is injective, because $ax=ay$ gives $a(x-y)=0$, so $x=y$.
An injective map of a finite set to itself is surjective, so $ab=1$ for some $b$.
\end{proof}

For example $\operatorname{char}\Z/6\Z=6$, and indeed $\Z/6\Z$ is not a domain: $2\cdot3=0$, exactly the factorisation used in (ii). The
rings $\F_p[x]$ and $\F_p(x)$ have characteristic $p$ and are infinite; $\F_p[x]$ is a domain that is not a field ($x$ has no inverse, by degrees),
so the finiteness in (iii) cannot be dropped. In the paper the Hecke algebra $\T$ has characteristic $0$ (it contains $\Z\cdot T_1=\Z$), while
its quotients $\T/\mathfrak m$ by maximal ideals containing $\ell$ have characteristic $\ell$.

\begin{definition}\label{def:rings:ideal}
An \emph{ideal} of $R$ is an additive subgroup $I$ with $ra\in I$ for all $r\in R$, $a\in I$. The ideal generated by
$a_1,\dots,a_n$ is $(a_1,\dots,a_n)=\{\sum r_ia_i\}$; an ideal $(a)$ is \emph{principal}. For ideals $I,J$: the sum
$I+J=\{a+b\}$, the product $IJ$ (the set of finite sums $\sum a_ib_i$, $a_i\in I$, $b_i\in J$) and the intersection $I\cap J$ are
ideals; $I$ and $J$ are \emph{comaximal} if $I+J=R$. For an $R$-module $M$ (modules are \Cref{ch:modules}; here $M$ may be any
abelian group on which $R$ acts) and $S\subseteq M$, the \emph{annihilator} is $\Ann_R(S)=\{r\in R:rs=0\text{ for all }s\in S\}$, an
ideal. The \emph{quotient ring} $R/I$ is the quotient group with $(a+I)(b+I)=ab+I$; the projection $R\to R/I$ is a
surjective homomorphism with kernel $I$.
\end{definition}

The multiplication on $R/I$ is well defined: if $a'=a+x$ and $b'=b+y$ with $x,y\in I$, then $a'b'-ab=ay+xb+xy\in I$.
This computation is the reason for the definition: an additive subgroup $I$ can serve as the set of elements declared to be $0$ in a quotient
ring exactly when it absorbs multiplication by all of $R$, since $ay$ and $xb$ must lie in $I$ for arbitrary $a,b$. Small examples: in $\Z$,
$(4,6)=2\Z$ and $4\Z\cdot6\Z=24\Z\subsetneq4\Z\cap6\Z=12\Z$; the ideal $(2,x)$ of $\Z[x]$ is the set of polynomials with even constant term;
$\Ann_\Z(\Z/6\Z)=6\Z$, and for the $\Z$-module $M=\Z/4\times\Z/6$ one has $\Ann_\Z(M)=12\Z$; $2\Z$ and $3\Z$ are comaximal since
$1=3-2$. The ideal $J=\Ann_\T(f^\perp)$ of Lemma lem:crit is an annihilator of exactly this kind, with $R=\T$ acting on the module $f^\perp$.

\begin{proposition}\label{prop:rings:ideals}
\begin{enumerate}[label=(\roman*)]
\item The kernel of a homomorphism is an ideal, and every ideal is the kernel of a homomorphism (the projection
$R\to R/I$). An ideal $I$ equals $R$ iff it contains a unit. A ring is a field iff it is nonzero and its only ideals
are $0$ and $R$.
\item $IJ\subseteq I\cap J$ for ideals $I,J$.
\item Every ideal of $\Z$ is $n\Z$ for a unique $n\ge0$. For $m,n\in\Z$, $m\Z+n\Z=\gcd(m,n)\Z$, and $m\Z\cap n\Z=\lcm(m,n)\Z$.
\end{enumerate}
\end{proposition}

\begin{proof}
(i) If $\phi(a)=\phi(b)=0$ then $\phi(a-b)=0$ and $\phi(ra)=\phi(r)\phi(a)=0$. The projection has kernel $I$. If $u\in I$ is a unit,
then $r=(ru^{-1})u\in I$ for every $r$. If $R$ is a field, a nonzero ideal contains a unit, so it is $R$; conversely, if $R\ne0$
has only the ideals $0$ and $R$ and $a\ne0$, then $(a)=R\ni1$, so $a$ is a unit.
(ii) Each $a_ib_i$ lies in $I$ (as $a_i\in I$) and in $J$ (as $b_i\in J$).
(iii) Let $I\ne0$ be an ideal of $\Z$ and $n$ its least positive element. For $a\in I$ write $a=bn+r$ with $0\le r<n$; then
$r=a-bn\in I$, so $r=0$ by minimality. Hence $I=n\Z$; $n$ is unique as the least positive element. Now $m\Z+n\Z=d\Z$ for
some $d\ge0$. Since $m,n\in d\Z$, $d$ is a common divisor; since $d=um+vn$ for some $u,v$, every common divisor divides $d$.
So $d=\gcd(m,n)$. Next, $m\Z\cap n\Z$ is an ideal (an intersection of ideals), so it is $e\Z$ for a unique $e\ge0$. Since $e\in m\Z\cap n\Z$,
$e$ is a common multiple of $m$ and $n$; and a common multiple $c$ lies in $m\Z\cap n\Z=e\Z$, so $e\mid c$. A nonnegative common multiple that
divides every common multiple is $\lcm(m,n)$ (with $\lcm(m,n)=0$ when $m$ or $n$ is $0$), so $e=\lcm(m,n)$.
\end{proof}

\begin{theorem}[isomorphism theorems]\label{thm:rings:isomorphism}
Let $\phi\colon R\to S$ be a homomorphism and $I$ an ideal of $R$.
\begin{enumerate}[label=(\roman*)]
\item $\phi$ induces an isomorphism $R/\ker\phi\simeq\phi(R)$.
\item (Correspondence.) $J\mapsto J/I$ is an inclusion-preserving bijection from the ideals of $R$ containing $I$ to the
ideals of $R/I$, with inverse $\bar J\mapsto$ its preimage, and $(R/I)/(J/I)\simeq R/J$. Under it, prime ideals correspond to
prime ideals and maximal ideals to maximal ideals (\Cref{thm:rings:prime-maximal}).
\item If $\phi$ is surjective with kernel $I$ and $J$ is any ideal of $R$, then $\phi(J)$ is an ideal of $S$ and $\phi$ induces an
isomorphism $R/(I+J)\simeq S/\phi(J)$.
\item Every quotient ring of $\Z$ is $\Z/r\Z$ for a unique $r\ge0$; as an additive group it is cyclic, generated by the
class of $1$.
\end{enumerate}
\end{theorem}

\begin{proof}
(i) The map $a+\ker\phi\mapsto\phi(a)$ is well defined and injective because $\phi(a)=\phi(b)$ iff $a-b\in\ker\phi$; it is
additive, multiplicative and sends $1$ to $1$, and it is onto $\phi(R)$.
(ii) For an ideal $\bar J$ of $R/I$ its preimage $J$ under the projection $\pi$ is an ideal containing $I$ with $\pi(J)=\bar J$, since $\pi$ is
onto. Conversely for $J\supseteq I$, $\pi(J)=J/I$ is an ideal of $R/I$ (an additive subgroup, stable under multiplication by
$\pi(r)$, as $\pi$ is onto), and its preimage is $J+I=J$. Both maps preserve inclusions. The composite $R\to R/I\to(R/I)/(J/I)$ is
onto with kernel $\pi^{-1}(J/I)=J$, so (i) gives $(R/I)/(J/I)\simeq R/J$. The statement on prime and maximal ideals follows
from this isomorphism and \Cref{thm:rings:prime-maximal}(i): $J$ is prime (maximal) iff $R/J$ is a domain (a field) iff
$(R/I)/(J/I)$ is.
(iii) $\phi(J)$ is an additive subgroup, and for $s=\phi(r)\in S$ (every $s$ is of this form) and $\phi(a)\in\phi(J)$,
$s\phi(a)=\phi(ra)\in\phi(J)$. The composite $R\to S\to S/\phi(J)$ is onto, and $a$ lies in its kernel iff $\phi(a)=\phi(b)$ for some $b\in J$,
iff $a-b\in I$ for some $b\in J$, iff $a\in I+J$. Now apply (i).
(iv) A quotient ring of $\Z$ is $\Z/I$ with $I=r\Z$, $r\ge0$ unique (\Cref{prop:rings:ideals}(iii)); every class is a multiple of the
class of $1$.
\end{proof}

\begin{proposition}\label{prop:rings:two-ideals}
Let $I,J$ be ideals of a ring $R$. The sequence of additive groups (indeed of $R$-modules)
\begin{equation}\label{eq:rings:two-ideals}
0\longrightarrow R/(I\cap J)\xrightarrow{\ x\mapsto(x,x)\ }R/I\times R/J\xrightarrow{\ (y,z)\mapsto y-z\ }R/(I+J)\longrightarrow0
\end{equation}
is exact. Consequently:
\begin{enumerate}[label=(\roman*)]
\item if $I\cap J=0$, then $R\to R/I\times R/J$ is injective, and its cokernel is isomorphic to $R/(I+J)$;
\item (Chinese remainder theorem for two ideals) if $I+J=R$, then $I\cap J=IJ$ and $R/IJ\to R/I\times R/J$ is a ring isomorphism;
\item (general Chinese remainder theorem) if $I_1,\dots,I_k$ are pairwise comaximal, then the natural map
$R/(I_1\cdots I_k)\to\prod_iR/I_i$ is a ring isomorphism.
\end{enumerate}
\end{proposition}

\begin{proof}
In \eqref{eq:rings:two-ideals} the first map is well defined (if $x\in I\cap J$ both components vanish) and the second also
(the class of $y-z$ modulo $I+J$ does not change if $y$ is moved by $I$ or $z$ by $J$). The first map is injective: $(x,x)=0$
means $x\in I$ and $x\in J$. The composite is zero. If $y-z\in I+J$, say $y-z=a+b$ with $a\in I$, $b\in J$, then $x:=y-a=z+b$
satisfies $x\equiv y\pmod I$ and $x\equiv z\pmod J$; so the kernel of the second map is the image of the first. The second map is
onto, since $(y,0)\mapsto y$.
(i) is the case $I\cap J=0$ read off from exactness.
(ii) If $I+J=R$ write $1=a+b$ with $a\in I$, $b\in J$. For $x\in I\cap J$, $x=xa+xb\in IJ$ (as $xa=ax\in JI$ and $xb\in IJ$); with
\Cref{prop:rings:ideals}(ii), $I\cap J=IJ$. Exactness with $R/(I+J)=0$ shows that $R/IJ\to R/I\times R/J$ is bijective; it is a ring
homomorphism because both coordinates are.
(iii) Induction on $k$, the case $k=2$ being (ii). For $k\ge3$, $I_1$ and $I_2\cdots I_k$ are comaximal: writing $1=a_j+b_j$ with
$a_j\in I_1$, $b_j\in I_j$ for $j\ge2$, the product $\prod_{j\ge2}(a_j+b_j)=1$ is $b_2\cdots b_k\in I_2\cdots I_k$ plus terms in $I_1$. By (ii) and
the induction hypothesis, $R/I_1\cdots I_k\simeq R/I_1\times R/I_2\cdots I_k\simeq R/I_1\times\prod_{j\ge2}R/I_j$.
\end{proof}

The maps in \eqref{eq:rings:two-ideals} commute with multiplication by elements of $R$, which is why the sequence is one of $R$-modules and not
only of groups. A worked case in $R=\Z$: for $I=4\Z$, $J=6\Z$ one has $I\cap J=12\Z$ and $I+J=2\Z$ (\Cref{prop:rings:ideals}(iii)), and the sequence is
\[
0\longrightarrow\Z/12\Z\longrightarrow\Z/4\Z\times\Z/6\Z\longrightarrow\Z/2\Z\longrightarrow0,
\]
with maps $x\mapsto(x\bmod4,x\bmod6)$ and $(y,z)\mapsto(y-z)\bmod2$.
The orders are consistent, $12\cdot2=24=4\cdot6$, and the cokernel $\Z/2\Z$ measures the failure of the Chinese remainder theorem: a pair
$(y,z)$ comes from an integer iff $y\equiv z\pmod2$. In (ii), the decomposition is computed by idempotents: for $I=3\Z$, $J=5\Z$, Bezout
gives $1=2\cdot5-3\cdot3$, and $a=-9\in I$, $b=10\in J$. The class of $b=10$ is $1$ modulo $3$ and $0$ modulo $5$, and $-9\equiv6\pmod{15}$ is $0$ modulo
$3$ and $1$ modulo $5$; both are idempotents of $\Z/15\Z$ ($10^2=100\equiv10$, $6^2=36\equiv6$). So the pair $(y\bmod3,z\bmod5)$ comes from
$10y+6z\bmod15$; for instance $(2,3)$ comes from $38\equiv8\pmod{15}$. This is the situation of Lemma lem:crit with the roles reversed: there
$I\cap J=0$ and the interesting group is the cokernel $R/(I+J)$, here it is $I+J=R$ and the cokernel vanishes.

\begin{proposition}\label{prop:rings:idempotents}
An \emph{idempotent} of $R$ is an element $e$ with $e^2=e$.
\begin{enumerate}[label=(\roman*)]
\item If $e$ is idempotent, so is $1-e$, and $e(1-e)=0$. The set $Re$ is a ring with identity $e$ (not a subring of $R$ unless
$e=1$), and $x\mapsto(xe,x(1-e))$ is a ring isomorphism $R\to Re\times R(1-e)$.
\item $R$ is isomorphic to a product of two nonzero rings iff it has an idempotent $e\ne0,1$.
\item The ideals of $R_1\times R_2$ are the products $I_1\times I_2$ of ideals; its prime ideals are the $P_1\times R_2$ and $R_1\times P_2$ with
$P_i$ prime, and the same holds for maximal ideals.
\item If $f=f_1\cdots f_k\in K[x]$ ($K$ a field) with pairwise coprime $f_i$, then $K[x]/(f)\simeq\prod_iK[x]/(f_i)$. For instance
$\Q[x]/(x^2-x)\simeq\Q\times\Q$, $x\mapsto(0,1)$.
\end{enumerate}
\end{proposition}

\begin{proof}
(i) $(1-e)^2=1-2e+e^2=1-e$ and $e(1-e)=e-e^2=0$. $Re$ is closed under addition and multiplication, and $e\cdot xe=xe^2=xe$, so
$e$ is its identity. The map $x\mapsto(xe,x(1-e))$ is additive and multiplicative ($xye=xe\cdot ye$ as $e^2=e$) and sends $1$ to
$(e,1-e)$, the identity of the product. It is injective, since $xe=x(1-e)=0$ gives $x=xe+x(1-e)=0$, and surjective, since
$(ae,b(1-e))$ is the image of $ae+b(1-e)$ (use $e(1-e)=0$).
(ii) If $R=R_1\times R_2$ with $R_i\ne0$, then $e=(1,0)$ is an idempotent different from $0$ and $1$. Conversely, given such $e$,
(i) gives $R\simeq Re\times R(1-e)$ with $Re\ni e\ne0$ and $R(1-e)\ni1-e\ne0$.
(iii) If $I$ is an ideal of $R_1\times R_2$, let $I_1$, $I_2$ be its projections; they are ideals. If $(a,b)\in I$ then
$(a,0)=(1,0)(a,b)\in I$ and $(0,b)\in I$; so $I=I_1\times I_2$. The quotient $(R_1\times R_2)/(I_1\times I_2)\simeq R_1/I_1\times R_2/I_2$ is a domain
(resp. a field) iff one factor is zero and the other a domain (resp. a field), because $(1,0)(0,1)=0$; by
\Cref{thm:rings:prime-maximal}(i) this gives the prime and maximal ideals.
(iv) This part uses two results proved later: $K[x]$ is a principal ideal domain (PID, \Cref{def:rings:factorisation}), by
\Cref{prop:rings:poly-division}(ii), and in a PID $(a)+(b)=(\gcd(a,b))$ (\Cref{thm:rings:euclid-pid-ufd}(ii)); neither proof uses this
proposition. So coprime $f_i$, $f_j$ satisfy $(f_i)+(f_j)=(\gcd(f_i,f_j))=K[x]$, and
$(f_1)\cdots(f_k)=(f)$; apply \Cref{prop:rings:two-ideals}(iii). For $x^2-x=x(x-1)$, the map is $g\mapsto(g(0),g(1))$, and the class of $x$
is the idempotent $(0,1)$.
\end{proof}

For example, in $R=\Z/6\Z$ the element $e=3$ is idempotent ($9\equiv3$), with $1-e=-2\equiv4$. Then $Re=\{0,3\}$, a ring with identity $3$
isomorphic to $\Z/2\Z$, and $R(1-e)=\{0,4,2\}$, a ring with identity $4$ isomorphic to $\Z/3\Z$; the map of (i) is $x\mapsto(3x,4x)$, and it
recovers $\Z/6\Z\simeq\Z/2\Z\times\Z/3\Z$. By (iii) the prime ideals of $\Z/6\Z$ are the two ideals $0\times\Z/3\Z$ and $\Z/2\Z\times0$, that is
$2\Z/6\Z=\{0,2,4\}$ and $3\Z/6\Z=\{0,3\}$, in accordance with the correspondence \Cref{thm:rings:isomorphism}(ii). In the paper the idempotents are
those of $\T\otimes\Q$ and $\T\otimes\Z_\ell$: $\T\otimes\Q=\Q\times\T'$ corresponds to the idempotent $e$ projecting onto the factor $\Q$ (the action on the line
spanned by $f$). The idempotent $e$ itself need not lie in $\T$: in the ring $R=\{(a,b)\in\Z^2:a\equiv b\pmod r\}$ of \Cref{ex:rings:toy-hecke}, the idempotent
$e=(1,0)$ of $R\otimes\Q=\Q\times\Q$ is not in $R$ for $r>1$, and the least positive integer $m$ with $me=(m,0)\in R$ is $m=r$, the order of $R/(I+J)$.

\begin{paperlink}\label{pl:rings:crit}
Lemma lem:crit (congruence criterion) is ring theory once one fact about Hecke algebras is granted. Let
$\phi_f\colon\T\to\Z$, $T_n\mapsto a_n(f)$, let $I=\ker\phi_f$ and $J=\Ann_\T(f^\perp)$. The paper writes $\T\otimes\Q=\Q\times\T'$, where
the first factor is the action on the line $\Q f$ and $\T'$ the (faithful) action on $f^\perp$; that this decomposition exists is
a statement about Hecke operators and the Petersson product (\Cref{ch:hecke}). Granting it, $\T$ embeds in $\T\otimes\Q$
($\T$ is torsion-free) and $t\mapsto(\phi_f(t),t')$. Then $t\in I$ iff its first coordinate vanishes, so $I=\T\cap(0\times\T')$; and $t\in J$ iff
$t$ kills $f^\perp$ iff $t'=0$, so $J=\T\cap(\Q\times0)$. Hence $I\cap J=0$. The map $\phi_f$ is onto $\Z$ because $T_1=1$ and
$\phi_f(T_1)=a_1(f)=1$; by \Cref{thm:rings:isomorphism}(iii),
\[
\T/(I+J)\simeq\Z/\phi_f(J),
\]
and this is cyclic (\Cref{thm:rings:isomorphism}(iv)): "$\T/(I+J)$ is cyclic of order $r$; equivalently, $r$ generates the
congruence ideal $\phi_f(J)$". By \Cref{prop:rings:two-ideals}(i), the inclusion $\T\subset\T/I\times\T/J$ has cokernel $\T/(I+J)$,
which is the sentence "Dualising the inclusion $\T\subset\T/I\times\T/J$, whose cokernel is $\T/(I+J)$" of the proof. The module
version of the same sequence, $L[\wp]\oplus L[I]\subseteq L$ in Lemma lem:ddt417, is \Cref{ch:modules}. In Step 1 of the proof of
Theorem thm:quant, and in Lemma lem:hecke, $\T\otimes\Z_\ell$ is a product of local rings cut out by idempotents
(\Cref{prop:rings:idempotents}; the general statement for finite $\Z_\ell$-algebras is \Cref{ch:modules}).
\end{paperlink}

\begin{proposition}\label{prop:rings:freshman}
Let $R$ be a ring of prime characteristic $p$.
\begin{enumerate}[label=(\roman*)]
\item $p\mid\binom pi$ for $0<i<p$; hence $(x+y)^p=x^p+y^p$ and $(x+y)^{p^k}=x^{p^k}+y^{p^k}$ for $x,y\in R$, $k\ge0$.
\item $F\colon x\mapsto x^p$ is a ring endomorphism of $R$ (the \emph{absolute Frobenius}). It is injective if $R$ is a domain and
bijective if $R$ is a finite domain, that is (\Cref{prop:rings:characteristic}(iii)) a finite field.
\item (Fermat) $a^p\equiv a\pmod p$ for every $a\in\Z$.
\end{enumerate}
\end{proposition}

\begin{proof}
(i) $\binom pi=\frac{p!}{i!(p-i)!}$ is an integer; $p$ divides the numerator and not the denominator when $0<i<p$ (a prime
dividing a product of integers smaller than $p$ would divide one of them: Euclid's lemma, \Cref{prop:rings:integers}(ii) below, whose
proof does not use this proposition). So the middle terms
of the binomial expansion vanish in $R$. The case $p^k$ follows by induction, applying the case $p$ to $x^{p^{k-1}}$ and $y^{p^{k-1}}$.
(ii) $F$ is additive by (i) and multiplicative by commutativity, and $F(1)=1$. In a domain $x^p=0$ forces $x=0$; an injective
map of a finite set to itself is bijective.
(iii) In $R=\Z/p\Z$, $F$ is additive and fixes $1$, so it fixes every $n\cdot1$: $n^p\equiv n\pmod p$.
\end{proof}

\section[Prime and maximal ideals, fractions, \texorpdfstring{$\Z_{(p)}$}{Z(p)}]{Prime and maximal ideals, fractions, and the ring \texorpdfstring{$\Z_{(p)}$}{Z(p)}}\label{sec:rings:primes}

Prime and maximal ideals are how the paper speaks about reduction: a maximal ideal $\mathfrak m$ of $\T\otimes\Z_\ell$ "containing
$\ell$ and $\ker\phi_f$" (Lemma lem:hecke, Theorem thm:quant) is a congruence modulo $\ell$ between $f$ and other eigenforms,
and the last step of Lemma lem:crit is that a maximal ideal of residue characteristic $\ell$ containing $I+J$ forces $\ell$
to divide $\#\T/(I+J)$. The ring $\Z_{(p)}$ of rationals with denominator prime to $p$ is the first local ring of the
book; it is where the paper's "nonzero rational number which is an $\ell$-adic integer" (Lemma lem:fpr) lives.

\begin{theorem}\label{thm:rings:prime-maximal}
An ideal $P\ne R$ is \emph{prime} if $ab\in P$ implies $a\in P$ or $b\in P$; an ideal $M\ne R$ is \emph{maximal} if the only ideals
containing $M$ are $M$ and $R$.
\begin{enumerate}[label=(\roman*)]
\item $R/P$ is a domain iff $P$ is prime; $R/M$ is a field iff $M$ is maximal. Every maximal ideal is prime.
\item The preimage of a prime ideal under a homomorphism is prime. The kernel of a homomorphism $\phi\colon R\to K$ to a field is
prime; it is maximal if $\phi$ is surjective, or more generally if $\phi(R)$ is finite.
\item If $R/P$ is finite, then $P$ is prime iff it is maximal.
\item The prime ideals of $\Z$ are $0$ and the $p\Z$; the maximal ideals are the $p\Z$.
\item Every nonzero ring has a maximal ideal; every non-unit of $R$ lies in a maximal ideal.
\end{enumerate}
\end{theorem}

\begin{proof}
(i) $R/P\ne0$ iff $P\ne R$, and $\bar a\bar b=0$ in $R/P$ means $ab\in P$; so "no zero divisors" is "prime". By the correspondence
\Cref{thm:rings:isomorphism}(ii) the ideals of $R/M$ are the $J/M$ with $M\subseteq J$, so $R/M$ has exactly two ideals iff $M$ is maximal,
and \Cref{prop:rings:ideals}(i) says this means $R/M$ is a field. A field is a domain.
(ii) If $Q\subset S$ is prime and $\phi\colon R\to S$, then $\phi^{-1}(Q)\ne R$ (as $\phi(1)=1\notin Q$), and $ab\in\phi^{-1}(Q)$ gives $\phi(a)\phi(b)\in Q$.
The zero ideal of a field is prime, so $\ker\phi$ is prime. By (i) of \Cref{thm:rings:isomorphism}, $R/\ker\phi\simeq\phi(R)$, a subring of a
field, hence a domain; if it is all of $K$ or finite, it is a field (\Cref{prop:rings:characteristic}(iii)), and $\ker\phi$ is maximal by (i).
(iii) $R/P$ is a finite ring; it is a domain iff it is a field (\Cref{prop:rings:characteristic}(iii)).
(iv) A nonzero ideal of $\Z$ is $n\Z$ with $n\ge1$ (\Cref{prop:rings:ideals}); $\Z/n$ is a domain iff $n$ is prime
(\Cref{prop:rings:zn}(i) below, proved from Bezout's identity and independent of this theorem), and is then finite, hence a field. $\Z/0=\Z$ is a domain and not a field.
(v) Let $a$ be a non-unit (for instance $a=0$ in a nonzero ring). The set $\Sigma$ of ideals $I\ne R$ containing $a$ is nonempty, as
$(a)\ne R$, and is partially ordered by inclusion. The union of a chain in $\Sigma$ is an ideal containing $a$, and it does not
contain $1$ because no member does; so it lies in $\Sigma$. By Zorn's lemma $\Sigma$ has a maximal element, which is a maximal ideal
containing $a$. (For Noetherian rings, \Cref{def:rings:factorisation} and \Cref{thm:rings:euclid-pid-ufd}(ii), Zorn's lemma can be
avoided: a nonempty set of ideals then has a maximal element.)
\end{proof}

Examples: in $\Z[x]$, $(x)$ is prime but not maximal, since $\Z[x]/(x)\simeq\Z$ (evaluation at $0$) is a domain but not a field; it lies in the maximal
ideal $(2,x)$, with $\Z[x]/(2,x)\simeq\F_2$. In $\Z$, $(4)$ is not prime ($2\cdot2\in(4)$). Part (ii) in its form "$\phi(R)$ finite" is the one used in the
paper: the map $\T\to\mathcal O/\lambda$ of \Cref{pl:rings:crit-last} need not be surjective, but its image is a subring of a finite field, hence
a finite domain, hence a field.

\begin{proposition}\label{prop:rings:residue-char}
Let $\mathfrak m$ be a maximal ideal of $R$ such that the field $R/\mathfrak m$ has characteristic $\ell>0$ (we say $\mathfrak m$ has
\emph{residue characteristic $\ell$}).
\begin{enumerate}[label=(\roman*)]
\item $\ell\cdot1\in\mathfrak m$.
\item If $\mathfrak a\subseteq\mathfrak m$ is an ideal with $R/\mathfrak a$ finite, then $\ell\mid\#(R/\mathfrak a)$.
\item For $r\ge1$, there is a ring homomorphism $\Z/r\Z\to k$ into a field $k$ of characteristic $\ell$ iff $\ell\mid r$.
\end{enumerate}
\end{proposition}

\begin{proof}
(i) The image of $\ell\cdot1$ in $R/\mathfrak m$ is $\ell\cdot1=0$.
(ii) The projection $R/\mathfrak a\to R/\mathfrak m$ is onto (\Cref{thm:rings:isomorphism}(ii)), so $\#(R/\mathfrak a)=\#(\mathfrak m/\mathfrak a)\cdot\#(R/\mathfrak m)$ by
Lagrange's theorem for the additive group. The field $R/\mathfrak m$ contains $\F_\ell$ (the image of $\Z$) and is a vector space over
it of dimension $\ge1$, finite as $R/\mathfrak a$ is; so $\#(R/\mathfrak m)$ is a positive power of $\ell$.
(iii) A homomorphism $\Z/r\to k$ composed with $\Z\to\Z/r$ is the unique map $\Z\to k$ (\Cref{prop:rings:characteristic}), whose kernel is
$\ell\Z$; it factors through $\Z/r$ iff $r\Z\subseteq\ell\Z$, i.e. iff $\ell\mid r$.
\end{proof}

\begin{paperlink}\label{pl:rings:crit-last}
The last step of the proof of Lemma lem:crit reads: "$\phi_g\equiv\phi_f\pmod\lambda$ on all of $\T$ \dots\ Thus $I+J$ lies in the
kernel of $\T\to\mathcal O/\lambda$, a maximal ideal of residue characteristic $\ell$, and $\ell$ divides $\#\T/(I+J)=r$." Here the
map $\T\to\mathcal O/\lambda$ is $t\mapsto\phi_f(t)\bmod\lambda$; its image is a subring of the finite field $\mathcal O/\lambda$ (whose
characteristic is $\ell$, as $\lambda\mid\ell$), so its kernel $\mathfrak m$ is maximal by \Cref{thm:rings:prime-maximal}(ii), of residue
characteristic $\ell$. It contains $I=\ker\phi_f$, and $J$ because $\phi_g$ kills $J$ and $\phi_g\equiv\phi_f\pmod\lambda$. Since
$\T/(I+J)\simeq\Z/r$ is finite (if $r\ne0$), \Cref{prop:rings:residue-char}(ii) gives $\ell\mid r$; equivalently, by (iii), $\Z/r$ maps to a field of
characteristic $\ell$. The same reasoning is behind "after enlarging $\mathcal O$ and $\lambda$" in Lemma lem:old.
\end{paperlink}

\begin{theorem}\label{thm:rings:fractions}
Let $R$ be a domain. On pairs $(a,b)\in R\times(R\setminus\{0\})$ put $(a,b)\sim(c,d)$ iff $ad=bc$, and write $a/b$ for the class. Then
$\Frac(R)=\{a/b\}$ with $a/b+c/d=(ad+bc)/bd$ and $(a/b)(c/d)=ac/bd$ is a field, $a\mapsto a/1$ is an injective homomorphism $R\to\Frac(R)$,
and every injective homomorphism $\iota\colon R\to K$ into a field extends uniquely to a homomorphism $\Frac(R)\to K$, namely
$a/b\mapsto\iota(a)\iota(b)^{-1}$. Examples: $\Frac(\Z)=\Q$, $\Frac(F[x])=F(x)$, $\Frac(\Z[i])=\Q(i)$, $\Frac(\Z[\zeta_3])=\Q(\zeta_3)=\Q(\sqrt{-3})$.
\end{theorem}

\begin{proof}
$\sim$ is reflexive and symmetric; if $ad=bc$ and $cf=de$ then $adf=bcf=bde$, so $d(af-be)=0$ and $af=be$ as $d\ne0$: it is
transitive. The operations are well defined: if $ab'=a'b$, then $(ad+bc)b'd=a'bd^2+bb'cd=(a'd+b'c)bd$ and $ac\cdot b'd=a'c\cdot bd$, so replacing $(a,b)$ by
$(a',b')$ does not change the classes of $(ad+bc,bd)$ and $(ac,bd)$ (by symmetry the same holds in the second argument). The ring axioms are
checked by computing both sides, using that $(x,y)\sim(tx,ty)$ for $t\ne0$; for instance for distributivity,
$\frac ab\bigl(\frac cd+\frac ef\bigr)=\frac{a(cf+de)}{bdf}$ and $\frac{ac}{bd}+\frac{ae}{bf}=\frac{acbf+bdae}{b^2df}=\frac{b\cdot a(cf+de)}{b\cdot bdf}$, which are equal; associativity of
the addition holds because both bracketings give $(adf+bcf+bde)/bdf$. Here $bd\ne0$ because $R$ is a domain, which is where the
hypothesis is used: in $\Z/6\Z$ the would-be fraction $1/2\cdot1/3$ has denominator $0$. A nonzero $a/b$ has $a\ne0$ and inverse $b/a$. The map
$a\mapsto a/1$ is a homomorphism, injective since $a/1=0/1$ means $a=0$. If $\iota$ is injective, $\iota(b)\ne0$ for $b\ne0$, so
$a/b\mapsto\iota(a)\iota(b)^{-1}$ makes sense; it is well defined ($ad=bc$ gives $\iota(a)\iota(d)=\iota(b)\iota(c)$) and a homomorphism.
Uniqueness: any extension $\tilde\iota$ satisfies $\tilde\iota(a/b)\iota(b)=\tilde\iota(a/b\cdot b/1)=\iota(a)$. For the examples: $\Q(i)$ and
$\Q(\zeta_3)$ are fields containing $\Z[i]$ and $\Z[\zeta_3]$ in which every element is a quotient (for $z=x+yi\in\Q(i)$ clear the
denominators of $x,y$); $\zeta_3=(-1+\sqrt{-3})/2$ gives $\Q(\zeta_3)=\Q(\sqrt{-3})$.
\end{proof}

\begin{proposition}\label{prop:rings:zp-local}
Let $p$ be a prime and $\Z_{(p)}=\{a/b\in\Q:a,b\in\Z,\ p\nmid b\}$.
\begin{enumerate}[label=(\roman*)]
\item $\Z_{(p)}$ is a subring of $\Q$, and $\Z_{(p)}^\times=\{a/b:p\nmid ab\}$; an integer prime to $p$ is a unit of $\Z_{(p)}$.
\item Every nonzero element is $p^ku$ with $k\ge0$ and $u\in\Z_{(p)}^\times$, uniquely; the nonzero ideals are the $p^k\Z_{(p)}$, $k\ge0$, and
$p\Z_{(p)}$ is the only maximal ideal (so $\Z_{(p)}$ is a \emph{local ring}; it is the first example of a discrete valuation ring,
\Cref{ch:modules}).
\item The reduction map $\Z_{(p)}\to\F_p$, $a/b\mapsto\bar a\bar b^{-1}$, is a well-defined surjective ring homomorphism with kernel
$p\Z_{(p)}$; it maps $\Z_{(p)}^\times$ onto $\F_p^\times$.
\item $\Z_{(p)}=\{x\in\Q:v_p(x)\ge0\}$ and $\Z_{(p)}^\times=\{x\in\Q:v_p(x)=0\}$.
\item $\Z=\bigcap_p\Z_{(p)}$.
\end{enumerate}
\end{proposition}

\begin{proof}
Every rational is $a/b$ in lowest terms with $b>0$, and $x\in\Z_{(p)}$ iff $p$ does not divide this $b$ (if $x=a'/b'$ with $p\nmid b'$, the
reduced denominator divides $b'$). (i) Sums and products of fractions with denominators prime to $p$ have denominators
$bd$ prime to $p$ (Euclid's lemma, \Cref{prop:rings:integers}). In lowest terms, $a/b$ has an inverse $b/a$ in $\Z_{(p)}$ iff $p\nmid a$.
(ii) Write $a=p^ka'$ with $p\nmid a'$ (unique factorisation in $\Z$); then $a/b=p^k\cdot a'/b$ with $a'/b$ a unit, and $k=v_p(a/b)$ is
determined (\Cref{def:rings:valuation}). If $I\ne0$ is an ideal, let $k$ be the least exponent of its elements; then $p^k=(p^ku)u^{-1}\in I$
and every element $p^{k'}u'$ of $I$, $k'\ge k$, lies in $p^k\Z_{(p)}$; so $I=p^k\Z_{(p)}$. These ideals form a chain, and the largest proper one
is $p\Z_{(p)}$, which is the set of non-units.
(iii) If $p\nmid b$ then $\bar b\in\F_p^\times$; if $a/b=a'/b'$ then $ab'=a'b$, so $\bar a\bar b^{-1}=\bar a'\bar b'^{-1}$. Additivity and
multiplicativity follow from the formulas for fractions. It is onto because it restricts to $\Z\to\F_p$. Its kernel is
$\{a/b:p\mid a\}=p\Z_{(p)}$, and units map to units.
(iv) is (i) and (ii) rewritten with $v_p$. (v) A rational in lowest terms $a/b$ with $b>1$ is not in $\Z_{(p)}$ for any prime $p\mid b$.
\end{proof}

For example, in $\Z_{(5)}$ the element $\frac12$ is a unit and reduces to $\bar2^{-1}=3\in\F_5$ (as $2\cdot3=6\equiv1$); the element $\frac{10}3=5\cdot\frac23$ has $v_5=1$ and
reduces to $0$; and $\frac15\notin\Z_{(5)}$, so it has no reduction. The word \emph{local} refers to the fact that $\Z_{(p)}$ only sees the prime $p$:
every other prime has become a unit, and the non-units form the single ideal $p\Z_{(p)}$. This is the setting of every statement of the form
"$x$ is an $\ell$-adic integer": it means $x\in\Z_{(\ell)}$, i.e. $\ell$ does not divide the reduced denominator of $x$.

\begin{paperlink}\label{pl:rings:fpr}
In Lemma lem:fpr, "$L_p(A_f,1)^{-1}=L_p(A_f(1),0)^{-1}$ is a nonzero rational number which is an $\ell$-adic integer" means: it lies in
$\Z_{(\ell)}$ (\Cref{ch:modules,ch:local} replace $\Z_{(\ell)}$ by $\Z_\ell$, and $\Z_{(\ell)}=\Q\cap\Z_\ell$). Indeed $L_p(A_f,1)^{-1}=P_p(p^{-2})$ with $P_p\in\Z[X]$,
a rational whose denominator is a power of $p\ne\ell$, and $p$ is a unit of $\Z_{(\ell)}$ by \Cref{prop:rings:zp-local}(i); see
\Cref{lem:rings:ell-integral}. The same remark gives sense to Definition def:terms ($q$ is an $\ell$-adic unit, so
$v_\ell(1-a_q^2q^{-2})$ and $v_\ell(P_q(q^{-2}))$ are valuations of elements of $\Z_{(\ell)}$). In Lemma lem:tamfactor, "the class of the
reduction of $u$ in $\F_q^\times$" is the image of the unit $u=q_Eq^{-n}\in\Z_q^\times$ under the reduction map, of which
\Cref{prop:rings:zp-local}(iii) is the rational version; \Cref{ch:local} gives the reduction $\Z_q\to\F_q$.
\end{paperlink}

\section[Euclidean domains, PIDs, UFDs]{Euclidean domains, PIDs and UFDs; the rings \texorpdfstring{$\Z[i]$}{Z[i]} and \texorpdfstring{$\Z[\zeta_3]$}{Z[zeta3]}}\label{sec:rings:ufd}

This section proves unique factorisation where the book needs it: in $\Z$ (the $\ell$-part of an integer, used in every
statement of the paper about $v_\ell(\degphi)$), in $F[x]$ (minimal polynomials, finite fields), and in the rings $\Z[i]$ and
$\Z[\zeta_3]$ of Gaussian and Eisenstein integers, whose characters $\chi_{-4}$ and $\chi_{-3}$ appear in the table of Section 1 and
whose $\ell$-adic versions $\Z_\ell\otimes\Z[\zeta_3]$ (a product of two copies of $\Z_\ell$, so not a domain, when
$\ell\equiv1\pmod3$) and $\Z_3[\zeta_3]$ appear in the proofs of Lemmas lem:euler and lem:IV. The primes of $\Z[i]$ and
$\Z[\zeta_3]$, the element $1-\zeta_3$ as a uniformiser and the ring $\Z_3[\zeta_3]$ are treated in \Cref{ch:ant}.

\begin{definition}\label{def:rings:factorisation}
Let $R$ be a domain. For $a,b\in R$, $a\mid b$ if $b=ac$ for some $c$; $a$ and $b$ are \emph{associates} if $b=ua$ with $u\in R^\times$. A nonzero
non-unit $a$ is \emph{irreducible} if $a=bc$ implies that $b$ or $c$ is a unit, and \emph{prime} if $a\mid bc$ implies $a\mid b$ or $a\mid c$. A
\emph{greatest common divisor} of $a$ and $b$ is a common divisor divisible by every common divisor (unique up to units if it
exists). $R$ is
\begin{itemize}
\item a \emph{Euclidean domain} if there is $N\colon R\setminus\{0\}\to\Z_{\ge0}$ such that for all $a$ and all $b\ne0$ there are $\kappa,\rho\in R$ with
$a=\kappa b+\rho$ and $\rho=0$ or $N(\rho)<N(b)$;
\item a \emph{principal ideal domain} (PID) if every ideal is principal;
\item a \emph{unique factorisation domain} (UFD) if every nonzero non-unit is a product of irreducibles, uniquely up to order and
associates.
\end{itemize}
A ring is \emph{Noetherian} if every ascending chain of ideals $I_1\subseteq I_2\subseteq\cdots$ is eventually constant.
\end{definition}

\begin{lemma}\label{lem:rings:prime-irreducible}
Let $R$ be a domain.
\begin{enumerate}[label=(\roman*)]
\item A prime element is irreducible.
\item For $a\ne0$, the ideal $(a)$ is prime iff $a$ is a prime element.
\item If $p_1\cdots p_r=uq_1\cdots q_s$ with $p_i$, $q_j$ prime elements and $u$ a unit, then $r=s$ and, after renumbering, $q_i$ is an
associate of $p_i$ for every $i$.
\end{enumerate}
\end{lemma}

\begin{proof}
(i) If $p=bc$ is prime then $p\mid bc$, say $p\mid b$, $b=pt$; then $p=ptc$, so $tc=1$ and $c$ is a unit.
(ii) $ab\in(a')$ means $a'\mid ab$; $(a')\ne R$ means $a'$ is not a unit. So the definitions coincide.
(iii) Induction on $r$. If $r=0$ the right side is a unit, so $s=0$ (a prime is not a unit). Otherwise $p_1$ divides
$uq_1\cdots q_s$, hence some $q_j$ (it does not divide $u$); $q_j=p_1t$ with $t$ a unit by (i). Renumber so $j=1$, cancel $p_1$ (in a
domain $p_1x=p_1y$ gives $x=y$), and apply induction to $p_2\cdots p_r=(ut)q_2\cdots q_s$.
\end{proof}

\begin{theorem}\label{thm:rings:euclid-pid-ufd}
\begin{enumerate}[label=(\roman*)]
\item A Euclidean domain is a PID.
\item A PID is Noetherian and a UFD. In a PID: $\gcd(a,b)$ exists, with $(a,b)=(\gcd(a,b))$ (Bezout's identity); irreducible elements are
prime; nonzero prime ideals are maximal.
\item In a UFD irreducible elements are prime, and greatest common divisors exist.
\item $\Z$ (with $N(a)=\abs a$) and $F[x]$ for a field $F$ (with $N=\deg$, \Cref{prop:rings:poly-division}) are Euclidean. $\Z[x]$ and $F[x,y]$ are UFDs
(\Cref{thm:rings:gauss-lemma}) but not PIDs: $(2,x)\subset\Z[x]$ and $(x,y)\subset F[x,y]$ are not principal.
\end{enumerate}
\end{theorem}

\begin{proof}
(i) Let $I\ne0$ be an ideal and $b\in I\setminus\{0\}$ with $N(b)$ minimal. For $a\in I$ write $a=\kappa b+\rho$; then $\rho=a-\kappa b\in I$, so $\rho=0$
by minimality. Hence $I=(b)$.
(ii) Let $I_1\subseteq I_2\subseteq\cdots$ be ideals of a PID. The union is an ideal, $(a)$ say; $a$ lies in some $I_n$, and then $I_m=(a)$ for
$m\ge n$. For $a,b$ write $(a,b)=(d)$; $d$ divides $a$ and $b$, and $d=ua+vb$ is divisible by every common divisor, so $d=\gcd(a,b)$.
Let $p$ be irreducible with $p\mid ab$ and $p\nmid a$. Then $\gcd(p,a)$ divides $p$, so it is a unit or an associate of $p$; the latter
would give $p\mid a$. So $1=up+va$ and $b=upb+vab$ is divisible by $p$: $p$ is prime. If $(p)\ne0$ is prime and $(p)\subseteq(m)\ne R$,
then $p=mt$, and as $p$ is irreducible (\Cref{lem:rings:prime-irreducible}) and $m$ is not a unit, $t$ is a unit and $(m)=(p)$; so $(p)$
is maximal. Existence of factorisations: suppose a nonzero non-unit $a_0$ is not a product of irreducibles. Then $a_0$ is not
irreducible, so $a_0=bc$ with $b,c$ non-units, and one of them, $a_1$ say, is not a product of irreducibles; $(a_0)\subsetneq(a_1)$
because the cofactor is not a unit. Repeating gives a strictly increasing chain $(a_0)\subsetneq(a_1)\subsetneq\cdots$, which is
impossible. Uniqueness follows from \Cref{lem:rings:prime-irreducible}(iii), since irreducibles are prime.
(iii) Let $p$ be irreducible with $p\mid ab$, $ab=pc$. Factoring $a$, $b$, $c$ into irreducibles, the two factorisations
$(\text{factors of }a)(\text{factors of }b)$ and $p(\text{factors of }c)$ of $ab$ agree up to order and units, so $p$ is an associate of an
irreducible factor of $a$ or of $b$. For gcds, fix one irreducible $p$ in each class of associates, write
$a=u\prod p^{e_p(a)}$, $b=u'\prod p^{e_p(b)}$, and take $\prod p^{\min(e_p(a),e_p(b))}$: by uniqueness, $d\mid a$ iff $e_p(d)\le e_p(a)$ for all $p$.
(iv) Division with remainder in $\Z$ is $a=\kappa b+\rho$ with $0\le\rho<\abs b$; in $F[x]$ it is \Cref{prop:rings:poly-division}(i), as every
nonzero polynomial over a field is a unit times a monic one. If $(2,x)=(d)$ in $\Z[x]$, then $d\mid2$ forces $\deg d=0$, and $d\mid x$
forces $d=\pm1$; but then $1=2g+xh$, and evaluating at $x=0$ gives $1=2g(0)$ in $\Z$, impossible. If $(x,y)=(d)$ in $F[x,y]$, then $d$
divides $x$ and $y$, so it is a nonzero constant, and $1=xg+yh$ evaluated at $(0,0)$ gives $1=0$.
\end{proof}

\begin{proposition}\label{prop:rings:integers}
\begin{enumerate}[label=(\roman*)]
\item (Bezout) For $m,n\in\Z$ there are $u,v\in\Z$ with $um+vn=\gcd(m,n)$.
\item (Euclid's lemma) If a prime $p$ divides $ab$, it divides $a$ or $b$.
\item (Fundamental theorem of arithmetic) Every $n\ne0$ is $n=\pm\prod_pp^{v_p(n)}$ with unique exponents $v_p(n)\ge0$, almost all zero.
\item $d\mid n$ iff $v_p(d)\le v_p(n)$ for all $p$; hence $\gcd(m,n)=\prod_pp^{\min(v_p(m),v_p(n))}$ and $\lcm(m,n)=\prod_pp^{\max(v_p(m),v_p(n))}$ for
$m,n\ne0$. The \emph{$p$-part} of $n\ne0$ is $p^{v_p(n)}$.
\end{enumerate}
\end{proposition}

\begin{proof}
$\Z$ is Euclidean, hence a PID (\Cref{thm:rings:euclid-pid-ufd}). (i) is Bezout's identity there (or \Cref{prop:rings:ideals}(iii)); (ii) says
that irreducible elements of $\Z$, the $\pm p$, are prime; (iii) is unique factorisation, the units being $\pm1$ and the positive
primes a system of representatives of the irreducibles. (iv) follows from (iii) as in the proof of \Cref{thm:rings:euclid-pid-ufd}(iii).
\end{proof}

\begin{theorem}\label{thm:rings:gauss-eisenstein-integers}
Let $R$ be $\Z[i]=\{a+bi:a,b\in\Z\}$ or $\Z[\zeta_3]=\{a+b\zeta_3:a,b\in\Z\}$, subrings of $\C$ (note $\zeta_3^2=-1-\zeta_3$), and $N(z)=z\bar z=\abs z^2$.
Then $N(a+bi)=a^2+b^2$ and $N(a+b\zeta_3)=a^2-ab+b^2$ are nonnegative integers, zero only at $0$, and $N(zw)=N(z)N(w)$.
\begin{enumerate}[label=(\roman*)]
\item For $\alpha,\beta\in R$, $\beta\ne0$, there is $\kappa\in R$ with $N(\alpha-\kappa\beta)\le\frac12N(\beta)$ if $R=\Z[i]$, and $\le\frac13N(\beta)$ if $R=\Z[\zeta_3]$.
Hence both rings are Euclidean for $N$, so PIDs and UFDs.
\item $z\in R$ is a unit iff $N(z)=1$; $\Z[i]^\times=\{\pm1,\pm i\}=\mu_4$ and $\Z[\zeta_3]^\times=\{\pm1,\pm\zeta_3,\pm\zeta_3^2\}=\mu_6$.
\item An element of prime norm is irreducible, hence prime.
\item $2=-i(1+i)^2$, $5=(2+i)(2-i)$, $3=-\zeta_3^2(1-\zeta_3)^2$, $7=(3+\zeta_3)(3+\zeta_3^2)$, $13=(4+\zeta_3)(4+\zeta_3^2)$, the factors being primes.
\end{enumerate}
\end{theorem}

\begin{proof}
Both sets are closed under multiplication ($i^2=-1$; $\zeta_3^2=-1-\zeta_3$ because $\zeta_3\ne1$ is a root of $x^3-1=(x-1)(x^2+x+1)$) and
under complex conjugation ($\bar\zeta_3=\zeta_3^2$). $N(a+b\zeta_3)=a^2+ab(\zeta_3+\bar\zeta_3)+b^2=a^2-ab+b^2$. Multiplicativity is that of $\abs\cdot^2$.
(i) Put $\gamma=\alpha/\beta\in\C$. It suffices to find $\kappa\in R$ with $\abs{\gamma-\kappa}^2\le\frac12$, resp. $\frac13$, and then
$N(\alpha-\kappa\beta)=\abs{\gamma-\kappa}^2N(\beta)$. In the first case $R$ is the square lattice, and $\C$ is tiled by the closed unit squares with
vertices in $R$. In the second, the parallelograms with vertices $\kappa,\kappa+1,\kappa+\zeta_3,\kappa+1+\zeta_3$ ($\kappa\in R$) tile $\C$, and the
diagonal from $\kappa$ to $\kappa+1+\zeta_3=\kappa-\zeta_3^2$ cuts each into two equilateral triangles of side $1$ (all sides have length $\abs1=\abs{\zeta_3}=\abs{1+\zeta_3}=1$).
So $\gamma$ lies in a square of side $1$, resp. an equilateral triangle of side $1$, with vertices in $R$. Join the centre $c$
(of the square, resp. the centroid of the triangle) to the midpoints of the sides: this cuts the polygon into convex
quadrilaterals, one at each vertex $w$, with corners $w$, two half-sides and $c$. The distance to $w$ is a convex function,
so on such a quadrilateral it is maximal at a corner; the corners are at distance $0$, $\frac12$, $\frac12$ and $\abs{c-w}$, which is
$\frac{\sqrt2}2$ for the square and the circumradius $\frac1{\sqrt3}$ for the triangle. Hence $\abs{\gamma-w}^2\le\frac12$, resp. $\frac13$, for the vertex $w$ of
the piece containing $\gamma$; take $\kappa=w$. The bound is attained at the centre, resp. the centroid. By
\Cref{thm:rings:euclid-pid-ufd}, both rings are PIDs and UFDs.
(ii) If $zw=1$ then $N(z)N(w)=1$ with both positive integers, so $N(z)=1$; conversely $N(z)=1$ gives $z\cdot\bar z=1$ with $\bar z\in R$. Now
$a^2+b^2=1$ has the solutions $(\pm1,0)$, $(0,\pm1)$; and $a^2-ab+b^2=1$, i.e. $(2a-b)^2+3b^2=4$, forces $b\in\{0,\pm1\}$ and gives $(a,b)=(\pm1,0)$, $(0,1)$,
$(1,1)$, $(0,-1)$, $(-1,-1)$, that is $\pm1$, $\zeta_3$, $1+\zeta_3=-\zeta_3^2$, $-\zeta_3$, $-1-\zeta_3=\zeta_3^2$. These are the six sixth roots of unity, $-\zeta_3$ having
order $6$.
(iii) If $z=xy$ with $N(z)$ prime, one of $N(x)$, $N(y)$ is $1$. In a PID irreducibles are prime.
(iv) Direct multiplication: $(1+i)^2=2i$; $(1-\zeta_3)^2=1-2\zeta_3+\zeta_3^2=-3\zeta_3$, so $-\zeta_3^2(1-\zeta_3)^2=3\zeta_3^3=3$; $(3+\zeta_3)(3+\zeta_3^2)=9+3(\zeta_3+\zeta_3^2)+1=7$;
$(4+\zeta_3)(4+\zeta_3^2)=16-4+1=13$. The norms of the factors are $2,5,5,3,7,7,13,13$, prime, so the factors are prime by (iii).
\end{proof}

\begin{example}[{the Euclidean algorithm in $\Z[i]$}]\label{ex:rings:euclid-gaussian}
To compute $\gcd(11+3i,1+8i)$ (norms $130$ and $65$): $\frac{11+3i}{1+8i}=\frac{(11+3i)(1-8i)}{65}=\frac{35-85i}{65}$, nearest Gaussian integer
$1-i$, and $11+3i=(1-i)(1+8i)+(2-4i)$ with $N(2-4i)=20<65$. Next $\frac{1+8i}{2-4i}=\frac{(1+8i)(2+4i)}{20}=\frac{-30+20i}{20}$, take $-2+i$:
$1+8i=(-2+i)(2-4i)+(1-2i)$, $N(1-2i)=5<20$. Finally $2-4i=2(1-2i)$. So $\gcd(11+3i,1+8i)=1-2i$, of norm $5$, up to units. Although $13$ divides both
norms, the factorisations $11+3i=(1-2i)(1+i)(3+2i)$ and $1+8i=-(1-2i)(3-2i)$ show that different primes of norm $13$ occur:
$13=(3+2i)(3-2i)$, and $3-2i$ is not an associate of $3+2i$.
\end{example}

\begin{example}[unique factorisation fails in general]\label{ex:rings:non-ufd}
In $\Z[\sqrt{-5}]=\{a+b\sqrt{-5}\}$, with $N(a+b\sqrt{-5})=a^2+5b^2$, one has $6=2\cdot3=(1+\sqrt{-5})(1-\sqrt{-5})$. The four factors have norms $4$, $9$, $6$, $6$;
no element has norm $2$ or $3$, so each factor is irreducible; the units are $\pm1$ (norm $1$), so no two of the factors are
associates. Thus $\Z[\sqrt{-5}]$ is not a UFD, hence not a PID; explicitly, $(2,1+\sqrt{-5})$ is not principal: it lies in the kernel of
the homomorphism $\Z[\sqrt{-5}]\to\F_2$, $a+b\sqrt{-5}\mapsto a+b\bmod2$ (well defined since $\Z[\sqrt{-5}]\simeq\Z[x]/(x^2+5)$ and $1^2+5=0$ in $\F_2$), so it is proper,
while a generator $d$ would have $N(d)\mid\gcd(4,6)=2$, so $N(d)=1$ and the ideal would be everything. Similarly
$4=2\cdot2=(1+\sqrt{-3})(1-\sqrt{-3})$ in $\Z[\sqrt{-3}]\subsetneq\Z[\zeta_3]$, with three non-associate irreducibles of norm $4$; in $\Z[\zeta_3]$ the
anomaly disappears because $1+\sqrt{-3}=2+2\zeta_3=-2\zeta_3^2$ is an associate of $2$. The ideal-theoretic repair is \Cref{ch:ant}.
\end{example}

\section{Polynomial rings and irreducibility}\label{sec:rings:polynomials}

Polynomials enter the paper as characteristic polynomials of Frobenius ($X^2-a_pX+p$), as the Euler factors $P_q(X)\in\Z[X]$ of
Definition def:terms, as the $3$-division polynomial of Proposition prop:ns, and, behind the scenes, through finite
fields and cyclotomic polynomials. This section proves division with remainder, the criterion for repeated roots,
Gauss's lemma, the criteria of Eisenstein and of reduction modulo $p$, and Kronecker's construction of a field in
which a given polynomial has a root. The input that a polynomial of degree $d$ over a field has at most $d$ roots
(\Cref{prop:rings:poly-division}(iv)) is what \Cref{ch:groups} uses to prove that finite subgroups of $K^\times$ are cyclic.

\begin{proposition}\label{prop:rings:poly-division}
Let $R$ be a ring and $R[x]$ the polynomial ring; $\deg0=-\infty$.
\begin{enumerate}[label=(\roman*)]
\item (Division by a monic polynomial.) If $g\in R[x]$ is monic of degree $m$, every $f\in R[x]$ is uniquely $f=hg+r$ with $\deg r<m$.
\item If $R$ is a domain, $\deg fg=\deg f+\deg g$, $R[x]$ is a domain and $R[x]^\times=R^\times$. If $F$ is a field, $F[x]$ is Euclidean for the
degree, hence a PID and a UFD, and its maximal ideals are the $(f)$ with $f$ irreducible.
\item (Factor theorem.) For $a\in R$, $f(a)=0$ iff $(x-a)\mid f$.
\item A nonzero polynomial of degree $d$ over a domain $R$ has at most $d$ roots in $R$.
\item $x^p-x\in\F_p[x]$ is nonzero but vanishes at every point of $\F_p$; over an infinite domain a polynomial is determined by
its values.
\end{enumerate}
\end{proposition}

\begin{proof}
(i) Existence by induction on $\deg f$: if $\deg f=n\ge m$ with leading coefficient $a$, then $f-ax^{n-m}g$ has degree $<n$. Uniqueness: if
$hg+r=h'g+r'$ with $h\ne h'$, then $(h-h')g=r'-r$; as $g$ is monic, the left side has degree $\deg(h-h')+m\ge m>\deg(r'-r)$.
(ii) The leading coefficient of $fg$ is the product of those of $f$ and $g$, nonzero in a domain. If $fg=1$ then $\deg f=\deg g=0$.
Over a field every nonzero $g$ is a unit times a monic polynomial, so (i) gives the Euclidean property; then
\Cref{thm:rings:euclid-pid-ufd} applies. A nonzero prime ideal of a PID is maximal, $(f)$ is prime iff $f$ is irreducible
(\Cref{lem:rings:prime-irreducible}, \Cref{thm:rings:euclid-pid-ufd}(ii)), and $0$ is not maximal since $F[x]$ is not a field.
(iii) Divide by the monic $x-a$: $f=h\cdot(x-a)+r$ with $r$ constant; evaluating at $a$ gives $r=f(a)$.
(iv) Induction on $d$. If $f(a)=0$ then $f=(x-a)h$ with $\deg h=d-1$, and a root $b\ne a$ of $f$ satisfies $(b-a)h(b)=0$, so $h(b)=0$.
(v) $a^p=a$ for $a\in\F_p$ (\Cref{prop:rings:freshman}(iii)). Over an infinite domain, two polynomials with the same values differ
by a polynomial with infinitely many roots, which is $0$ by (iv).
\end{proof}

\begin{proposition}\label{prop:rings:derivative}
The formal derivative of $f=\sum a_ix^i$ is $f'=\sum ia_ix^{i-1}$; it is additive and $(fg)'=f'g+fg'$. Let $F$ be a field, $f\in F[x]$ nonzero and
$K\supseteq F$ a field.
\begin{enumerate}[label=(\roman*)]
\item $a\in K$ is a multiple root of $f$ (that is, $(x-a)^2\mid f$ in $K[x]$) iff $f(a)=f'(a)=0$.
\item The gcd of $f$ and $f'$ is the same in $F[x]$ and in $K[x]$; $f$ has no repeated root in any extension field iff $\gcd(f,f')=1$
(we then call $f$ \emph{separable}).
\item In characteristic $p$, $x^{p^r}-x$ is separable. For $n\ge1$, $x^n-1$ is separable over $F$ iff $\operatorname{char}F\nmid n$.
\end{enumerate}
\end{proposition}

\begin{proof}
The Leibniz rule holds for monomials, $(x^{m+n})'=(m+n)x^{m+n-1}=mx^{m-1}x^n+x^mnx^{n-1}$, and extends by bilinearity.
(i) If $f(a)=0$ write $f=(x-a)g$; then $f'=g+(x-a)g'$ and $f'(a)=g(a)$, which vanishes iff $(x-a)\mid g$.
(ii) Let $d=\gcd(f,f')$ in $F[x]$, $d=uf+vf'$. Then $d$ divides $f$, $f'$ in $K[x]$, and every common divisor in $K[x]$ divides $d$; so $d$ is
also the gcd in $K[x]$. If $f$ has a multiple root $a$ in some $K$, then $x-a$ divides $f$ and $f'$ in $K[x]$, so $d\ne1$. If $d$ is
nonconstant, it has a root $a$ in some extension $K$ (\Cref{thm:rings:adjoin-root}(iii)), and $f(a)=f'(a)=0$.
(iii) $(x^{p^r}-x)'=p^rx^{p^r-1}-1=-1$. If $p=\operatorname{char}F\nmid n$, a common root of $x^n-1$ and $nx^{n-1}$ would be $0$, which is not a root of
$x^n-1$. If $p\mid n$, the derivative is $0$ and $x^n-1=(x^{n/p}-1)^p$ by \Cref{prop:rings:freshman}.
\end{proof}

\begin{theorem}[Gauss's lemma]\label{thm:rings:gauss-lemma}
Let $R$ be a UFD with fraction field $K$. The \emph{content} $c(f)$ of a nonzero $f\in R[x]$ is a gcd of its coefficients (defined up to
units), and $f$ is \emph{primitive} if $c(f)$ is a unit. Every nonzero $g\in K[x]$ is $g=cg_1$ with $c\in K^\times$ and $g_1\in R[x]$ primitive, and $c$
is unique up to $R^\times$.
\begin{enumerate}[label=(\roman*)]
\item $c(fg)=c(f)c(g)$ up to units; a product of primitive polynomials is primitive.
\item If $f\in R[x]$ factors as $f=gh$ in $K[x]$, then $f=g_1h_1$ in $R[x]$ with $g_1\in K^\times g$ and $h_1\in K^\times h$. Hence a primitive $f$ of positive
degree is irreducible in $R[x]$ iff it is irreducible in $K[x]$.
\item $R[x]$ is a UFD. Its prime elements are the primes of $R$ and the primitive polynomials of positive degree irreducible in
$K[x]$. Consequently $\Z[x_1,\dots,x_n]$ and $F[x_1,\dots,x_n]$ ($F$ a field) are UFDs.
\item A monic $f\in\Z[x]$ which is a product of monic polynomials in $\Q[x]$ has those factors in $\Z[x]$. If $a/b$ in lowest terms is a
root of $a_nx^n+\dots+a_0\in\Z[x]$, then $b\mid a_n$ and $a\mid a_0$; a rational root of a monic integer polynomial is an integer.
\end{enumerate}
\end{theorem}

\begin{proof}
For $g\in K[x]$, multiply by a common denominator $b\in R$ and write $bg=c(bg)g_1$; then $g=(c(bg)/b)g_1$. If $cg_1=c'g_1'$ with $g_1,g_1'$
primitive, write $c/c'=a/b$ with $a,b\in R$; then $ag_1=bg_1'$, and taking contents $a=b\cdot$unit, so $c/c'\in R^\times$.
(i) Suppose $f$, $g$ are primitive and $fg$ is not. Some irreducible, hence prime (\Cref{thm:rings:euclid-pid-ufd}(iii)), $p\in R$
divides every coefficient of $fg$. Reduce modulo $p$: $R/p$ is a domain (\Cref{lem:rings:prime-irreducible}(ii)), so $(R/p)[x]$ is a
domain, and $\bar f\bar g=0$ with $\bar f\ne0\ne\bar g$, a contradiction. In general write $f=c(f)f_1$, $g=c(g)g_1$ with $f_1,g_1$ primitive; then
$fg=c(f)c(g)f_1g_1$ with $f_1g_1$ primitive, so the gcd of the coefficients of $fg$ is $c(f)c(g)$ times the gcd of those of $f_1g_1$, a unit:
$c(fg)=c(f)c(g)$ up to units. (For instance $(2x+2)(3x+3)=6x^2+12x+6$ has content $6=2\cdot3$.)
(ii) Write $g=cg_1$, $h=c'h_1$ with $g_1,h_1$ primitive. Then $f=(cc')g_1h_1$ with $g_1h_1$ primitive by (i), and also $f=c(f)f_1$; by
uniqueness, $cc'=c(f)u$ with $u\in R^\times$, so $cc'\in R$. Put $g_1':=cc'g_1\in R[x]$; then $f=g_1'h_1$. For the last claim: if $f$ is primitive and
$f=gh$ in $R[x]$ with $g$ a non-unit, then $g$ is not a constant (a constant factor divides $c(f)$, a unit), so a factorisation in $R[x]$
into non-units has both factors of positive degree; and a factorisation in $K[x]$ into factors of positive degree gives one in
$R[x]$ with the same degrees.
(iii) Existence: write $f=c(f)f_1$ and factor $c(f)$ in $R$. If the primitive $f_1$ is not irreducible, $f_1=gh$ with $g,h$ non-units,
necessarily of positive degree as just seen and primitive (as $c(g)c(h)=c(f_1)$ is a unit); induction on the degree. Primes: if $p$
is prime in $R$, then $R[x]/pR[x]\simeq(R/p)[x]$ is a domain, so $p$ is prime in $R[x]$. If $f$ is primitive of positive degree and
irreducible in $R[x]$, it is irreducible in $K[x]$ by (ii), hence prime there ($K[x]$ is a PID). If $f\mid gh$ in $R[x]$, then $f$ divides $g$,
say, in $K[x]$: $g=fk$ with $k=ck_1$, $k_1$ primitive; by (i) $fk_1$ is primitive, so $c$ equals $c(g)$ up to a unit and lies in $R$; hence
$k\in R[x]$ and $f\mid g$ in $R[x]$. These two kinds exhaust the irreducibles of $R[x]$: an irreducible of degree $0$ is an irreducible
element of $R$ (a factorisation in $R$ is one in $R[x]$, and $R[x]^\times=R^\times$), hence prime in $R$; an irreducible $f$ of positive degree
is primitive, since otherwise $f=c(f)f_1$ would be a factorisation into two non-units. So every irreducible is prime and uniqueness
follows from \Cref{lem:rings:prime-irreducible}(iii). The polynomial rings in several variables follow by induction,
$R[x_1,\dots,x_n]=R[x_1,\dots,x_{n-1}][x_n]$.
(iv) By (ii), $f=g_1h_1$ in $\Z[x]$ with $g_1=\lambda g$, $h_1=\lambda'h$; comparing leading coefficients, $\lambda\lambda'=1$ with $\lambda,\lambda'\in\Z$ the leading
coefficients of $g_1,h_1$, so $\lambda=\pm1$ and $g=\pm g_1\in\Z[x]$. If $f(a/b)=0$ then $\sum a_ia^ib^{n-i}=0$, so $b\mid a_na^n$ and $a\mid a_0b^n$; since
$\gcd(a,b)=1$, Euclid's lemma gives $b\mid a_n$ and $a\mid a_0$.
\end{proof}

\begin{theorem}[Eisenstein's criterion]\label{thm:rings:eisenstein}
Let $R$ be a domain, $P$ a prime ideal and $f=a_nx^n+\dots+a_0\in R[x]$, $n\ge1$, with $a_n\notin P$, $a_i\in P$ for $i<n$, and $a_0\notin P^2$. Then $f$ is
not a product of two polynomials of positive degree in $R[x]$. If $R$ is a UFD with fraction field $K$ and $f$ is primitive, then $f$
is irreducible in $R[x]$ and in $K[x]$. Examples: $x^n-p$; $x^2+3x+3=\Phi_3(x+1)$; $\Phi_p(x+1)=((x+1)^p-1)/x$ and $\Phi_{p^k}(x+1)$
(\Cref{thm:rings:cyclotomic}(vi)).
\end{theorem}

\begin{proof}
Suppose $f=gh$ with $\deg g=r\ge1$, $\deg h=s\ge1$, $r+s=n$. Let $k$ be the fraction field of the domain $R/P$ and reduce modulo $P$:
$\bar g\bar h=\bar a_nx^n$ in $k[x]$ with $\bar a_n\ne0$. Then $\deg\bar g=r$ and $\deg\bar h=s$ (their degrees are at most $r$, $s$ and add up to $n$). In
the UFD $k[x]$ the prime $x$ and units are the only factors of $\bar a_nx^n$, so $\bar g=bx^r$ and $\bar h=cx^s$. As $r,s\ge1$, the constant terms
$g(0)$, $h(0)$ lie in $P$, and $a_0=g(0)h(0)\in P^2$, a contradiction. If $R$ is a UFD and $f$ is primitive, a factorisation $f=gh$ into non-units
has both factors of positive degree (a constant non-unit factor would divide $c(f)$), so $f$ is irreducible in $R[x]$, and in $K[x]$
by \Cref{thm:rings:gauss-lemma}(ii). The examples: $x^n-p$ at $P=p\Z$; $\Phi_p(x+1)=\sum_{i=1}^p\binom pix^{i-1}$ has leading coefficient $1$,
the other coefficients divisible by $p$ (\Cref{prop:rings:freshman}(i)) and constant term $\binom p1=p$.
\end{proof}

For $p=5$: $\Phi_5(x+1)=\frac{(x+1)^5-1}x=x^4+5x^3+10x^2+10x+5$, Eisenstein at $5$. The condition $a_0\notin P^2$ cannot be dropped:
$x^2-4=(x-2)(x+2)$ has all lower coefficients in $2\Z$ but $a_0=-4\in4\Z$. The substitution $x\mapsto x+1$ is often needed: $\Phi_5$ itself
has all coefficients equal to $1$ and no prime divides them, and the shift replaces the roots $\zeta_5^a$ by the numbers $\zeta_5^a-1$, which are all
multiples of $\zeta_5-1$ in $\Z[\zeta_5]$, since $\zeta_5^a-1=(\zeta_5-1)(1+\zeta_5+\dots+\zeta_5^{a-1})$; their product is $\Phi_5(1)=5$ up to sign.

\begin{proposition}[reduction modulo $p$]\label{prop:rings:reduction-mod-p}
Let $f\in\Z[x]$ have degree $n$ and let $p$ be a prime not dividing the leading coefficient. If $f=gh$ in $\Z[x]$, then $\bar f=\bar g\bar h$ in $\F_p[x]$
with $\deg\bar g=\deg g$ and $\deg\bar h=\deg h$. Hence:
\begin{enumerate}[label=(\roman*)]
\item if $f$ is primitive and $\bar f$ is irreducible in $\F_p[x]$, then $f$ is irreducible in $\Q[x]$ (and in $\Z[x]$);
\item the degrees of the factors in any factorisation of $f$ in $\Q[x]$ are sums of degrees of irreducible factors of $\bar f$
(disjoint sub-multisets); so if two primes give incompatible patterns, $f$ is irreducible over $\Q$.
\end{enumerate}
\end{proposition}

\begin{proof}
Reduction $\Z[x]\to\F_p[x]$ is a ring homomorphism. The leading coefficients of $g$ and $h$ multiply to that of $f$, which is
nonzero modulo $p$, so neither leading coefficient vanishes modulo $p$. (i) A factorisation in $\Q[x]$ into factors of positive
degree gives one in $\Z[x]$ with the same degrees (\Cref{thm:rings:gauss-lemma}(ii)), and then one of $\bar f$ with the same degrees.
(ii) Write $f=g_1\cdots g_k$ in $\Q[x]$, hence (Gauss) in $\Z[x]$ with the same degrees; then $\bar f=\bar g_1\cdots\bar g_k$, $\deg\bar g_i=\deg g_i$, and by
unique factorisation in $\F_p[x]$ the irreducible factors of $\bar f$ are distributed among the $\bar g_i$.
\end{proof}

\begin{warning}[$x^4+1$]\label{warn:rings:x4plus1}
The test works in one direction only. $x^4+1$ is irreducible over $\Q$: $(x+1)^4+1=x^4+4x^3+6x^2+4x+2$ is Eisenstein at $2$. But it is
reducible modulo every prime $p$. For $p=2$, $x^4+1=(x+1)^4$. For odd $p$, $8\mid p^2-1$, so the cyclic group $\F_{p^2}^\times$
(\Cref{thm:rings:finite-field-exists,prop:rings:finite-field-basic}) contains an element $\zeta$ of order $8$, and $\zeta^4=-1$. If $x^4+1$
were irreducible over $\F_p$, then $\F_p[x]/(x^4+1)\simeq\F_p[\zeta]\subseteq\F_{p^2}$ (\Cref{thm:rings:adjoin-root}(ii)) would have $p^4\le p^2$ elements.
(Checked with PARI for all $p<500$: every factor has degree at most $2$.)
\end{warning}

\begin{example}[the $3$-division polynomial of 11a1]\label{ex:rings:psi3}
For $E=\curve{11a1}$: $y^2+y=x^3-x^2-10x-20$, the $3$-division polynomial, whose roots are the $x$-coordinates of the points of order $3$
(\Cref{ch:ellcurves}), is $\psi_3=3x^4-4x^3-60x^2-237x-21$. It is primitive, $5$ does not divide its leading coefficient, and modulo $5$
it is $3(x^4+2x^3+x+3)$, which is irreducible in $\F_5[x]$: the monic factor has no root (its values at $0,1,2,3,4$ are $3,2,2,1,1$), and it is not
$(x^2+ax+b)(x^2+cx+d)$, since $bd=3$ leaves $(b,d)\in\{(1,3),(3,1),(2,4),(4,2)\}$, and in each case $a+c=2$, $ad+bc=1$ determine $a,c$ and give
$b+d+ac\in\{4,4,2,2\}$, never $0$. By \Cref{prop:rings:reduction-mod-p}(i), $\psi_3$ is irreducible over $\Q$.
\begin{lstlisting}[language=GP]
E = ellinit([0,-1,1,-10,-20]);      \\ 11a1
f = elldivpol(E, 3)                  \\ 3*x^4 - 4*x^3 - 60*x^2 - 237*x - 21
factormod(f, 5)                      \\ a single factor x^4 + 2*x^3 + x + 3
polisirreducible(f)                  \\ 1
\end{lstlisting}
\end{example}

\begin{paperlink}\label{pl:rings:ns-divpol}
Proposition prop:ns ends: "The $3$-division polynomial of $E$ then factors as a product of two irreducible quadratics,
and otherwise it is irreducible", and Section sec:anomaly checks that the $3$-division polynomial of every violating curve (all of
image $\Ns$) factors as two quadratics.
\Cref{prop:rings:reduction-mod-p} is how such a statement is tested for a given curve: one prime $p$ with $\bar\psi_3$ irreducible proves
irreducibility, as for 11a1, whose mod $3$ image is $\GL_2(\F_3)$. The converse direction (that a factorisation over $\Q$ is
detected) needs the factorisation over $\Q$ itself, which PARI computes; the Galois-theoretic meaning of the pattern is
\Cref{ch:galois}.
\end{paperlink}

\begin{theorem}[Kronecker's construction]\label{thm:rings:adjoin-root}
Let $F$ be a field and $f\in F[x]$ irreducible of degree $d$.
\begin{enumerate}[label=(\roman*)]
\item $F[x]/(f)$ is a field containing (a copy of) $F$, of dimension $d$ over $F$ with basis $1,\bar x,\dots,\bar x^{d-1}$, and $f(\bar x)=0$.
\item If $\alpha$ is a root of $f$ in a field $K\supseteq F$, then $x\mapsto\alpha$ induces an isomorphism $F[x]/(f)\simeq F[\alpha]$; in particular $F[\alpha]$ is a
field of dimension $d$ over $F$.
\item For every nonconstant $g\in F[x]$ there is a field $K\supseteq F$, of finite dimension over $F$, in which $g$ is a product of linear factors.
\end{enumerate}
\end{theorem}

\begin{proof}
(i) $(f)$ is maximal (\Cref{prop:rings:poly-division}(ii)), so the quotient is a field. $F\to F[x]/(f)$ is injective, as a nonzero constant is not
divisible by $f$. Division by $f$ (made monic) shows that $1,\bar x,\dots,\bar x^{d-1}$ span, and they are independent because a nonzero
polynomial of degree $<d$ is not divisible by $f$. Clearly $f(\bar x)=\overline{f(x)}=0$.
(ii) Evaluation $F[x]\to K$, $x\mapsto\alpha$, has image $F[\alpha]$ and a kernel containing $f$, not equal to $F[x]$; as $(f)$ is maximal, the kernel is
$(f)$. Apply \Cref{thm:rings:isomorphism}(i).
(iii) Induction on $\deg g$. Let $f$ be an irreducible factor of $g$ and $K_1=F[x]/(f)$, in which $f$, hence $g$, has a root $\alpha_1$; then
$g=(x-\alpha_1)g_1$ in $K_1[x]$ (\Cref{prop:rings:poly-division}(iii)). By induction $g_1$ splits in some $K\supseteq K_1$ with $\dim_{K_1}K$ finite, and
$\dim_FK=\dim_{K_1}K\cdot\dim_FK_1$ (a basis of $K$ over $F$ is given by the products of bases).
\end{proof}

In (i) the inverse of a nonzero class is computed by Bezout's identity: if $g\notin(f)$ then $\gcd(f,g)=1$ ($f$ is irreducible), so $uf+vg=1$
for some $u,v\in F[x]$, and $\bar v=\bar g^{-1}$. For $F=\Q$ and $f=x^2-2$ this gives the familiar formula: in $\Q[x]/(x^2-2)\simeq\Q(\sqrt2)$, with
$t=\bar x$,
\[
(a+bt)(a-bt)=a^2-2b^2,\qquad\text{so}\qquad(a+bt)^{-1}=\frac{a-bt}{a^2-2b^2},
\]
and $a^2-2b^2\ne0$ for rational $(a,b)\ne(0,0)$ precisely because $x^2-2$ has no rational root. The point of the construction is that it
needs no ambient field such as $\C$: over $\F_p$, where there is no such field to start from, it is the only way to produce a root.

\begin{example}[$\F_9$ and the paper's matrix $J$]\label{ex:rings:f9}
$x^2+1$ has no root in $\F_3$ (its values at $0,1,2$ are $1,2,2$), so it is irreducible and $\F_9:=\F_3[x]/(x^2+1)=\F_3[i]$, $i=\bar x$, is a field with $9$
elements (by \Cref{thm:rings:finite-field-exists} every field with $9$ elements is isomorphic to it). In the basis $1,i$ of $\F_9$ over $\F_3$,
multiplication by $-i$ sends $1\mapsto-i$ and $i\mapsto1$, so its matrix is the paper's $J=\sm01{-1}0$ of Proposition prop:ns: $J=m_{-i}$, as in
the standard models of \Cref{ch:groups}. Its characteristic polynomial $X^2+1$ is irreducible over $\F_3$, so $J$ has no eigenvalue in
$\F_3$ and the eigenvalues $\pm i$ in $\F_9$. This is why the character $\psi$ of order $4$ in Proposition prop:ns takes values in $\F_9^\times$ and
why the paper writes $\bar\rho_{E,3}\otimes\F_9\simeq\Ind\psi$ (over $\F_3$ itself $\bar\rho_{E,3}$ is irreducible). More generally
$\F_q[x]/(x^2-c)$, $c$ a non-square, is the field $\F_q[\sqrt c]$ with $q^2$ elements used in \Cref{ch:groups}.
\end{example}

\section[The ring \texorpdfstring{$\Z/n\Z$}{Z/nZ}, units and primitive roots]{The ring \texorpdfstring{$\Z/n\Z$}{Z/nZ}, the Chinese remainder theorem, units and primitive roots}\label{sec:rings:zn}

The groups $(\Z/n\Z)^\times$ are where Dirichlet characters live (\Cref{sec:rings:dirichlet}), and the groups
$(\Z/\ell^k\Z)^\times=\Z_\ell^\times/(1+\ell^k\Z_\ell)$ ($k\ge1$) are the finite levels of $\Z_\ell^\times$, the image of the cyclotomic character. The paper uses one fact about them in an essential way:
in Lemma lem:h0, an isomorphism $\rho\bmod\ell^k\simeq(\rho\otimes\chi_\ell)\bmod\ell^k$ forces $\chi_\ell^2\equiv1\pmod{\ell^k}$, and this is possible
for $k\ge1$ only when $\ell^k=3$ (\Cref{cor:rings:squares-one}). The lifting-the-exponent lemma computes valuations such as $v_\ell(q^2-1)$ in the
table of Section 1.

\begin{proposition}\label{prop:rings:zn}
Let $n\ge1$.
\begin{enumerate}[label=(\roman*)]
\item $a\bmod n$ is a unit of $\Z/n\Z$ iff $\gcd(a,n)=1$. For $n\ge2$, $\Z/n\Z$ is a field iff it is a domain iff $n$ is prime; we write
$\F_p=\Z/p\Z$.
\item $\#(\Z/n\Z)^\times=\varphi(n)$, the number of $1\le a\le n$ prime to $n$, and $a^{\varphi(n)}\equiv1\pmod n$ for $\gcd(a,n)=1$ (Euler). Moreover
$\varphi(p^k)=p^{k-1}(p-1)$ and $\sum_{d\mid n}\varphi(d)=n$.
\item The map $a\mapsto(x\mapsto ax)$ is an isomorphism $(\Z/n\Z)^\times\simeq\Aut(\Z/n\Z)$ (automorphisms of the additive group).
\end{enumerate}
For $k\ge2$, $\Z/p^k\Z$ is not a field ($p$ is nilpotent in it), and in particular it is not $\F_{p^k}$.
\end{proposition}

\begin{proof}
(i) If $\gcd(a,n)=1$, Bezout gives $ua+vn=1$, so $ua\equiv1$. If $d=\gcd(a,n)>1$ then $a\cdot(n/d)\equiv0$ with $n/d\not\equiv0$, so $a$ is a zero divisor
and not a unit. If $n=p$ is prime, every nonzero class is a unit; if $n=ab$ with $1<a,b<n$, then $\bar a\bar b=0$ and $\Z/n$ is not a domain.
A field is a domain.
(ii) The first claim is (i); Euler's congruence is Lagrange's theorem in the group $(\Z/n)^\times$ (\Cref{ch:groups}). Among $1,\dots,p^k$
exactly the $p^{k-1}$ multiples of $p$ are not prime to $p^k$. For the sum, $\Z/n$ is cyclic of order $n$ and has, for each $d\mid n$,
exactly $\varphi(d)$ elements of order $d$ (the generators of its unique subgroup of order $d$, \Cref{ch:groups}); count the
elements by their order.
(iii) An endomorphism of the cyclic group $\Z/n$ is determined by the image $a$ of $1$, and it is $x\mapsto ax$; it is bijective iff $a$
generates, iff $\gcd(a,n)=1$. Composition corresponds to multiplication.
\end{proof}

For $n=12$: $(\Z/12\Z)^\times=\{1,5,7,11\}$, so $\varphi(12)=4$, and each of these units squares to $1$ ($25=2\cdot12+1$, $49=4\cdot12+1$, $121=10\cdot12+1$), so
$(\Z/12\Z)^\times$ is not cyclic. The divisors $1,2,3,4,6,12$ of $12$ have $\varphi=1,1,2,2,2,4$, with sum $12$. The class $2$ is a zero divisor, as in the
proof of (i): $2\cdot6\equiv0$.

\begin{theorem}[Chinese remainder theorem]\label{thm:rings:crt}
For $m,n\ge1$ the reduction map $\Z/mn\Z\to\Z/m\Z\times\Z/n\Z$ is a ring isomorphism iff $\gcd(m,n)=1$. Hence for $n=\prod_ip_i^{k_i}$ the reduction maps
give ring isomorphisms
\[
\Z/n\Z\simeq\prod_i\Z/p_i^{k_i}\Z,\qquad(\Z/n\Z)^\times\simeq\prod_i(\Z/p_i^{k_i}\Z)^\times,
\]
and $\varphi(mn)=\varphi(m)\varphi(n)$ for coprime $m,n$, so $\varphi(n)=\prod_ip_i^{k_i-1}(p_i-1)$.
\end{theorem}

\begin{proof}
If $\gcd(m,n)=1$ then $m\Z+n\Z=\Z$ and $m\Z\cdot n\Z=mn\Z$ (\Cref{prop:rings:ideals}(iii)), and \Cref{prop:rings:two-ideals}(ii) applies; (iii) of that
proposition gives the decomposition over the prime powers $p_i^{k_i}$, which are pairwise coprime. If $d=\gcd(m,n)>1$, every element
of $\Z/m\times\Z/n$ has additive order dividing $\lcm(m,n)=mn/d<mn$, while $1\in\Z/mn$ has order $mn$. The unit group of a product of
rings is the product of the unit groups. (The group-theoretic version is in \Cref{ch:groups}.)
\end{proof}

The inverse of the reduction map is explicit once Bezout's identity $um+vn=1$ is known: the pair $(y\bmod m,z\bmod n)$ comes from
$x=y\cdot vn+z\cdot um$, since $vn\equiv1\pmod m$, $vn\equiv0\pmod n$, and symmetrically for $um$. For $m=3$, $n=5$ this is the computation after
\Cref{prop:rings:two-ideals}. For $\gcd(m,n)>1$ the failure is visible already for $m=n=2$: $\Z/4\Z$ is cyclic of order $4$, while every element of
$\Z/2\Z\times\Z/2\Z$ has order at most $2$. The unit statement gives, for example, $(\Z/15\Z)^\times\simeq(\Z/3\Z)^\times\times(\Z/5\Z)^\times$, of order $2\cdot4=8=\varphi(15)$.

\begin{lemma}[lifting the exponent]\label{lem:rings:lte}
\begin{enumerate}[label=(\roman*)]
\item Let $\ell$ be an odd prime and $a,b\in\Z$ with $\ell\mid a-b$ and $\ell\nmid ab$. Then $v_\ell(a^n-b^n)=v_\ell(a-b)+v_\ell(n)$ for every $n\ge1$.
\item If $a$ is odd and $n\ge2$ is even, then $v_2(a^n-1)=v_2(a-1)+v_2(a+1)+v_2(n)-1$.
\item For $\ell$ odd and $j\ge0$, $v_\ell\bigl((1+\ell)^{\ell^j}-1\bigr)=j+1$; hence $1+\ell$ has order $\ell^{k-1}$ in $(\Z/\ell^k\Z)^\times$.
\item Let $q\ne\ell$ be primes, $\ell$ odd, and $d$ the order of $q$ modulo $\ell$. Then $v_\ell(q^f-1)=0$ if $d\nmid f$, and
$v_\ell(q^f-1)=v_\ell(q^d-1)+v_\ell(f/d)$ if $d\mid f$.
\item $v_\ell(q^2-1)=v_\ell(q-1)+v_\ell(q+1)$, and for $\ell$ odd at most one of the two terms is positive.
\end{enumerate}
\end{lemma}

\begin{proof}
(i) If $\ell\nmid n$: $a^n-b^n=(a-b)S$ with $S=\sum_{i=0}^{n-1}a^{n-1-i}b^i\equiv na^{n-1}\not\equiv0\pmod\ell$, as $a\equiv b$. If $n=\ell$: write $a=b+t$ with $\ell\mid t$.
By the binomial theorem, $a^{j}=(b+t)^{j}=b^{j}+jtb^{j-1}+(\text{terms divisible by }t^2)$, so modulo $t^2$,
$a^{\ell-1-i}b^i\equiv b^{\ell-1}+(\ell-1-i)tb^{\ell-2}$, and summing over $i$ (with $\sum_{i=0}^{\ell-1}(\ell-1-i)=0+1+\dots+(\ell-1)$),
\[
S=\sum_{i=0}^{\ell-1}a^{\ell-1-i}b^i\equiv\ell b^{\ell-1}+tb^{\ell-2}\sum_{i=0}^{\ell-1}(\ell-1-i)=\ell b^{\ell-1}+tb^{\ell-2}\frac{\ell(\ell-1)}2\pmod{t^2}.
\]
Since $\ell$ is odd, $\ell\mid\frac{\ell(\ell-1)}2$, so the second term and $t^2$ are divisible by $\ell^2$, and $S\equiv\ell b^{\ell-1}\pmod{\ell^2}$: $v_\ell(S)=1$. For
general $n=\ell^sm$, $\ell\nmid m$, apply the first case to $a,b$ and $m$: $v_\ell(a^m-b^m)=v_\ell(a-b)$. Then apply the second case $s$ times,
to the pairs $a^{m\ell^j}$, $b^{m\ell^j}$ for $j=0,\dots,s-1$. These still satisfy the hypotheses: $a^{m\ell^j}\equiv b^{m\ell^j}\pmod\ell$ because $a\equiv b$, and
neither is divisible by $\ell$ because $\ell\nmid ab$. Each application adds exactly $1$ to the valuation, since
$a^{m\ell^{j+1}}-b^{m\ell^{j+1}}=(a^{m\ell^j})^\ell-(b^{m\ell^j})^\ell$; so $v_\ell(a^n-b^n)=v_\ell(a-b)+s$.
(ii) Write $n=2^sm$ with $m$ odd and $s\ge1$. As in (i), $a^m-1=(a-1)(a^{m-1}+\dots+1)$ and $a^m+1=(a+1)(a^{m-1}-a^{m-2}+\dots+1)$, the
second factors being sums of $m$ odd terms, hence odd; so $v_2(a^m-1)=v_2(a-1)$, $v_2(a^m+1)=v_2(a+1)$, and
$v_2(a^{2m}-1)=v_2\bigl((a^m-1)(a^m+1)\bigr)=v_2(a-1)+v_2(a+1)$. For odd $x$, $x^2\equiv1\pmod8$ (as $x^2-1=(x-1)(x+1)$ is a product of two consecutive
even numbers, one divisible by $4$), so $x^2+1\equiv2\pmod8$ and $v_2(x^2+1)=1$. For $j\ge2$ put $x=a^{2^{j-2}m}$, which is odd; then
$a^{2^jm}-1=(x^2-1)(x^2+1)$, so each further squaring adds exactly $1$: $v_2(a^{2^jm}-1)=v_2(a^{2^{j-1}m}-1)+1$. After the $s-1$ squarings from
$2m$ to $2^sm=n$ the valuation is $v_2(a-1)+v_2(a+1)+(s-1)$.
(iii) Apply (i) with $a=1+\ell$, $b=1$: $v_\ell((1+\ell)^{\ell^j}-1)=1+j$. So $(1+\ell)^{\ell^{k-1}}\equiv1\pmod{\ell^k}$ but $(1+\ell)^{\ell^{k-2}}\not\equiv1\pmod{\ell^k}$ for
$k\ge2$; as the order divides $\ell^{k-1}$, it equals $\ell^{k-1}$.
(iv) $\ell\mid q^f-1$ iff $q^f\equiv1\pmod\ell$ iff $d\mid f$. If $d\mid f$, apply (i) to $a=q^d$, $b=1$, $n=f/d$.
(v) $q^2-1=(q-1)(q+1)$, and an odd prime cannot divide both $q-1$ and $q+1$, whose difference is $2$.
\end{proof}

Two checks. With $\ell=3$, $a=4$, $b=1$, $n=9$: (i) predicts $v_3(4^9-1)=v_3(3)+v_3(9)=3$, and indeed $4^9-1=262143=3^3\cdot7\cdot19\cdot73$. At $\ell=2$ the
naive formula fails, which is why (ii) has the extra term: $v_2(3^2-1)=v_2(8)=3$, not $v_2(3-1)+v_2(2)=2$; (ii) gives $1+2+1-1=3$, and for $n=4$ it gives
$v_2(3^4-1)=v_2(80)=1+2+2-1=4$. The failure comes from the step $n=\ell$ of the proof: for $\ell=2$ the number $\frac{\ell(\ell-1)}2=1$ is not divisible by $\ell$.

\begin{theorem}[primitive roots]\label{thm:rings:primitive-roots}
\begin{enumerate}[label=(\roman*)]
\item $(\Z/p\Z)^\times=\F_p^\times$ is cyclic of order $p-1$; a generator is a \emph{primitive root} modulo $p$.
\item For $p$ odd and $k\ge1$, $(\Z/p^k\Z)^\times$ is cyclic of order $p^{k-1}(p-1)$. More precisely, if $g$ is a primitive root modulo $p$ with
$g^{p-1}\not\equiv1\pmod{p^2}$ (true for $g$ or for $g+p$), then $g$ generates $(\Z/p^k\Z)^\times$ for every $k$.
\item $(\Z/2\Z)^\times$ and $(\Z/4\Z)^\times$ are cyclic. For $k\ge3$, $(\Z/2^k\Z)^\times=\{\pm1\}\times\langle5\rangle\simeq\Z/2\times\Z/2^{k-2}$, which is not cyclic.
\item $(\Z/n\Z)^\times$ is cyclic iff $n\in\{1,2,4,p^k,2p^k\}$ with $p$ an odd prime and $k\ge1$.
\end{enumerate}
\end{theorem}

\begin{proof}
(i) A finite subgroup of the multiplicative group of a field is cyclic (\Cref{ch:groups}; the proof uses
\Cref{prop:rings:poly-division}(iv)).
(ii) Let $g$ be a primitive root modulo $p$. If $g^{p-1}\equiv1\pmod{p^2}$, then
$(g+p)^{p-1}\equiv g^{p-1}+(p-1)pg^{p-2}\equiv1-pg^{p-2}\not\equiv1\pmod{p^2}$, and $g+p$ is also a primitive root modulo $p$. So assume $g^{p-1}=1+pt$
with $p\nmid t$. The order $o$ of $g$ modulo $p^k$ is a multiple of $p-1$ (its order modulo $p$) and divides $\varphi(p^k)=p^{k-1}(p-1)$, so
$o=(p-1)p^j$ with $j\le k-1$. By \Cref{lem:rings:lte}(i), $v_p(g^{(p-1)p^j}-1)=v_p(g^{p-1}-1)+j=1+j$, which is $\ge k$ only if $j\ge k-1$. So
$o=\varphi(p^k)$.
(iii) $(\Z/2)^\times=\{1\}$ and $(\Z/4)^\times=\{\pm1\}$. Let $k\ge3$. By \Cref{lem:rings:lte}(ii), $v_2(5^{2^j}-1)=2+1+j-1=j+2$ for $j\ge1$, and $v_2(5-1)=2$;
so $5^{2^j}\equiv1\pmod{2^k}$ iff $j\ge k-2$, and $5$ has order $2^{k-2}$. All powers of $5$ are $\equiv1\pmod4$, so $-1\notin\langle5\rangle$; the subgroup
$\{\pm1\}\times\langle5\rangle$ has $2^{k-1}=\varphi(2^k)$ elements and is everything. It has the two distinct elements $-1$ and $5^{2^{k-3}}$ of order $2$,
while a cyclic group has at most one.
(iv) By \Cref{thm:rings:crt}, $(\Z/n)^\times\simeq\prod_i(\Z/p_i^{k_i})^\times$. For odd $p$ the factor has even order $p^{k-1}(p-1)$, and $(\Z/2^k)^\times$ has
even order for $k\ge2$. If two factors have even order, each contains an element of order $2$, and the product contains a
subgroup $\Z/2\times\Z/2$, impossible in a cyclic group; if $8\mid n$, the factor $(\Z/2^k)^\times$ is not cyclic and neither is a group
containing it. This excludes every $n$ outside the list. Conversely, $(\Z/2p^k)^\times\simeq(\Z/2)^\times\times(\Z/p^k)^\times\simeq(\Z/p^k)^\times$, and the other cases
are (ii), (iii). (Checked numerically for all $n<200$.)
\end{proof}

The proviso in (ii) is real. The class of $7$ is a primitive root modulo $5$ (it is $\equiv2$, of order $4$), but $7^4=2401=96\cdot25+1\equiv1\pmod{25}$, so $7$ has order $4$
modulo $25$ and does not generate the group of order $20$; the replacement $7+5=12$ has $12^4\equiv11\pmod{25}$ and is a generator. On the other
hand $2$ is a primitive root modulo $3$ with $2^2=4\not\equiv1\pmod9$, so $2$ generates $(\Z/3^k\Z)^\times$ for every $k$: its orders modulo $9$ and $27$ are $6$
and $18$. (Orders computed with PARI's \texttt{znorder}.)

\begin{corollary}\label{cor:rings:squares-one}
Let $\ell$ be an odd prime and $k\ge1$.
\begin{enumerate}[label=(\roman*)]
\item The solutions of $a^2\equiv1\pmod{\ell^k}$ are $a\equiv\pm1$.
\item $a^2\equiv1\pmod{\ell^k}$ holds for every $a$ prime to $\ell$ iff $\varphi(\ell^k)=2$ iff $\ell^k=3$.
\item If $\chi\colon G\to(\Z/\ell^k)^\times$ is a homomorphism from a group with $\chi(g)^2=1$ for all $g$, and either $\chi$ is surjective or its image
contains the class of $2$, then $\ell^k=3$. Example: $2^2=4\not\equiv1\pmod9$, and $4\not\equiv1\pmod\ell$ for $\ell\ge5$.
\end{enumerate}
\end{corollary}

\begin{proof}
(i) $(\Z/\ell^k)^\times$ is cyclic of even order (\Cref{thm:rings:primitive-roots}(ii)), so it has exactly $\gcd(2,\varphi(\ell^k))=2$ elements with $x^2=1$
(\Cref{ch:groups}): $1$ and $-1$, distinct as $\ell^k\ge3$. (Directly: $\ell^k\mid(a-1)(a+1)$, and the odd prime $\ell$ cannot divide both factors,
whose difference is $2$, so $\ell^k$ divides one of them.) (ii) By (i), every unit squares to $1$ iff there are only two units. And
$\varphi(\ell^k)=\ell^{k-1}(\ell-1)=2$ iff $\ell=3$, $k=1$. (iii) Every element of the image squares to $1$, so it is $\pm1$ by (i). If $\chi$ is
surjective, $(\Z/\ell^k)^\times=\{\pm1\}$ and $\ell^k=3$ by (ii). If $2$ is in the image, $2\equiv\pm1\pmod{\ell^k}$, so $\ell^k\mid3$.
\end{proof}

\begin{paperlink}\label{pl:rings:h0}
Lemma lem:h0 (in the revised paper, following the correction recorded in the book's notes on the paper): an element of
exact order $\ell^k$ of $H^0(\Q,\ad T(1)\otimes\Q_\ell/\Z_\ell)$ gives an isomorphism $\rho\bmod\ell^k\simeq(\rho\otimes\chi_\ell)\bmod\ell^k$, and comparing
determinants, $\chi_\ell(g)^2\equiv1\pmod{\ell^k}$ for all $g\in G_\Q$. The cyclotomic character maps $G_\Q$ onto $\Z_\ell^\times$ (\Cref{ch:galois}), hence onto
$(\Z/\ell^k)^\times$; so \Cref{cor:rings:squares-one}(iii) gives $\ell^k\in\{1,3\}$: "Since $\chi_\ell(G_\Q)=\Z_\ell^\times$ contains $2$, and $4\not\equiv1$
modulo $\ell$ for $\ell\ge5$ and modulo $9$, this forces $k=0$ for $\ell\ge5$ and $k\le1$ for $\ell=3$." Hence $h_\ell(E)=0$ for $\ell\ge5$ and
$h_3(E)\le1$. The case $\ell^k=3$ is the identity $\bar\chi_3^{\mathrm{cyc}}=\chi_{-3}$, whose Dirichlet side is \Cref{prop:rings:paper-characters}(iii).
\end{paperlink}

\section{Finite fields}\label{sec:rings:finite-fields}

In this section and the next $q=p^r$ is a prime power (\Cref{conv:rings:primes}). Finite fields are the residue fields of
everything in the book: $\F_q$ at a bad prime $q$ of $E$, $\F_\ell$ and $\F_9$ as coefficient fields of $\bar\rho$, and $\F_{q^f}$ in the
unramified extensions of \Cref{ch:local}. We prove existence and uniqueness, the Frobenius automorphism and the subfields,
count irreducible polynomials, and determine the roots of unity in $\F_q$; the last item is the arithmetic input of Lemma
lem:euler. That $\Gal(\F_{q^n}/\F_q)$ is cyclic generated by $x\mapsto x^q$ is proved in \Cref{ch:galois}.

\begin{proposition}\label{prop:rings:finite-field-basic}
Let $F$ be a finite field.
\begin{enumerate}[label=(\roman*)]
\item $\operatorname{char}F=p$ is a prime, $F$ contains the image $\F_p$ of $\Z$, and $\#F=p^r$ with $r=\dim_{\F_p}F$.
\item $F^\times$ is cyclic of order $p^r-1$.
\item Every $a\in F$ satisfies $a^{p^r}=a$, and $x^{p^r}-x=\prod_{a\in F}(x-a)$ in $F[x]$.
\item $F=\F_p[\alpha]$ for every generator $\alpha$ of $F^\times$.
\end{enumerate}
\end{proposition}

\begin{proof}
(i) $\Z\to F$ is not injective, as $F$ is finite, so $\operatorname{char}F$ is a prime $p$ (\Cref{prop:rings:characteristic}); the image is $\Z/p=\F_p$, and $F$ is a
finite-dimensional $\F_p$-vector space. (ii) is the cyclicity of finite subgroups of $K^\times$ (\Cref{ch:groups}). (iii) For $a\ne0$, Lagrange
gives $a^{p^r-1}=1$; and $0^{p^r}=0$. So the $p^r$ distinct linear polynomials $x-a$ divide $x^{p^r}-x$; they are pairwise coprime, so their
product, monic of degree $p^r$, divides it and equals it. (iv) $\F_p[\alpha]$ contains $0$ and all powers of $\alpha$.
\end{proof}

\begin{theorem}[existence and uniqueness of finite fields]\label{thm:rings:finite-field-exists}
Let $q=p^r$.
\begin{enumerate}[label=(\roman*)]
\item There is a field with $q$ elements.
\item Any two fields with $q$ elements are isomorphic. We write $\F_q$ for any of them.
\item $\F_q\simeq\F_p[x]/(f)$ for every irreducible $f\in\F_p[x]$ of degree $r$ (such $f$ exist by \Cref{prop:rings:irreducible-count}).
\item The isomorphism in (ii) is not unique when $r>1$ (compose with powers of the Frobenius, \Cref{prop:rings:frobenius}); inside a fixed
algebraic closure of $\F_p$ the subfield with $q$ elements is unique (\Cref{ch:galois}).
\end{enumerate}
\end{theorem}

\begin{proof}
(i) Let $K\supseteq\F_p$ be a field in which $x^q-x$ splits (\Cref{thm:rings:adjoin-root}(iii)). Its roots are $q$ distinct elements, since
$x^q-x$ is separable (\Cref{prop:rings:derivative}(iii)). They form a subfield: they contain $0$, $1$; if $a^q=a$ and $b^q=b$ then
$(a\pm b)^q=a^q\pm b^q=a\pm b$ (\Cref{prop:rings:freshman}; for $p=2$, $-b=b$; for odd $p$, $(-b)^q=-b^q$), $(ab)^q=ab$, and $(a^{-1})^q=a^{-1}$.
(ii) Let $K$, $K'$ have $q$ elements. Write $K=\F_p[\alpha]$ (\Cref{prop:rings:finite-field-basic}(iv)). The kernel of $\F_p[x]\to K$, $x\mapsto\alpha$, is a nonzero
($\F_p[x]$ is infinite) prime ideal, so it is $(f)$ with $f$ monic irreducible, and $K\simeq\F_p[x]/(f)$. Since $\alpha^q=\alpha$, $f\mid x^q-x$. In $K'[x]$,
$x^q-x=\prod_{b\in K'}(x-b)$, so by unique factorisation in $K'[x]$ the divisor $f$ is a product of some of these linear factors and has a
root $\beta\in K'$. By \Cref{thm:rings:adjoin-root}(ii), $\F_p[\beta]\simeq\F_p[x]/(f)\simeq K$, so $\F_p[\beta]$ has $q$ elements and equals $K'$.
(iii) $\F_p[x]/(f)$ is a field with $p^r$ elements (\Cref{thm:rings:adjoin-root}(i)); apply (ii).
(iv) If $\sigma\colon K\to K'$ is an isomorphism, so is $\sigma\circ\Frob_p^s$ for every $s$, and these are distinct for $0\le s<r$ by
\Cref{prop:rings:frobenius}(i).
\end{proof}

\begin{proposition}[Frobenius and subfields]\label{prop:rings:frobenius}
Let $q=p^r$.
\begin{enumerate}[label=(\roman*)]
\item $\Frob_p\colon x\mapsto x^p$ is a field automorphism of $\F_q$, its fixed points are exactly $\F_p$, and it has order exactly $r$.
\item For $d\mid r$, $\{x\in\F_q:x^{p^d}=x\}$ is a subfield with $p^d$ elements, and it is the only subfield with $p^d$ elements. A subfield with
$p^e$ elements exists only if $e\mid r$.
\item The fixed field of $\Frob_p^s$ is the subfield with $p^{\gcd(s,r)}$ elements.
\end{enumerate}
The fact that every automorphism of $\F_q$ is a power of $\Frob_p$, so that $\Gal(\F_{q^n}/\F_q)$ is cyclic of order $n$ generated by $x\mapsto x^q$,
is proved in \Cref{ch:galois}.
\end{proposition}

\begin{proof}
(i) $\Frob_p$ is an injective endomorphism of the finite field $\F_q$, hence bijective (\Cref{prop:rings:freshman}(ii)). Its fixed points are
the roots of $x^p-x$, at most $p$ of them (\Cref{prop:rings:poly-division}(iv)), and every element of $\F_p$ is one (Fermat). By
\Cref{prop:rings:finite-field-basic}(iii), $\Frob_p^r=\id$. If $\Frob_p^s=\id$ with $1\le s<r$, all $p^r$ elements would be roots of $x^{p^s}-x$, of degree
$p^s<p^r$.
(ii) The set is the fixed set of the automorphism $\Frob_p^d$, hence a subfield. As $d\mid r$, $p^d-1$ divides $p^r-1$ (write $r=dt$ and
$y=p^d$; then $y^t-1=(y-1)(y^{t-1}+\dots+1)$), and for the same reason $x^{p^d-1}-1$ divides $x^{p^r-1}-1$; multiplying by $x$, $x^{p^d}-x$ divides $x^{p^r}-x$; since $x^{p^r}-x$ splits into distinct linear factors over $\F_q$, so does its divisor $x^{p^d}-x$, which therefore has exactly $p^d$
roots in $\F_q$. A subfield $L$ with $p^d$ elements satisfies $x^{p^d}=x$ on $L$ (\Cref{prop:rings:finite-field-basic}(iii)), so it lies in this set,
and equals it by counting. If $L$ has $p^e$ elements, $\F_q$ is an $L$-vector space of some dimension $t$, and $p^r=p^{et}$.
(iii) Let $g=\gcd(s,r)=as+br$ with $a,b\in\Z$. In the cyclic group generated by $\Frob_p$ (of order $r$), $\Frob_p^g=(\Frob_p^s)^a$, since
$\Frob_p^{br}=\id$. So an element fixed by $\Frob_p^s$ is fixed by $\Frob_p^g$; the converse holds as $g\mid s$. Apply (ii) to $g\mid r$.
\end{proof}

\begin{proposition}\label{prop:rings:irreducible-count}
Let $q$ be a prime power and $r\ge1$. In $\F_q[x]$,
\[
x^{q^r}-x=\prod_{d\mid r}\ \prod_{f\text{ monic irreducible},\ \deg f=d}f,
\]
each factor occurring once. Hence $q^r=\sum_{d\mid r}dN_d$, where $N_d$ is the number of monic irreducible polynomials of degree $d$ over $\F_q$, and
\[
N_r=\frac1r\sum_{d\mid r}\mu(d)\,q^{r/d}>0 .
\]
Over $\F_2$, $N_1,\dots,N_8=2,1,2,3,6,9,18,30$; over $\F_3$, $N_1,\dots,N_5=3,3,8,18,48$, and the three monic irreducible quadratics are $x^2+1$,
$x^2+x+2$, $x^2+2x+2$.
\end{proposition}

\begin{proof}
The polynomial $x^{q^r}-x$ is separable, so no irreducible factor occurs twice ($f^2\mid g$ implies $f\mid g'$). Let $f$ be monic irreducible of
degree $d$, $K=\F_q[x]/(f)$, a field with $q^d$ elements, and $\alpha=\bar x$; $f$ is the minimal polynomial of $\alpha$, so $f\mid x^{q^r}-x$ iff $\alpha^{q^r}=\alpha$.
The map $\phi\colon y\mapsto y^q$ is an automorphism of $K$ fixing $\F_q$ (\Cref{prop:rings:finite-field-basic}(iii) for $\F_q$), and it has order exactly $d$ by the
argument of \Cref{prop:rings:frobenius}(i). If $\alpha^{q^r}=\alpha$, then $\phi^r$ fixes $\F_q$ and $\alpha$, hence all of $K=\F_q[\alpha]$, so $d\mid r$. If $d\mid r$, then
$\alpha^{q^d}=\alpha$ and so $\alpha^{q^r}=\alpha$ (apply $y\mapsto y^{q^d}$ repeatedly, $r/d$ times). Comparing degrees gives $q^r=\sum_{d\mid r}dN_d$. For Moebius
inversion recall $\sum_{e\mid m}\mu(e)=0$ for $m>1$: if $p_1,\dots,p_s$ ($s\ge1$) are the distinct primes dividing $m$, only squarefree $e$ contribute, and
the sum is $\sum_{T\subseteq\{1,\dots,s\}}(-1)^{\#T}=(1-1)^s=0$. Hence
\[
\sum_{d\mid r}\mu(d)q^{r/d}=\sum_{d\mid r}\mu(d)\sum_{e\mid r/d}eN_e=\sum_{e\mid r}eN_e\sum_{d\mid r/e}\mu(d)=rN_r,
\]
since the inner sum is $1$ for $e=r$ and $0$ otherwise. For $r\ge2$ the term $d=1$ is $q^r$, and the others have absolute value $q^{r/d}$ with
$r/d\le r/2$ distinct integers; so
\[
rN_r\ge q^r-\sum_{j=0}^{\lfloor r/2\rfloor}q^j>q^r-q^{\lfloor r/2\rfloor+1}\ge0,
\]
using $\sum_{j=0}^{M}q^j=(q^{M+1}-1)/(q-1)<q^{M+1}$ for $q\ge2$, and $\lfloor r/2\rfloor+1\le r$; and $N_1=q$. The numerical values follow from the formula; for the
quadratics over $\F_3$, each of the three listed has no root in $\F_3$ and $N_2=(9-3)/2=3$.
\end{proof}

For $q=2$, $r=2$ the factorisation reads $x^4-x=x(x+1)(x^2+x+1)$ in $\F_2[x]$: the two linear polynomials and the unique irreducible quadratic, and
$4=1\cdot2+2\cdot1$. The formula for $N_r$ is the finite-field analogue of the prime number theorem: $N_r=\frac{q^r}r+O\bigl(\frac{q^{r/2}}r\bigr)$, so a random
monic polynomial of degree $r$ is irreducible with probability about $1/r$.

\begin{proposition}[roots of unity in $\F_q$]\label{prop:rings:mu-in-fq}
Let $q=p^r$ and $n\ge1$.
\begin{enumerate}[label=(\roman*)]
\item $\mu_n(\F_q)$ is cyclic of order $\gcd(n,q-1)$; it has $n$ elements iff $n\mid q-1$.
\item If $p\nmid n$, the least $f\ge1$ with $\#\mu_n(\F_{q^f})=n$ is the order of $q$ in $(\Z/n\Z)^\times$.
\item In characteristic $p$, $x^{p^km}-1=(x^m-1)^{p^k}$; there are no roots of unity of order divisible by $p$.
\item Let $q$ be odd. Then $-1$ is a square in $\F_q$ iff $q\equiv1\pmod4$. If moreover $p\ne3$, then $-3$ is a square in $\F_q$ iff $\F_q$ contains a
primitive cube root of unity iff $q\equiv1\pmod3$; indeed $(1+2\zeta)^2=-3$ for a primitive cube root of unity $\zeta$, and $\zeta=(-1+s)/2$ is one if
$s^2=-3$.
\item Let $\sigma$ be an element of order $3$ of any group and $m$ an integer prime to $3$. Then $\sigma^m=\sigma$ if $m\equiv1\pmod3$ and
$\sigma^m=\sigma^{-1}$ if $m\equiv2\pmod3$. Likewise if $\sigma$ has order $4$ (or $2$) and $m$ is odd, $\sigma^m=\sigma$ or $\sigma^{-1}$ according as $m\equiv1$ or
$3\pmod4$. In particular this applies to $\mu_3(K)$ and $\mu_4(K)$ for any field $K$ and to $\zeta\mapsto\zeta^q$; in the notation of
\Cref{prop:rings:paper-characters}, $\sigma^m=\sigma^{\chi_{-3}(m)}$, resp. $\sigma^m=\sigma^{\chi_{-4}(m)}$.
\end{enumerate}
\end{proposition}

\begin{proof}
(i) $\mu_n(\F_q)$ is the kernel of $x\mapsto x^n$ on the cyclic group $\F_q^\times$ of order $q-1$, of order $\gcd(n,q-1)$ (\Cref{ch:groups}).
(ii) By (i), $\#\mu_n(\F_{q^f})=n$ iff $n\mid q^f-1$ iff $q^f\equiv1\pmod n$.
(iii) $(x^m-1)^{p^k}=x^{p^km}-1$ by \Cref{prop:rings:freshman}. If $\zeta$ has order $p^km$ with $k\ge1$, then $\eta=\zeta^m$ satisfies $\eta^{p^k}=1$, so $(\eta-1)^{p^k}=0$
and $\eta=1$, contradicting the order.
(iv) An element $y$ with $y^2=-1$ has order $4$ (as $-1\ne1$), and conversely if $y$ has order $4$ then $y^2$ has order $2$, so $y^2=-1$, the
unique element of order $2$ of the cyclic group $\F_q^\times$. Such $y$ exists iff $4\mid q-1$. Next, a primitive cube root of unity $\zeta$ satisfies
$\zeta^2+\zeta+1=0$ (as $\zeta\ne1$ and $x^3-1=(x-1)(x^2+x+1)$), so $(1+2\zeta)^2=1+4(\zeta+\zeta^2)=-3$. Conversely if $s^2=-3$ and $\zeta=(-1+s)/2$ ($2$ is
invertible), then $4(\zeta^2+\zeta+1)=(1-2s+s^2)+(-2+2s)+4=s^2+3=0$, and $\zeta\ne1$ because $1+1+1=3\ne0$ ($p\ne3$); so $\zeta$ has order $3$. By (i),
such $\zeta$ exists iff $3\mid q-1$.
(v) $\sigma^m$ depends only on $m$ modulo the order of $\sigma$, and $\sigma^2=\sigma^{-1}$, resp. $\sigma^3=\sigma^{-1}$. The value $\chi_{-3}(m)\equiv m\pmod3$, resp.
$\chi_{-4}(m)\equiv m\pmod4$, is \Cref{prop:rings:paper-characters}(i).
\end{proof}

\begin{paperlink}\label{pl:rings:euler}
In Lemma lem:IV the paper writes "$\rho(\Frob_q)\sigma\rho(\Frob_q)^{-1}=\sigma^q=\sigma^{\chi_{-3}(q)}$", for an element $\sigma$ of order $3$ generating the
image of inertia, and in the proof of Lemma lem:euler: "Frobenius centralises $\sigma$ if $q\equiv1\pmod3$ and inverts it if $q\equiv2$, for either
choice of Frobenius, so it acts on the line by $\chi_{-3}(q)$"; for types $\mathrm{III}$, $\mathrm{III}^*$ the same with an element $\tau$ of order $4$ and
$\chi_{-4}(q)$. The Galois input (Frobenius acts on the tame quotient of inertia by $x\mapsto x^q$) is \Cref{ch:local,ch:ant}; the arithmetic
input is \Cref{prop:rings:mu-in-fq}(v): $\sigma^q=\sigma^{\pm1}$ according to $q$ modulo $3$, and $\pm1=\chi_{-3}(q)$. Conjugation by Frobenius on the
invariant line spanned by $\sigma-\sigma^{-1}$ (resp. by $\tau$) is then multiplication by $\chi_{-3}(q)$ (resp. $\chi_{-4}(q)$), and it does not matter
whether the arithmetic or the geometric Frobenius is used, since $\pm1$ is its own inverse.
\end{paperlink}

\section{Powers in \texorpdfstring{$\F_q^\times$}{Fq*}}\label{sec:rings:powers}

This section is the arithmetic of Lemma lem:tamfactor (the Tamagawa factor at a multiplicative prime). There $q$ is a
prime, $u\in\Z_q^\times$ is the unit part of the Tate parameter, and the Tamagawa factor of $\ad T$ exceeds $t_\ell(E,q)=v_\ell(n)$ by
$\min(v_\ell(n),v_\ell(q-1),v_\ell(c))$, where $c$ is the class of the reduction $\bar u$ in $\F_q^\times/(\F_q^\times)^{\ell^{v_\ell(q-1)}}$. We prove that this group is
cyclic of order $\ell^{v_\ell(q-1)}$, that $v_\ell(c)$ does not depend on the identification with $\Z/\ell^{v_\ell(q-1)}$, and that $v_\ell(c)\ge k$ iff $\bar u$
is an $\ell^k$-th power. The statements hold for any prime power $q$; the group theory is that of cyclic groups (\Cref{ch:groups}).

\begin{theorem}\label{thm:rings:power-index}
Let $q$ be a prime power, $m\ge1$ and $d=\gcd(m,q-1)$.
\begin{enumerate}[label=(\roman*)]
\item The $m$-th power map of $\F_q^\times$ has kernel $\mu_d(\F_q)$, of order $d$, and image $(\F_q^\times)^m=(\F_q^\times)^d$, the unique subgroup of index $d$.
\item (Generalised Euler criterion.) $u\in\F_q^\times$ is an $m$-th power iff $u^{(q-1)/d}=1$.
\item For a prime $\ell$ and $k\ge0$,
\[
\bigl[\F_q^\times:(\F_q^\times)^{\ell^k}\bigr]=\ell^{\min(k,v_\ell(q-1))}.
\]
In particular, if $q\not\equiv1\pmod\ell$, the $\ell$-th power map is bijective and every element is an $\ell^k$-th power for every $k$; if $q\equiv1\pmod\ell$,
exactly a fraction $1/\ell$ of the elements are $\ell$-th powers.
\item ($m=2$, $q$ odd.) The squares form the subgroup of index $2$, and $u^{(q-1)/2}=\pm1$, with $+1$ iff $u$ is a square.
\end{enumerate}
\end{theorem}

\begin{proof}
$\F_q^\times$ is cyclic of order $q-1$ (\Cref{prop:rings:finite-field-basic}). (i) and (iii) are the statements on cyclic groups of \Cref{ch:groups}:
for $C$ cyclic of order $N$, $x\mapsto x^m$ has kernel $C[\gcd(m,N)]$ and image $C^{\gcd(m,N)}$ of index $\gcd(m,N)$; and $\gcd(\ell^k,q-1)=\ell^{\min(k,v_\ell(q-1))}$.
(ii) The image $C^d$ is the unique subgroup of order $(q-1)/d$, which is $C[(q-1)/d]=\{u:u^{(q-1)/d}=1\}$. (iv) is (i), (ii) with $m=2$, $d=2$: $u^{(q-1)/2}$
squares to $1$, so it is $\pm1$ (the roots of $x^2-1$).
\end{proof}

The hypothesis is $d=\gcd(m,q-1)$, not $m$: every element of $\F_{11}^\times$ is a cube, since $\gcd(3,10)=1$, although "$u^{10/3}$" is meaningless.

\begin{proposition}[the class of a unit modulo $\ell$-power powers]\label{prop:rings:class-c}
Let $q$ be a prime power, $\ell\nmid q$ a prime, $v=v_\ell(q-1)$, and $A=\F_q^\times/(\F_q^\times)^{\ell^v}$.
\begin{enumerate}[label=(\roman*)]
\item $A$ is cyclic of order $\ell^v$. The projection restricted to $\mu_{\ell^v}(\F_q)$, the $\ell$-primary part of $\F_q^\times$, is an isomorphism onto $A$,
and $u\mapsto u^{(q-1)/\ell^v}$ induces an isomorphism $A\simeq\mu_{\ell^v}(\F_q)$.
\item A generator $g$ of $\F_q^\times$ identifies $A$ with $\Z/\ell^v$: $u=g^a\mapsto c(u)=a\bmod\ell^v$. Replacing $g$ by another generator multiplies
$c(u)$ by a unit of $\Z/\ell^v$; so $v_\ell(c(u))$, read as in \Cref{conv:rings:valuations}, does not depend on $g$.
\item For $0\le k\le v$: $u\in(\F_q^\times)^{\ell^k}$ iff $v_\ell(c(u))\ge k$. For $k\ge v$: $(\F_q^\times)^{\ell^k}=(\F_q^\times)^{\ell^v}$.
\item For $n\in\Z\setminus\{0\}$,
\[
\min\bigl(v_\ell(n),v,v_\ell(c(u))\bigr)=\max\bigl\{k\le\min(v_\ell(n),v):u\in(\F_q^\times)^{\ell^k}\bigr\},
\]
and this is $\ge1$ iff $\ell\mid n$, $q\equiv1\pmod\ell$ and $u$ is an $\ell$-th power in $\F_q^\times$.
\item If $q$ is prime and $u\in\Z_{(q)}^\times$ (or $u\in\Z_q^\times$, \Cref{ch:local}), the class of $u$ means the class of its reduction $\bar u\in\F_q^\times$
(\Cref{prop:rings:zp-local}(iii)); we write $c(u)=c(\bar u)$.
\end{enumerate}
\end{proposition}

\begin{proof}
Write $C=\F_q^\times$, cyclic of order $q-1=\ell^vm$ with $\ell\nmid m$. (i) $A=C/C^{\ell^v}$ has order $\gcd(\ell^v,q-1)=\ell^v$ and is cyclic, as a quotient of a cyclic
group. $C^{\ell^v}$ is the unique subgroup of order $m$, i.e. $C[m]$, and $C=C[\ell^v]\times C[m]$ (\Cref{ch:groups}), so $C[\ell^v]=\mu_{\ell^v}(\F_q)$ maps
isomorphically onto $C/C[m]=A$. The map $u\mapsto u^m$ has kernel $C[m]=C^{\ell^v}$ and image $C^m=C[\ell^v]$.
(ii) Every $u$ is $g^a$ with $a$ well defined modulo $q-1$, hence modulo $\ell^v$; the map $u\mapsto a\bmod\ell^v$ is a surjective homomorphism
$C\to\Z/\ell^v$ whose kernel is $\{g^a:\ell^v\mid a\}=\langle g^{\ell^v}\rangle=C^{\ell^v}$. Another generator is $g'=g^b$ with $\gcd(b,q-1)=1$; then $u=g'^{a'}=g^{ba'}$, so
$c_g(u)=b\,c_{g'}(u)$ with $b$ a unit modulo $\ell^v$. Finally $\min(v_\ell(\tilde b\tilde c),v)=\min(v_\ell(\tilde c),v)$ when $\ell\nmid\tilde b$.
(iii) For $k\le v$, $\gcd(\ell^k,q-1)=\ell^k$, so $C^{\ell^k}=\langle g^{\ell^k}\rangle=\{g^a:\ell^k\mid a\}$ (this condition on $a$ is well defined modulo $q-1$), and $\ell^k\mid a$ iff
$c(u)\in\ell^k\Z/\ell^v\Z$ iff $v_\ell(c(u))\ge k$. For $k\ge v$, $C^{\ell^k}=C^{\gcd(\ell^k,q-1)}=C^{\ell^v}$.
(iv) By (iii) the set of $k\in[0,v]$ with $u\in C^{\ell^k}$ is $\{0,1,\dots,v_\ell(c(u))\}$, as $v_\ell(c(u))\le v$. Intersecting with $k\le v_\ell(n)$, the maximum is
$\min(v_\ell(n),v_\ell(c(u)))=\min(v_\ell(n),v,v_\ell(c(u)))$. It is $\ge1$ iff $v_\ell(n)\ge1$, $v\ge1$ and $v_\ell(c(u))\ge1$, the last meaning $u\in C^\ell$ by (iii).
(v) is a definition, justified by \Cref{prop:rings:zp-local}(iii).
\end{proof}

When $v=0$ the group $A$ is trivial, $c(u)=0\in\Z/1$ and $v_\ell(c(u))=0$; the minimum in (iv) is then $0$, as it should be. The
identification $A\simeq\F_q^\times\otimes\Z_\ell\simeq\Z_\ell(1)/(q-1)$ used in the proof of Lemma lem:tamfactor is \Cref{ch:modules,ch:galcoh}; the Kummer
interpretation $\bar u^{(q-1)/\ell^k}=$ the value of a Kummer character at Frobenius is \Cref{ch:galois}.

\begin{example}\label{ex:rings:powers}
(a) $q=7$, $\ell=3$: $v=1$; the cubes in $\F_7^\times$ are $\{1,6\}$. So $\bar u=-1$ is a cube and $\bar u=2$ is not. A multiplicative prime $q=7$ with $3\mid n_7$
and $u\equiv-1\pmod7$ therefore contributes $v_3(n_7)+1$ to $v_3(\Tam_7(\ad T))$ by Lemma lem:tamfactor.
(b) $q=19$, $\ell=3$: $v=2$, and $2$ is a primitive root modulo $19$. The cubes are $\{1,7,8,11,12,18\}$ and the ninth powers $\{1,18\}$. With $g=2$:
$c(2)=1$, $c(8)=3$, $c(-1)=c(2^9)=9\bmod9=0$, so $v_3(c)=0,1,2$ respectively: $2$ is not a cube, $8$ is a cube but not a ninth power, and $-1$ is a
ninth power.
(c) $q=13$, $\ell=3$: $v=1$, $(q-1)/3=4$, and $u$ is a cube iff $u^4=1$; the cubes are $\{1,5,8,12\}$, the fourth roots of unity.
(d) $q=5$, $\ell=3$: $v=0$, every element is a cube. $q=11$, $\ell=5$: $v=1$, the fifth powers are $\{\pm1\}$.
(All lists recomputed with PARI's \texttt{znprimroot} and \texttt{znlog}.)
\end{example}

\begin{paperlink}\label{pl:rings:tamfactor}
Lemma lem:tamfactor states: "$c$ the class of the reduction of $u$ in $\F_q^\times/(\F_q^\times)^{\ell^v}$, $v=v_\ell(q-1)$, identified with an element of
$\Z/\ell^v$ by the choice of a generator; $v_\ell(c)$ does not depend on that choice, $v_\ell(0)$ is read as $v$, and $v_\ell(c)\ge k$ if and only if
the reduction of $u$ is an $\ell^k$-th power in $\F_q^\times$, for $k\le v$." This is \Cref{prop:rings:class-c}(i) to (iii) and (v), with \Cref{conv:rings:valuations}.
The "In particular" clause ("with equality unless $q\equiv1\pmod\ell$, $\ell\mid n$ and the reduction of $u$ is an $\ell$-th power") is (iv). Remark
rem:ko(1) and Section sec:tam test exactly this condition on the tables; \Cref{ex:rings:powers} shows the three cases at $q=19$.
Here $q$ is prime, so $\F_q=\Z/q\Z$ and "the class of $u$" is computed as a discrete logarithm of $u\bmod q$ to a primitive root.
\end{paperlink}

\section{Roots of unity and cyclotomic polynomials}\label{sec:rings:cyclotomic}

Roots of unity appear in the paper as $\zeta_3$ and $i$ (the images of inertia of types $\mathrm{II}$ to $\mathrm{IV}^*$ have eigenvalues in $\mu_3$, $\mu_4$,
$\mu_6$), as $\mu_{\ell^k}$ (the cyclotomic character, the Tate module $\Z_\ell(1)$), and through the condition $q\equiv1\pmod\ell$ of Lemma
lem:tamfactor, which says that $\F_q$ contains the $\ell$-th roots of unity. This section proves what is needed about cyclotomic
polynomials over $\Z$ and modulo $p$; the irreducibility of all $\Phi_n$ and $\Gal(\Q(\zeta_n)/\Q)\simeq(\Z/n)^\times$ are \Cref{ch:galois}.

\begin{definition}\label{def:rings:roots-of-unity}
For a field $K$, $\mu_n(K)=\{x\in K:x^n=1\}$ and $\mu(K)=\bigcup_{n\ge1}\mu_n(K)$, the group of roots of unity of $K$. A \emph{primitive $n$-th root of
unity} is an element of order exactly $n$ in $K^\times$.
\end{definition}

\begin{proposition}\label{prop:rings:roots-of-unity}
\begin{enumerate}[label=(\roman*)]
\item $\mu_n(K)$ is cyclic of order dividing $n$; it has order $n$ iff $K$ contains a primitive $n$-th root of unity, and this requires
$\operatorname{char}K\nmid n$.
\item $\mu_n(\C)=\langle\zeta_n\rangle$; its primitive elements are the $\zeta_n^a$ with $\gcd(a,n)=1$, $\varphi(n)$ in number.
\item $\mu(\Q)=\mu(\R)=\{\pm1\}$, $\mu(\Q(i))=\mu_4$ and $\mu(\Q(\zeta_3))=\mu_6$.
\end{enumerate}
\end{proposition}

\begin{proof}
(i) $\mu_n(K)$ is a finite subgroup of $K^\times$ ($x^n-1$ has at most $n$ roots), hence cyclic (\Cref{ch:groups}), and its elements have order
dividing $n$; a cyclic group of order dividing $n$ has order $n$ iff it has an element of order $n$. If $\operatorname{char}K=p\mid n$, then
$x^n-1=(x^{n/p}-1)^p$ and $\mu_n(K)=\mu_{n/p}(K)$ has order at most $n/p$.
(ii) The roots of $x^n-1$ in $\C$ are the $n$ numbers $e^{2\pi ia/n}=\zeta_n^a$, $0\le a<n$; the order of $\zeta_n^a$ is $n/\gcd(a,n)$.
(iii) A real root of unity has absolute value $1$, so it is $\pm1$. Let $K$ be $\Q(i)$ or $\Q(\zeta_3)=\Q(\sqrt{-3})$, a two-dimensional $\Q$-vector space,
and $z\in\mu(K)$ of order $m$. The elements $1,z,z^2$ are linearly dependent over $\Q$, so the minimal polynomial $h$ of $z$ over $\Q$ (the
monic generator of the kernel of $\Q[x]\to\C$, $x\mapsto z$) has degree at most $2$, and it divides $x^m-1$. By \Cref{thm:rings:gauss-lemma}(iv), $h\in\Z[x]$.
If $\deg h=2$, $h=x^2+bx+c$ with roots $z$ and $\bar z$ (both roots of $x^m-1$, of absolute value $1$): $\abs c=\abs{z\bar z}=1$ and $\abs b=\abs{z+\bar z}\le2$.
If $c=-1$ the roots are real, hence $\pm1$, and $h=x^2-1$ is not irreducible; if $c=1$, $b=\pm2$ gives $(x\pm1)^2$, not irreducible. So $h$ is
$x\mp1$, $x^2+1$ or $x^2\pm x+1$, and $m\in\{1,2,4,3,6\}$. Now $\Q(i)$ does not contain $\zeta_3$: if $(a+bi)^2=-3$ with $a,b\in\Q$ then $ab=0$, $a^2-b^2=-3$, so
$a=0$ and $b^2=3$, impossible. And $\Q(\sqrt{-3})$ does not contain $i$: $(a+b\sqrt{-3})^2=-1$ forces $a=0$, $3b^2=1$. Hence $\mu(\Q(i))=\mu_4$ and
$\mu(\Q(\zeta_3))=\{\pm\zeta_3^j\}=\mu_6$.
\end{proof}

\begin{theorem}[cyclotomic polynomials]\label{thm:rings:cyclotomic}
For $n\ge1$ let $\Phi_n(x)=\prod_{1\le a\le n,\ \gcd(a,n)=1}(x-\zeta_n^a)\in\C[x]$.
\begin{enumerate}[label=(\roman*)]
\item $x^n-1=\prod_{d\mid n}\Phi_d(x)$.
\item $\Phi_n\in\Z[x]$, monic of degree $\varphi(n)$.
\item $\Phi_n(x)=\prod_{d\mid n}(x^d-1)^{\mu(n/d)}$, with $\mu$ the Moebius function.
\item $\Phi_p=x^{p-1}+\dots+x+1$; $\Phi_{p^k}(x)=\Phi_p(x^{p^{k-1}})$; $\Phi_{2m}(x)=\Phi_m(-x)$ for odd $m>1$.
\item $\Phi_1(1)=0$; $\Phi_n(1)=p$ if $n=p^k$, $k\ge1$; $\Phi_n(1)=1$ if $n$ has at least two distinct prime factors; $\Phi_n(0)=1$ for $n\ge2$.
\item $\Phi_{p^k}$ is irreducible over $\Q$.
\end{enumerate}
\end{theorem}

\begin{proof}
(i) Group the $n$-th roots of unity by their order $d\mid n$; those of order $d$ are the primitive $d$-th roots, the roots of $\Phi_d$
(\Cref{prop:rings:roots-of-unity}(ii)).
(ii) Induction on $n$: $\Phi_1=x-1$. By (i), $\Phi_n=(x^n-1)/g$ with $g=\prod_{d\mid n,d<n}\Phi_d$ monic in $\Z[x]$ by induction. Division of $x^n-1$ by the
monic $g$ in $\Z[x]$ (\Cref{prop:rings:poly-division}(i)) gives a quotient in $\Z[x]$ and a remainder, which must vanish by uniqueness of
division in $\C[x]$. The degree is the number of primitive $n$-th roots of unity, $\varphi(n)$.
(iii) In the field $\C(x)$, using (i),
$\prod_{d\mid n}(x^d-1)^{\mu(n/d)}=\prod_{d\mid n}\prod_{e\mid d}\Phi_e^{\mu(n/d)}=\prod_{e\mid n}\Phi_e^{\sum_{e\mid d\mid n}\mu(n/d)}$, and $\sum_{e\mid d\mid n}\mu(n/d)=\sum_{k\mid n/e}\mu(k)$ (put $k=n/d$, which
runs over the divisors of $n/e$) is $1$ if $e=n$ and $0$ otherwise (the identity $\sum_{k\mid m}\mu(k)=0$ for $m>1$ is proved in \Cref{prop:rings:irreducible-count}).
For example $\Phi_6=\frac{(x^6-1)(x-1)}{(x^3-1)(x^2-1)}=\frac{x^3+1}{x+1}=x^2-x+1$.
(iv) By (iii), $\Phi_{p^k}=(x^{p^k}-1)/(x^{p^{k-1}}-1)=\Phi_p(x^{p^{k-1}})$, and $\Phi_p=(x^p-1)/(x-1)$. For odd $m>1$ the divisors of $2m$ are the $d$ and
$2d$ with $d\mid m$, and $\mu(2m/d)=-\mu(m/d)$, $\mu(2m/2d)=\mu(m/d)$; so $\Phi_{2m}=\prod_{d\mid m}\bigl((x^{2d}-1)/(x^d-1)\bigr)^{\mu(m/d)}=\prod_{d\mid m}(x^d+1)^{\mu(m/d)}$.
On the other hand $\Phi_m(-x)=\prod_{d\mid m}(-(x^d+1))^{\mu(m/d)}$ ($d$ is odd) $=(-1)^{\sum_{d\mid m}\mu(m/d)}\prod_{d\mid m}(x^d+1)^{\mu(m/d)}$, and the sign is $+1$
as $m>1$.
(v) $\Phi_1(1)=0$ and $\Phi_{p^k}(1)=\Phi_p(1)=p$ by (iv). For $n>1$, (i) gives $1+x+\dots+x^{n-1}=\prod_{d\mid n,\,d>1}\Phi_d(x)$, and at $x=1$:
$n=\prod_{d\mid n,d>1}\Phi_d(1)$. The prime power divisors $d=p^j>1$ contribute $\prod_pp^{v_p(n)}=n$, so the product of $\Phi_d(1)$ over the
divisors $d>1$ that are not prime powers is $1$; by induction on $n$ all of them except $d=n$ are $1$, hence so is $\Phi_n(1)$ when $n$ is
not a prime power. Similarly $x^n-1$ at $x=0$ gives $-1=\Phi_1(0)\prod_{d\mid n,d>1}\Phi_d(0)=-\prod_{d\mid n,d>1}\Phi_d(0)$, and induction gives $\Phi_n(0)=1$.
(vi) Put $g(x)=\Phi_{p^k}(x+1)$, monic of degree $\varphi(p^k)$. Modulo $p$, $\Phi_{p^k}=(x^{p^k}-1)/(x^{p^{k-1}}-1)\equiv(x-1)^{p^k}/(x-1)^{p^{k-1}}=(x-1)^{\varphi(p^k)}$
(\Cref{prop:rings:freshman}; division by a monic polynomial commutes with reduction), so $g\equiv x^{\varphi(p^k)}\pmod p$, and $g(0)=\Phi_{p^k}(1)=p$. By
Eisenstein (\Cref{thm:rings:eisenstein}) $g$ is irreducible, and so is $\Phi_{p^k}$, since $x\mapsto x+1$ is an automorphism of $\Q[x]$ preserving degrees.
\end{proof}

\begin{center}
\small
\begin{tabular}{lll@{\qquad}lll}
\toprule
$n$ & $\Phi_n(x)$ & $\Phi_n(1)$ & $n$ & $\Phi_n(x)$ & $\Phi_n(1)$\\
\midrule
$1$ & $x-1$ & $0$ & $7$ & $x^6+x^5+x^4+x^3+x^2+x+1$ & $7$\\
$2$ & $x+1$ & $2$ & $8$ & $x^4+1$ & $2$\\
$3$ & $x^2+x+1$ & $3$ & $9$ & $x^6+x^3+1$ & $3$\\
$4$ & $x^2+1$ & $2$ & $10$ & $x^4-x^3+x^2-x+1$ & $1$\\
$5$ & $x^4+x^3+x^2+x+1$ & $5$ & $11$ & $x^{10}+x^9+\dots+x+1$ & $11$\\
$6$ & $x^2-x+1$ & $1$ & $12$ & $x^4-x^2+1$ & $1$\\
\bottomrule
\end{tabular}
\end{center}

\begin{proposition}\label{prop:rings:cyclotomic-mod-p}
Let $p$ be a prime with $p\nmid n$, and write $\bar\Phi_n$ for the reduction of $\Phi_n$ modulo $p$.
\begin{enumerate}[label=(\roman*)]
\item $\bar\Phi_n$ divides $x^n-1$ in $\F_p[x]$, so it is separable; in any field of characteristic $p$ its roots are exactly the elements of
order $n$.
\item For $a\in\Z$: $p\mid\Phi_n(a)$ iff $a\bmod p$ has order exactly $n$ in $\F_p^\times$. Hence every prime divisor $p\nmid n$ of $\Phi_n(a)$ satisfies
$p\equiv1\pmod n$.
\item There are infinitely many primes $p\equiv1\pmod n$.
\item $\bar\Phi_n$ is a product of distinct linear factors in $\F_p[x]$ iff $p\equiv1\pmod n$.
\end{enumerate}
The complete factorisation pattern ($\varphi(n)/f$ irreducible factors of degree $f$, $f$ the order of $p$ modulo $n$) is \Cref{ch:galois}.
\end{proposition}

\begin{proof}
(i) Reduce \Cref{thm:rings:cyclotomic}(i) modulo $p$: $x^n-1=\prod_{d\mid n}\bar\Phi_d$, and $x^n-1$ is separable over $\F_p$
(\Cref{prop:rings:derivative}(iii)). Let $K$ have characteristic $p$; enlarging $K$ (\Cref{thm:rings:adjoin-root}(iii)) we may assume $x^n-1$
splits in $K$, which does not change orders of elements of $K$. The roots of $x^n-1$ are the $n$ elements of $\mu_n(K)$, those of order $d$
being, by induction on $n$, the roots of $\bar\Phi_d$ for $d<n$; since the roots of the $\bar\Phi_d$ are pairwise distinct, the roots of $\bar\Phi_n$ are
the remaining elements, of order exactly $n$.
(ii) $p\mid\Phi_n(a)$ iff $\bar a$ is a root of $\bar\Phi_n$ in $\F_p$, iff $\bar a$ has order $n$ in $\F_p^\times$ by (i) ($\bar a\ne0$ because $\bar\Phi_n(0)=1$ for $n\ge2$; for
$n=1$ the claim is clear). Then $n\mid p-1$.
(iii) For $n=1$ this is Euclid's theorem; let $n\ge2$ and $S$ a finite set of primes, $N=n\prod_{p\in S}p$. For an integer $t$, $\Phi_n(tN)\equiv\Phi_n(0)=1\pmod N$,
so $\Phi_n(tN)$ is prime to $n$ and to every $p\in S$. For $t$ large, $\abs{\Phi_n(tN)}>1$ ($\Phi_n$ is monic of positive degree), so it has a prime
divisor $p\notin S$, $p\nmid n$, and $p\equiv1\pmod n$ by (ii).
(iv) If $p\equiv1\pmod n$, the cyclic group $\F_p^\times$ contains $\varphi(n)$ elements of order $n$, all roots of $\bar\Phi_n$, which has degree $\varphi(n)$.
Conversely a root in $\F_p$ has order $n$, so $n\mid p-1$.
\end{proof}

For example $\Phi_{12}(2)=16-4+1=13\equiv1\pmod{12}$ and $\Phi_9(2)=\Phi_3(8)=73\equiv1\pmod9$; and $2$ has order $12$ modulo $13$ and order $9$ modulo $73$.
In the paper the primes $q\equiv1\pmod\ell$ of Lemma lem:tamfactor are exactly those with $\mu_\ell\subset\F_q$ (\Cref{prop:rings:mu-in-fq}(i)).

\section{Quadratic residues and quadratic reciprocity}\label{sec:rings:qr}

The paper's quadratic characters are Kronecker symbols: $\chi_{-3}=\bigl(\frac{-3}\cdot\bigr)$, $\chi_{-4}=\bigl(\frac{-4}\cdot\bigr)$, and $\chi_{q^*}=\bigl(\frac{q^*}\cdot\bigr)$, the
character of $\Q(\sqrt{q^*})$ by which Lemma lem:twist twists. That $\chi_{q^*}$ is "the quadratic character of conductor $q$", i.e. that
$\bigl(\frac{q^*}n\bigr)=\bigl(\frac nq\bigr)$ for all $n$, is quadratic reciprocity (\Cref{prop:rings:qstar}). We prove reciprocity by Gauss sums in
characteristic $p$, extend it to Jacobi and Kronecker symbols, and show that $\bigl(\frac D\cdot\bigr)$ is a Dirichlet character modulo $\abs D$ for every
discriminant $D$. \Cref{ch:ant} gives a second proof and the meaning of $\chi_D(p)$ ($D$ fundamental) as a splitting law.

\begin{definition}\label{def:rings:legendre}
For an odd prime $p$ and $a\in\Z$ the \emph{Legendre symbol} is $\bigl(\frac ap\bigr)=0$ if $p\mid a$, $1$ if $a$ is a nonzero square modulo $p$ (a
\emph{quadratic residue}), and $-1$ otherwise (a \emph{non-residue}).
\end{definition}

\begin{proposition}\label{prop:rings:legendre}
Let $p$ be an odd prime.
\begin{enumerate}[label=(\roman*)]
\item (Euler's criterion) $\bigl(\frac ap\bigr)\equiv a^{(p-1)/2}\pmod p$.
\item $\bigl(\frac{ab}p\bigr)=\bigl(\frac ap\bigr)\bigl(\frac bp\bigr)$; $a\mapsto\bigl(\frac ap\bigr)$ is the unique homomorphism of order $2$ from $(\Z/p)^\times$ to $\{\pm1\}$; there are $(p-1)/2$
residues and $(p-1)/2$ non-residues.
\item $\bigl(\frac{-1}p\bigr)=(-1)^{(p-1)/2}$.
\item For $p=3$: $\bigl(\frac a3\bigr)\equiv a\pmod3$ for all $a$.
\end{enumerate}
\end{proposition}

\begin{proof}
(i) For $p\mid a$ both sides are $0$; otherwise it is \Cref{thm:rings:power-index}(iv) with $q=p$. (ii) Multiplicativity follows from (i), as
$1\not\equiv-1\pmod p$. A homomorphism onto $\{\pm1\}$ has a kernel of index $2$, and the cyclic group $(\Z/p)^\times$ has only one subgroup of
index $2$, the squares (\Cref{thm:rings:power-index}(i)). (iii) is (i) for $a=-1$. (iv) The nonzero squares modulo $3$ are $\{1\}$, so $\bigl(\frac23\bigr)=-1\equiv2$.
\end{proof}

For $p=7$ the squares of $1,\dots,6$ are $1,4,2,2,4,1$, so the residues are $\{1,2,4\}$ and the non-residues $\{3,5,6\}$. Euler's criterion agrees:
$3^3=27\equiv-1$ and $2^3=8\equiv1\pmod7$. In particular $\bigl(\frac{-1}7\bigr)=-1$, as (iii) says for $7\equiv3\pmod4$.

\begin{proposition}[the quadratic Gauss sum]\label{prop:rings:quadratic-gauss-sum}
Let $q$ be an odd prime, $K$ a field of characteristic $\ne q$ containing a primitive $q$-th root of unity $\zeta$, and
$g=\sum_{a=1}^{q-1}\bigl(\frac aq\bigr)\zeta^a\in K$.
\begin{enumerate}[label=(\roman*)]
\item $g^2=q^*=(-1)^{(q-1)/2}q$ in $K$; in particular $g\ne0$.
\item For $b$ prime to $q$, $\sum_{a=1}^{q-1}\bigl(\frac aq\bigr)\zeta^{ab}=\bigl(\frac bq\bigr)g$.
\item $\sqrt{q^*}\in\Q(\zeta_q)$: over $\C$ with $\zeta=\zeta_q$, $g\in\Z[\zeta_q]$ and $g^2=q^*$.
\end{enumerate}
\end{proposition}

\begin{proof}
For $b\in\Z$, $\sum_{a=0}^{q-1}\zeta^{ab}$ is $q$ if $q\mid b$ and $0$ otherwise: if $q\nmid b$, $\zeta^b\ne1$ is a root of $x^q-1=(x-1)(1+x+\dots+x^{q-1})$, hence of the
second factor. So $\sum_{a=1}^{q-1}\zeta^{ab}$ is $q-1$ if $q\mid b$ and $-1$ otherwise. Also $\sum_{c=1}^{q-1}\bigl(\frac cq\bigr)=0$ by
\Cref{prop:rings:legendre}(ii).
(ii) As $a$ runs over $(\Z/q)^\times$ so does $c=ab$, and $\bigl(\frac aq\bigr)=\bigl(\frac cq\bigr)\bigl(\frac bq\bigr)$ since $\bigl(\frac bq\bigr)^{-1}=\bigl(\frac bq\bigr)$.
(i) $g^2=\sum_{a,b}\bigl(\frac{ab}q\bigr)\zeta^{a+b}$. Substitute $b=ac$ ($c$ runs over $(\Z/q)^\times$ for fixed $a$): $\bigl(\frac{a\cdot ac}q\bigr)=\bigl(\frac{a^2}q\bigr)\bigl(\frac cq\bigr)=\bigl(\frac cq\bigr)$. The
inner sum over $a$ is $\sum_{a=1}^{q-1}\zeta^{a(1+c)}$, which by the first paragraph is $q-1$ for the one value $c\equiv-1$ and $-1$ for the others; then
$\sum_{c\ne-1}\bigl(\frac cq\bigr)=0-\bigl(\frac{-1}q\bigr)$. So
\[
g^2=\sum_{c=1}^{q-1}\Bigl(\frac cq\Bigr)\sum_{a=1}^{q-1}\zeta^{a(1+c)}=\Bigl(\frac{-1}q\Bigr)(q-1)-\sum_{c\ne-1}\Bigl(\frac cq\Bigr)=\Bigl(\frac{-1}q\Bigr)(q-1)+\Bigl(\frac{-1}q\Bigr)=\Bigl(\frac{-1}q\Bigr)q,
\]
which is $q^*$ by \Cref{prop:rings:legendre}(iii). It is nonzero because $\operatorname{char}K\ne q$. (iii) is (i) over $\C$.
\end{proof}

For $q=3$ over $\C$: $g=\zeta_3-\zeta_3^2=1+2\zeta_3=\sqrt{-3}$, and $g^2=-3=3^*$. For $q=5$: $g=\zeta_5-\zeta_5^2-\zeta_5^3+\zeta_5^4=2\cos\frac{2\pi}5-2\cos\frac{4\pi}5=\sqrt5$
(\Cref{ex:rings:gauss-sums}). The point of allowing an arbitrary field $K$ is the proof of reciprocity below, where $K$ is a finite field of
characteristic $p\ne q$ and the identity $g^2=q^*$ is used inside $K$.

\begin{remark}[the sign of the Gauss sum]\label{rem:rings:gauss-sign}
Over $\C$ with $\zeta=\zeta_q$, Gauss proved $g=\sqrt q$ if $q\equiv1\pmod4$ and $g=i\sqrt q$ if $q\equiv3\pmod4$. We do not prove this and do not use it;
see \cite[Chapter 6, Section 6.4]{IrelandRosen}. (Checked numerically for all $q<60$.)
\end{remark}

\begin{theorem}[quadratic reciprocity]\label{thm:rings:qr}
Let $p\ne q$ be odd primes.
\begin{enumerate}[label=(\roman*)]
\item $\bigl(\frac pq\bigr)\bigl(\frac qp\bigr)=(-1)^{\frac{p-1}2\cdot\frac{q-1}2}$; equivalently $\bigl(\frac{q^*}p\bigr)=\bigl(\frac pq\bigr)$.
\item (First supplement) $\bigl(\frac{-1}p\bigr)=(-1)^{(p-1)/2}$.
\item (Second supplement) $\bigl(\frac2p\bigr)=(-1)^{(p^2-1)/8}$, that is, $2$ is a square modulo $p$ iff $p\equiv\pm1\pmod8$.
\end{enumerate}
\end{theorem}

\begin{proof}
(i) Let $f$ be the order of $p$ modulo $q$. By \Cref{prop:rings:mu-in-fq}(ii) (and \Cref{thm:rings:finite-field-exists}) the field $K=\F_{p^f}$ contains a
primitive $q$-th root of unity $\zeta$; its characteristic $p$ is different from $q$. Let $g=\sum_a\bigl(\frac aq\bigr)\zeta^a\in K$. By
\Cref{prop:rings:quadratic-gauss-sum}(i), $g^2=q^*$, an element of $\F_p\subseteq K$. By \Cref{prop:rings:freshman}, and because
$\bigl(\frac aq\bigr)^p=\bigl(\frac aq\bigr)$ ($p$ odd),
\[
g^p=\sum_a\Bigl(\frac aq\Bigr)^p\zeta^{ap}=\sum_a\Bigl(\frac aq\Bigr)\zeta^{ap}=\Bigl(\frac pq\Bigr)g
\]
by \Cref{prop:rings:quadratic-gauss-sum}(ii). As $g\ne0$, $g^{p-1}=\bigl(\frac pq\bigr)$. But $g^{p-1}=(g^2)^{(p-1)/2}=(q^*)^{(p-1)/2}=\bigl(\frac{q^*}p\bigr)$ in $\F_p$ by Euler's criterion. So
$\bigl(\frac{q^*}p\bigr)=\bigl(\frac pq\bigr)$ in $\F_p$, hence as integers ($1\ne-1$ in $\F_p$). Finally
$\bigl(\frac{q^*}p\bigr)=\bigl(\frac{-1}p\bigr)^{(q-1)/2}\bigl(\frac qp\bigr)=(-1)^{\frac{p-1}2\frac{q-1}2}\bigl(\frac qp\bigr)$.
(ii) is \Cref{prop:rings:legendre}(iii).
(iii) As $p$ is odd, $8\mid p^2-1$, so $\F_{p^2}$ contains an element $\zeta$ of order $8$ (\Cref{prop:rings:mu-in-fq}(i)). Then $\zeta^4=-1$ and $\zeta^{-2}=\zeta^6=-\zeta^2$.
Put $\tau=\zeta+\zeta^{-1}$; then $\tau^2=\zeta^2+2+\zeta^{-2}=2$, so $\tau\ne0$. Now $\tau^p=\zeta^p+\zeta^{-p}$, which is $\tau$ if $p\equiv\pm1\pmod8$, and if $p\equiv\pm3\pmod8$
it is $\zeta^3+\zeta^{-3}=\zeta^3+\zeta^5=\zeta^4(\zeta^{-1}+\zeta)=-\tau$. Hence $2^{(p-1)/2}=\tau^{p-1}=\pm1$ accordingly, and Euler's criterion concludes. Note that
$(-1)^{(p^2-1)/8}$ is $1$ for $p\equiv\pm1$ and $-1$ for $p\equiv\pm3\pmod8$.
\end{proof}

\begin{definition}\label{def:rings:jacobi}
For an odd integer $n=\prod_ip_i^{e_i}>0$ and $a\in\Z$, the \emph{Jacobi symbol} is $\bigl(\frac an\bigr)=\prod_i\bigl(\frac a{p_i}\bigr)^{e_i}$ (Legendre symbols; $\bigl(\frac a1\bigr)=1$).
\end{definition}

\begin{proposition}\label{prop:rings:jacobi}
Let $m,n>0$ be odd.
\begin{enumerate}[label=(\roman*)]
\item $\bigl(\frac an\bigr)$ is multiplicative in $a$ and in $n$, depends only on $a\bmod n$, and vanishes iff $\gcd(a,n)>1$.
\item (Reciprocity and supplements.) If $\gcd(m,n)=1$, then $\bigl(\frac mn\bigr)\bigl(\frac nm\bigr)=(-1)^{\frac{m-1}2\cdot\frac{n-1}2}$. Moreover
$\bigl(\frac{-1}n\bigr)=(-1)^{(n-1)/2}$ and $\bigl(\frac2n\bigr)=(-1)^{(n^2-1)/8}$.
\item $\bigl(\frac an\bigr)=1$ does not imply that $a$ is a square modulo $n$: $\bigl(\frac2{15}\bigr)=\bigl(\frac23\bigr)\bigl(\frac25\bigr)=(-1)(-1)=1$, but $2$ is not a square modulo $3$, hence
not modulo $15$. (If $\bigl(\frac an\bigr)=-1$, then $a$ is not a square modulo $n$.)
\item Reciprocity computes $\bigl(\frac an\bigr)$ without factoring $n$: reduce $a$ modulo $n$, remove the factors $-1$ and $2$ with (ii), and invert
the remaining odd symbol. For instance, $1009$ being prime,
$\bigl(\frac{713}{1009}\bigr)=\bigl(\frac{1009}{713}\bigr)=\bigl(\frac{296}{713}\bigr)=\bigl(\frac2{713}\bigr)^3\bigl(\frac{37}{713}\bigr)=\bigl(\frac{713}{37}\bigr)=\bigl(\frac{10}{37}\bigr)=\bigl(\frac2{37}\bigr)\bigl(\frac5{37}\bigr)=(-1)\bigl(\frac{37}5\bigr)=-\bigl(\frac25\bigr)=1$,
so $713$ is a square modulo $1009$ (PARI's \texttt{kronecker(713,1009)} returns $1$).
\end{enumerate}
\end{proposition}

\begin{proof}
(i) is inherited from the Legendre symbol factor by factor. (ii) Let $e(n)=\frac{n-1}2$ and $w(n)=\frac{n^2-1}8$ modulo $2$, for odd $n$. Both are additive:
$mn-1=(m-1)(n-1)+(m-1)+(n-1)$ with $4\mid(m-1)(n-1)$, and $m^2n^2-1=(m^2-1)(n^2-1)+(m^2-1)+(n^2-1)$ with $64\mid(m^2-1)(n^2-1)$. Write $m=\prod p_i$, $n=\prod q_j$
(with repetitions). By \Cref{thm:rings:qr},
$\bigl(\frac mn\bigr)\bigl(\frac nm\bigr)=\prod_{i,j}\bigl(\frac{p_i}{q_j}\bigr)\bigl(\frac{q_j}{p_i}\bigr)=(-1)^{\sum_{i,j}e(p_i)e(q_j)}=(-1)^{e(m)e(n)}$, and likewise
$\bigl(\frac{-1}n\bigr)=(-1)^{\sum_je(q_j)}=(-1)^{e(n)}$, $\bigl(\frac2n\bigr)=(-1)^{\sum_jw(q_j)}=(-1)^{w(n)}$. (iii) A square modulo $15$ is a square modulo $3$. If $\bigl(\frac an\bigr)=-1$,
some $\bigl(\frac a{p_i}\bigr)=-1$. (iv) In the example $713\equiv1\pmod8$, $37\equiv5\pmod8$, and $713\equiv1$, $37\equiv1\pmod4$, so no signs arise from (ii) in
the inversions.
\end{proof}

\begin{definition}\label{def:rings:kronecker}
The \emph{Kronecker symbol} $\bigl(\frac an\bigr)$, for $a,n\in\Z$, extends the Jacobi symbol using \Cref{conv:rings:symbols}: for $n=s2^e\prod_ip_i^{e_i}\ne0$
with a sign $s=\pm1$ and odd primes $p_i$,
\[
\Bigl(\frac an\Bigr)=\Bigl(\frac as\Bigr)\Bigl(\frac a2\Bigr)^e\prod_i\Bigl(\frac a{p_i}\Bigr)^{e_i},
\]
and $\bigl(\frac a0\bigr)=1$ if $a=\pm1$, $0$ otherwise. A \emph{discriminant} is an integer $D\ne0$ with $D\equiv0$ or $1\pmod4$. A \emph{fundamental
discriminant} is a discriminant $D$ with either $D\equiv1\pmod4$ squarefree, or $D=4m$ with $m\equiv2,3\pmod4$ squarefree. The \emph{prime
discriminants} are $-4$, $8$, $-8$ and $q^*=(-1)^{(q-1)/2}q$ for odd primes $q$. For a fundamental discriminant $D$ put $\chi_D(a):=\bigl(\frac Da\bigr)$ for all
$a\in\Z$; for a discriminant that is not fundamental no name $\chi_D$ is introduced, and the function $\bigl(\frac D\cdot\bigr)$ is written as a Kronecker
symbol (\Cref{conv:rings:characters}).
\end{definition}

The definition is made so that $n\mapsto\bigl(\frac Dn\bigr)$, a function of the lower entry, is a character; for the Jacobi symbol it is the upper
entry that varies. For $D=-3$ and $n=1,\dots,6$ the values are $1,-1,0,1,-1,0$: at $n=2$ because $-3\equiv5\pmod8$, at $n=4$ because
$\bigl(\frac{-3}2\bigr)^2=1$, and at $n=5$ because $-3\equiv2$ is not a square modulo $5$. These are the values of $\bigl(\frac n3\bigr)$, a first instance of
\Cref{prop:rings:qstar}. The congruence conditions on fundamental discriminants say that $D$ is the discriminant of the quadratic field
$\Q(\sqrt D)$ (\Cref{ch:ant}): $-3$, $-4$, $5$, $8$, $-8$, $12$ are fundamental, while $-12=-3\cdot2^2$ and $20=5\cdot2^2$ are not.

\begin{proposition}\label{prop:rings:kronecker}
Let $D$ be a discriminant.
\begin{enumerate}[label=(\roman*)]
\item $n\mapsto\bigl(\frac Dn\bigr)$ is completely multiplicative on $\Z\setminus\{0\}$ and vanishes exactly when $\gcd(D,n)>1$. For nonzero $a,b$, $\bigl(\frac{ab}n\bigr)=\bigl(\frac an\bigr)\bigl(\frac bn\bigr)$.
\item If $D\equiv1\pmod4$, then $\bigl(\frac Dn\bigr)=\bigl(\frac n{\abs D}\bigr)$ (Jacobi symbol) for every $n>0$.
\item For $n,n'>0$ with $n\equiv n'\pmod{\abs D}$, $\bigl(\frac Dn\bigr)=\bigl(\frac D{n'}\bigr)$; and $\bigl(\frac D{-n}\bigr)=\sgn(D)\bigl(\frac Dn\bigr)$. Hence $\bigl(\frac D\cdot\bigr)$ is a
Dirichlet character modulo $\abs D$ (\Cref{def:rings:dirichlet}), real (with values in $\{0,\pm1\}$), with value $\sgn D$ at $-1$. For fundamental $D$
this character is $\chi_D$.
\item Values at primes: $\bigl(\frac Dp\bigr)$ is the Legendre symbol for odd $p$; $\bigl(\frac D2\bigr)=0$ if $4\mid D$, $1$ if $D\equiv1\pmod8$, $-1$ if $D\equiv5\pmod8$.
\item For discriminants $D,D'$, $\bigl(\frac{DD'}n\bigr)=\bigl(\frac Dn\bigr)\bigl(\frac{D'}n\bigr)$ for all $n\in\Z$; and for $f\ne0$, $\bigl(\frac{Df^2}n\bigr)=\bigl(\frac Dn\bigr)$ for $\gcd(n,f)=1$.
So for fundamental $D$ the character $\bigl(\frac{Df^2}\cdot\bigr)$ modulo $\abs{D}f^2$ is $\chi_D$ made imprimitive at the primes dividing $f$
(\Cref{prop:rings:conductor}), and in particular $\chi_{DD'}=\chi_D\chi_{D'}$ when $D$, $D'$ and $DD'$ are fundamental.
\end{enumerate}
\end{proposition}

\begin{proof}
(i) Multiplicativity in $n$ holds by construction (for the sign, $\bigl(\frac a{-1}\bigr)^2=1$ when $a\ne0$). Multiplicativity in nonzero $a$ holds for each factor:
Legendre symbols, $\bigl(\frac\cdot2\bigr)$ (on odd $a$ it is the character of $(\Z/8)^\times$ with kernel $\{\pm1\}$, and it is $0$ on even $a$), and $\bigl(\frac\cdot{-1}\bigr)=\sgn$. A
factor vanishes iff $p_i\mid D$, or $e\ge1$ and $D$ is even; that is, iff a prime dividing $n$ divides $D$.
(ii) Let $n>0$, $n=2^em$ with $m$ odd. If $\gcd(D,n)>1$ both sides vanish; assume not. For the odd part, write $D=s\abs D$, $s=\sgn D$. By (i) and
\Cref{prop:rings:jacobi}(ii), $\bigl(\frac Dm\bigr)=\bigl(\frac sm\bigr)\bigl(\frac{\abs D}m\bigr)=\bigl(\frac sm\bigr)(-1)^{\frac{m-1}2\frac{\abs D-1}2}\bigl(\frac m{\abs D}\bigr)$. If $s=1$ then $\abs D\equiv1\pmod4$ and the sign is $1$. If $s=-1$
then $\abs D\equiv3\pmod4$ and $(-1)^{\frac{m-1}2}=\bigl(\frac{-1}m\bigr)$ cancels $\bigl(\frac sm\bigr)$. So $\bigl(\frac Dm\bigr)=\bigl(\frac m{\abs D}\bigr)$. For the factor $2$:
$\bigl(\frac D2\bigr)=1$ iff $D\equiv1\pmod8$ iff $\abs D\equiv\pm1\pmod8$ (as $D\equiv1\pmod4$), iff $\bigl(\frac2{\abs D}\bigr)=1$ (\Cref{prop:rings:jacobi}(ii)). Multiply.
(iii) For $D=1$ the symbol $\bigl(\frac1n\bigr)$ is $1$ for every $n$ (each factor is), and everything is trivial; let $D\ne1$, so that $\abs D\ge3$
and $n=\abs D-1$ is positive. If $D\equiv1\pmod4$, (ii) shows that $\bigl(\frac Dn\bigr)$ depends only on $n\bmod\abs D$ for $n>0$, and $\bigl(\frac D{\abs D-1}\bigr)=\bigl(\frac{-1}{\abs D}\bigr)=(-1)^{(\abs D-1)/2}=\sgn D$. If $D=2^tD'$ with $D'$
odd and $t\ge2$: for even $n$ both $\bigl(\frac Dn\bigr)$ and $\bigl(\frac D{n'}\bigr)$ vanish ($n'\equiv n\pmod4$ is even). For odd $n>0$, as in (ii),
$\bigl(\frac Dn\bigr)=\bigl(\frac2n\bigr)^t\bigl(\frac{\sgn D'}n\bigr)(-1)^{\frac{n-1}2\frac{\abs{D'}-1}2}\bigl(\frac n{\abs{D'}}\bigr)$. Here $\bigl(\frac2n\bigr)$ depends on $n\bmod8$, the next two factors on $n\bmod4$, and the last on
$n\bmod\abs{D'}$; if $t\ge3$ the period $8\abs{D'}$ divides $\abs D$, and if $t=2$ the factor $\bigl(\frac2n\bigr)^2=1$ disappears and the period $4\abs{D'}=\abs D$. Evaluating at
$n=\abs D-1$ (so $n\equiv-1$ modulo $4$, modulo $\abs{D'}$, and modulo $8$ if $t\ge3$): $\bigl(\frac2n\bigr)^t=1$; $\bigl(\frac n{\abs{D'}}\bigr)=(-1)^{(\abs{D'}-1)/2}$ cancels the sign
$(-1)^{\frac{n-1}2\frac{\abs{D'}-1}2}=(-1)^{(\abs{D'}-1)/2}$; and $\bigl(\frac{\sgn D'}n\bigr)=\sgn D'$. So in all cases $\bigl(\frac D{\abs D-1}\bigr)=\sgn D$. Now $\bigl(\frac D{-n}\bigr)=\bigl(\frac D{-1}\bigr)\bigl(\frac Dn\bigr)=\sgn(D)\bigl(\frac Dn\bigr)$
by (i). Hence for $n>0$, $\bigl(\frac D{-n}\bigr)=\bigl(\frac D{\abs D-1}\bigr)\bigl(\frac Dn\bigr)=\bigl(\frac D{(\abs D-1)n}\bigr)$, and $(\abs D-1)n\equiv-n\pmod{\abs D}$: the symbol is periodic modulo $\abs D$
on $\Z\setminus\{0\}$; at multiples of $\abs D$ (including $0$) it vanishes, as $\abs D>1$. It is completely multiplicative and vanishes exactly off
$(\Z/\abs D)^\times$, so it comes from the homomorphism $(\Z/\abs D)^\times\to\{\pm1\}$ it induces.
(iv) is the definition. (v) By (i) for $n\ne0$; at $n=0$ both sides vanish unless $D=D'=1$. And $\bigl(\frac{Df^2}n\bigr)=\bigl(\frac Dn\bigr)\bigl(\frac fn\bigr)^2=\bigl(\frac Dn\bigr)$ when
$\gcd(f,n)=1$ and $n\ne0$; for $n=0$ the condition forces $f=\pm1$. The conductor statement is \Cref{prop:rings:conductor}(i) and (iii), the two characters
agreeing at the integers prime to $Df$.
\end{proof}

\begin{proposition}\label{prop:rings:qstar}
Let $q$ be an odd prime and $q^*=(-1)^{(q-1)/2}q$.
\begin{enumerate}[label=(\roman*)]
\item $q^*\equiv1\pmod4$; it is the prime discriminant attached to $q$.
\item For every $n\in\Z$, $\bigl(\frac{q^*}n\bigr)=\bigl(\frac nq\bigr)$ (the Legendre symbol of $n\bmod q$). So $\chi_{q^*}=\bigl(\frac\cdot q\bigr)$ as functions on $\Z$.
\item $\chi_{q^*}$ is the unique Dirichlet character of order $2$ and conductor $q$; this is why the paper calls it "the quadratic character of
conductor $q$". By contrast, for $q\equiv3\pmod4$ the function $n\mapsto\bigl(\frac qn\bigr)$ agrees on odd $n$ with the character $\chi_{4q}=\chi_{-4}\chi_{q^*}$, of
conductor $4q$, but on all of $\Z$ it is not periodic modulo $4q$: $\bigl(\frac q{4q+2}\bigr)=-\bigl(\frac q2\bigr)\ne0$. For $q=3$:
$\bigl(\frac32\bigr)=-1$ but $\bigl(\frac3{14}\bigr)=1$.
\item The same holds for the coefficient prime $\ell$: $\ell^*\equiv1\pmod4$, and for primes $p\ne\ell$, $\ell^*$ is a square modulo $p$ (for odd $p$), resp.
$\ell^*\equiv1\pmod8$ (for $p=2$), iff $p$ is a square modulo $\ell$.
\item $3^*=-3$, $5^*=5$, $7^*=-7$, $11^*=-11$, $13^*=13$; so the characters of conductor $3$, $5$, $7$, $11$, $13$ are $\chi_{-3}$, $\chi_5$, $\chi_{-7}$, $\chi_{-11}$,
$\chi_{13}$.
\end{enumerate}
\end{proposition}

\begin{proof}
(i) If $q\equiv1\pmod4$, $q^*=q$; if $q\equiv3$, $q^*=-q\equiv1\pmod4$. (ii) Both sides are completely multiplicative in $n\ne0$
(\Cref{prop:rings:kronecker}(i), \Cref{prop:rings:legendre}(ii)), so it suffices to compare them at $n=-1$, $2$, odd primes $p$, and to check $n=0$. At $-1$:
$\bigl(\frac{q^*}{-1}\bigr)=\sgn q^*=(-1)^{(q-1)/2}=\bigl(\frac{-1}q\bigr)$. At $2$: $\bigl(\frac{q^*}2\bigr)=1$ iff $q^*\equiv1\pmod8$ iff $q\equiv1\pmod8$ (when $q\equiv1\pmod4$) or $q\equiv7\pmod8$ (when
$q\equiv3\pmod4$), i.e. iff $q\equiv\pm1\pmod8$ iff $\bigl(\frac2q\bigr)=1$ (\Cref{thm:rings:qr}(iii)). At odd $p\ne q$ it is \Cref{thm:rings:qr}(i); at $p=q$ and at $n=0$
both sides vanish. (Checked numerically for all $q<200$ and $0<\abs n\le300$.)
(iii) By (ii), $\chi_{q^*}$ is a nontrivial character modulo the prime $q$, so its conductor is $q$ (the only other divisor, $1$, belongs to the
trivial character). The characters modulo $q$ form a group isomorphic to the cyclic group $(\Z/q)^\times$ (\Cref{prop:rings:dirichlet}), which has
a unique element of order $2$. For $q\equiv3\pmod4$, $q=(-1)\cdot q^*$, so for odd $n>0$, $\bigl(\frac qn\bigr)=\bigl(\frac{-1}n\bigr)\bigl(\frac{q^*}n\bigr)=\chi_{-4}(n)\chi_{q^*}(n)$ (for odd $n>0$,
$\bigl(\frac{-1}n\bigr)=(-1)^{(n-1)/2}=\bigl(\frac{-4}n\bigr)$), and $\chi_{-4}\chi_{q^*}=\chi_{4q}$ is primitive of conductor $4q$ by \Cref{cor:rings:fundamental-primitive}. For odd $n<0$ both sides equal their
values at $-n$: $\bigl(\frac q{-1}\bigr)=\sgn q=1$, and $\chi_{4q}(-1)=\chi_{-4}(-1)\chi_{q^*}(-1)=(-1)(-1)=1$ since $q^*<0$. Non-periodicity: $2q+1$ is odd, positive,
prime to $q$ and $\equiv3\pmod4$, so by \Cref{prop:rings:jacobi}(ii) $\bigl(\frac q{2q+1}\bigr)=-\bigl(\frac{2q+1}q\bigr)=-\bigl(\frac1q\bigr)=-1$, and
$\bigl(\frac q{4q+2}\bigr)=\bigl(\frac q2\bigr)\bigl(\frac q{2q+1}\bigr)=-\bigl(\frac q2\bigr)$, which is nonzero and differs from $\bigl(\frac q2\bigr)$. For $q=3$: $3\equiv3\pmod8$ gives
$\bigl(\frac32\bigr)=-1$, and $\bigl(\frac37\bigr)=-\bigl(\frac73\bigr)=-1$ gives $\bigl(\frac3{14}\bigr)=1$.
(iv) is (ii) with $q=\ell$, read at $n=p$. (v) is arithmetic.
\end{proof}

\begin{paperlink}\label{pl:rings:twist}
Lemma lem:twist begins: "Let $q$ be an odd prime, $\chi=\chi_{q^*}$ the quadratic character of conductor $q$ and $E^{(q)}$ the twist of $E$ by $\chi$, that
is, by $\Q(\sqrt{q^*})$." The sign in $q^*$ is forced: by \Cref{prop:rings:qstar}, $\chi_{q^*}=\bigl(\frac\cdot q\bigr)$ is ramified only at $q$, while for $q\equiv3\pmod4$ the
field $\Q(\sqrt q)$ has discriminant $4q$ and its character $\chi_{4q}$ is ramified at $2$ as well; the twist by it would in general change the reduction at
$2$ and violate part (i) of the lemma. The same character appears in Theorem thm:main through $\ell^*$ (the field $\Q(\sqrt{\ell^*})\subset\Q(\zeta_\ell)$, by
\Cref{prop:rings:quadratic-gauss-sum}(iii)), in Theorems thm:twist and thm:quant through $\chi_S=\prod_{q\in S}\chi_{q^*}$, and in Example ex:qual, where
the twists of 121a1 and 539c1 are by $\chi_{-11}$ and $\chi_{-7}$ ($11^*=-11$, $7^*=-7$). The earlier version of the paper wrote $\chi_q$ for
$\chi_{q^*}$; the revised version, like this book, does not.
\end{paperlink}

\section[Dirichlet characters and Gauss sums]{Dirichlet characters, conductors and Gauss sums; the paper's characters}\label{sec:rings:dirichlet}

The paper twists newforms by Dirichlet characters, $a_n(g\otimes\chi)=\chi(n)a_n(g)$ (Lemma lem:twist(iii), Theorems thm:twist and
thm:quant, Example ex:qual), and the level of the twist depends on the conductor of $\chi$ (Atkin-Li). This section defines
Dirichlet characters and their conductors, proves orthogonality and the basic properties of Gauss sums, and
assembles the paper's characters in one table (\Cref{prop:rings:paper-characters}). Their interpretation as characters of
$G_\Q$, with $\chi(\Frob_p)=\chi(p)$, is \Cref{ch:ant}; here they are functions on $\Z$.

\begin{definition}\label{def:rings:dirichlet}
A \emph{Dirichlet character modulo $n\ge1$} is a function $\chi\colon\Z\to\C$ of the form $\chi(a)=\tilde\chi(a\bmod n)$ for $\gcd(a,n)=1$ and $\chi(a)=0$ for
$\gcd(a,n)>1$, where $\tilde\chi\colon(\Z/n)^\times\to\C^\times$ is a homomorphism. The \emph{principal character} $\chi_0$ modulo $n$ has $\tilde\chi=1$. $\chi$ is \emph{even} or
\emph{odd} as $\chi(-1)=1$ or $-1$. The characters modulo $n$ form a group $X(n)$ under pointwise multiplication (of the $\tilde\chi$). Characters
with values in another field $k$ (for instance $\F_\ell$, $\F_9$) are defined in the same way with $\tilde\chi\colon(\Z/n)^\times\to k^\times$. The example used later is
$\omega_\ell\colon\Z\to\F_\ell$, $\omega_\ell(a)=a\bmod\ell$, an $\F_\ell^\times$-valued character modulo $\ell$ of order $\ell-1$: the Dirichlet character of the reduced
cyclotomic character $\bar\chi_\ell$ (\Cref{ch:galois,ch:ant}).
\end{definition}

\begin{proposition}\label{prop:rings:dirichlet}
Let $\chi$ be a Dirichlet character modulo $n$.
\begin{enumerate}[label=(\roman*)]
\item $\chi$ is completely multiplicative and $n$-periodic; for $\gcd(a,n)=1$, $\chi(a)$ is a root of unity whose order divides the exponent of
$(\Z/n)^\times$ (in particular divides $\varphi(n)$), and $\overline{\chi(a)}=\chi(a)^{-1}$; $\chi(-1)=\pm1$.
\item $X(n)$ is a finite abelian group isomorphic (non-canonically) to $(\Z/n)^\times$; in particular $\#X(n)=\varphi(n)$. For $a\in(\Z/n)^\times$, $a\ne1$, there
is $\chi\in X(n)$ with $\chi(a)\ne1$.
\end{enumerate}
\end{proposition}

\begin{proof}
(i) $\chi(ab)=\chi(a)\chi(b)$ when both are units modulo $n$, and both sides vanish otherwise ($\gcd(ab,n)>1$ iff $\gcd(a,n)>1$ or $\gcd(b,n)>1$).
If $a^e\equiv1$ then $\chi(a)^e=1$, so $\abs{\chi(a)}=1$ and $\bar z=z^{-1}$. $\chi(-1)^2=\chi(1)=1$.
(ii) For a finite abelian group $A$ let $\hat A=\Hom(A,\C^\times)$. If $A=\Z/m$, a homomorphism is determined by the image of $1$, an element of
$\mu_m(\C)$, so $\hat A\simeq\mu_m(\C)\simeq\Z/m$. For $A=A_1\times A_2$, $\hat A\simeq\hat A_1\times\hat A_2$ by restriction. By the structure theorem
(\Cref{ch:groups}), $A$ is a product of cyclic groups, so $\hat A\simeq A$. If $a\ne1$, some coordinate $a_i\in\Z/m_i$ of $a$ is nonzero, and the
character $x\mapsto\zeta_{m_i}^{x_i}$ is $\ne1$ at $a$. Apply this to $A=(\Z/n)^\times$.
\end{proof}

For $n=5$: $(\Z/5\Z)^\times$ is cyclic generated by $2$ ($2,4,3,1$ are its powers), so a character is determined by $\chi(2)\in\mu_4(\C)$ and there are four. The
choice $\chi(2)=i$ gives $\chi(4)=-1$, $\chi(3)=\chi(2^3)=-i$, $\chi(1)=1$: a character of order $4$, odd since $\chi(-1)=\chi(4)=-1$. The choice $\chi(2)=-1$ gives
the Legendre symbol $\bigl(\frac\cdot5\bigr)=\chi_5$, which is even. This is the non-canonical isomorphism of (ii): it depends on the choice of the generator $2$ and
of $i$ among the primitive fourth roots of unity.

\begin{theorem}[orthogonality]\label{thm:rings:orthogonality}
For Dirichlet characters modulo $n$:
\begin{enumerate}[label=(\roman*)]
\item $\sum_{a\bmod n}\chi(a)=\varphi(n)$ if $\chi=\chi_0$, and $0$ otherwise;
\item $\sum_{\chi\in X(n)}\chi(a)=\varphi(n)$ if $a\equiv1\pmod n$, and $0$ otherwise;
\item for $\gcd(a,n)=1$, $\frac1{\varphi(n)}\sum_\chi\overline{\chi(a)}\chi(b)$ is $1$ if $b\equiv a\pmod n$ and $0$ otherwise;
\item the characters modulo $n$ (restricted to $(\Z/n)^\times$) form an orthonormal basis of the space of functions on $(\Z/n)^\times$ for
$\langle f,h\rangle=\varphi(n)^{-1}\sum_{a\in(\Z/n)^\times}f(a)\overline{h(a)}$.
\end{enumerate}
\end{theorem}

\begin{proof}
(i) If $\chi\ne\chi_0$, pick $b$ with $\chi(b)\ne1$; $a\mapsto ab$ permutes $(\Z/n)^\times$, so $S=\sum\chi(a)=\sum\chi(ab)=\chi(b)S$, and $S=0$. (ii) If $a$ is not a
unit every term vanishes. If $a\not\equiv1$ is a unit, pick $\psi$ with $\psi(a)\ne1$ (\Cref{prop:rings:dirichlet}(ii)); $\chi\mapsto\psi\chi$ permutes $X(n)$, so
$T=\sum\chi(a)=\psi(a)T$ and $T=0$. (iii) Apply (ii) to $a^{-1}b$, using $\overline{\chi(a)}=\chi(a^{-1})$. (iv) $\langle\chi,\psi\rangle=\varphi(n)^{-1}\sum(\chi\bar\psi)(a)$, which is $1$ or
$0$ by (i) applied to $\chi\psi^{-1}$. Orthonormal vectors are independent, and there are $\varphi(n)$ of them in a space of dimension $\varphi(n)$.
\end{proof}

This is the character table of the finite abelian group $(\Z/n)^\times$, with its row and column orthogonality as in
Riley-Hobson-Bence; the expansion of a function in characters is finite Fourier analysis on $(\Z/n)^\times$.

\begin{definition}\label{def:rings:conductor}
Let $\chi$ be a character modulo $n$ and $d\mid n$. $\chi$ is \emph{defined modulo $d$} (induced from a character $\chi'$ modulo $d$) if $\chi(a)=\chi'(a)$ for
all $a$ with $\gcd(a,n)=1$. The \emph{conductor} $f_\chi$ is the least $d\mid n$ modulo which $\chi$ is defined; $\chi$ is \emph{primitive} if $f_\chi=n$.
\end{definition}

A modulus is a choice, the conductor is not: the same function on the integers prime to $12$ can be presented as a character modulo $4$,
$12$ or $24$, and the conductor records the least modulus that is really needed. For example the character modulo $12$ given by
$a\mapsto(-1)^{(a-1)/2}$ on units ($1,5\mapsto1$; $7,11\mapsto-1$) is induced from $\chi_{-4}$ modulo $4$, so its conductor is $4$. The conductor is what
matters in the paper: the level of a twist $g\otimes\chi$ and the ramification of the corresponding Galois character are governed by $f_\chi$,
not by the modulus used to write $\chi$ down.

\begin{proposition}\label{prop:rings:conductor}
Let $\chi$ be a character modulo $n$.
\begin{enumerate}[label=(\roman*)]
\item For $d\mid n$ the reduction $(\Z/n)^\times\to(\Z/d)^\times$ is surjective, and $\chi$ is defined modulo $d$ iff $\chi(a)=1$ for all $a$ with $\gcd(a,n)=1$ and
$a\equiv1\pmod d$; the inducing $\chi'$ is then unique.
\item If $\chi$ is defined modulo $d_1$ and modulo $d_2$, it is defined modulo $\gcd(d_1,d_2)$.
\item $f_\chi$ divides every $d$ modulo which $\chi$ is defined. $\chi$ is induced by a unique character $\chi^{\mathrm{prim}}$ modulo $f_\chi$, and $\chi^{\mathrm{prim}}$ is primitive.
\item $\chi$ and $\chi^{\mathrm{prim}}$ agree at integers prime to $n$ but may differ at $a$ with $\gcd(a,n)>1=\gcd(a,f_\chi)$: $\chi_{-3}$ regarded modulo $6$ vanishes at $2$,
while $\chi_{-3}(2)=-1$.
\item Write $n=\prod_pp^{k_p}$. Then $\chi=\prod_p\chi^{(p)}$ with unique characters $\chi^{(p)}$ modulo $p^{k_p}$ (the \emph{$p$-components} of $\chi$), and
$f_\chi=\prod_pf_{\chi^{(p)}}$. A product of primitive
characters with pairwise coprime conductors is primitive, of conductor the product.
\item The principal character has conductor $1$; the nontrivial character modulo $4$, $\chi_{-4}$, has conductor $4$ (not $2$); $\bigl(\frac\cdot q\bigr)$ has conductor
$q$; $\chi_8$ ($=1$ iff $n\equiv\pm1\pmod8$) and $\chi_{-8}$ ($=1$ iff $n\equiv1,3\pmod8$) have conductor $8$; for a finite set $S$ of odd primes,
$\chi_S=\prod_{q\in S}\chi_{q^*}$ is primitive of conductor $D=\prod_{q\in S}q$ and vanishes exactly at the $n$ with $\gcd(n,D)>1$.
\end{enumerate}
\end{proposition}

\begin{proof}
(i) Given $b$ prime to $d$, let $n'$ be the product of the primes dividing $n$ but not $d$; by the Chinese remainder theorem there is $a$ with
$a\equiv b\pmod d$ and $a\equiv1\pmod{n'}$, and $a$ is prime to every prime dividing $n$. So the reduction is onto, with kernel
$\{a\equiv1\pmod d\}$. $\chi$ is defined modulo $d$ iff $\tilde\chi$ factors through this surjection, iff it is trivial on the kernel; the factorisation
is unique because the map is onto.
(ii) Let $g=\gcd(d_1,d_2)$, $L=\lcm(d_1,d_2)$, and $a$ prime to $n$ with $a\equiv1\pmod g$. Since $1\equiv a\pmod g$ and $d_1\Z+d_2\Z=g\Z$, write $a-1=d_1u+d_2w$; then
$b_0=1+d_1u$ satisfies $b_0\equiv1\pmod{d_1}$ and $b_0\equiv a\pmod{d_2}$. Adjust $b_0$ by a multiple of $L$ so that $b\equiv1$ modulo every prime dividing $n$ but not
$L$; then $b$ is prime to $n$ (primes dividing $d_1$ do not divide $b\equiv1$, primes dividing $d_2$ do not divide $b\equiv a$). Now $\chi(b)=1$ as
$b\equiv1\pmod{d_1}$, and $\chi(ab^{-1})=1$ as $ab^{-1}\equiv1\pmod{d_2}$ ($b^{-1}$ an inverse modulo $n$). So $\chi(a)=1$, and (i) applies.
(iii) If $\chi$ is defined modulo $d$, it is defined modulo $\gcd(d,f_\chi)\le f_\chi$, which must equal $f_\chi$, so $f_\chi\mid d$. $\chi^{\mathrm{prim}}$ is the $\chi'$ of (i)
for $d=f_\chi$. If $\chi^{\mathrm{prim}}$ were defined modulo $d'\mid f_\chi$, $d'<f_\chi$, so would be $\chi$ (integers prime to $n$ are prime to $f_\chi$).
(iv) By definition, and $2\mid6$ while $\chi_{-3}(2)=\bigl(\frac{-3}2\bigr)=-1$ as $-3\equiv5\pmod8$.
(v) By the Chinese remainder theorem $(\Z/n)^\times=\prod_p(\Z/p^{k_p})^\times$, and $\tilde\chi$ is the product of its restrictions $\tilde\chi^{(p)}$ to the factors.
Under this identification $\{a\equiv1\pmod d\}$ corresponds to $\prod_p\{a_p\equiv1\pmod{p^{v_p(d)}}\}$, and $\chi$ is trivial on it iff each $\chi^{(p)}$ is trivial
on its factor. So $\chi$ is defined modulo $d$ iff each $\chi^{(p)}$ is defined modulo $p^{v_p(d)}$, and the least such $d$ is $\prod_pf_{\chi^{(p)}}$. If
$\chi=\prod_i\chi^{(i)}$ with $\chi^{(i)}$ primitive of pairwise coprime conductors $f_i$, regard it modulo $\prod f_i$; its $p$-component for $p\mid f_i$ is
the $p$-component $(\chi^{(i)})^{(p)}$ of $\chi^{(i)}$, so $f_\chi=\prod_i\prod_{p\mid f_i}f_{(\chi^{(i)})^{(p)}}=\prod_if_{\chi^{(i)}}=\prod_if_i$ (the second equality is the
first statement of (v) applied to $\chi^{(i)}$).
(vi) The trivial character is defined modulo $1$. $(\Z/2)^\times$ is trivial, so a nontrivial character modulo $4$ is not defined modulo $2$.
$\bigl(\frac\cdot q\bigr)$ is nontrivial modulo the prime $q$. $\chi_8(5)=\bigl(\frac85\bigr)=\bigl(\frac35\bigr)=-1$ and $\chi_{-8}(5)=\bigl(\frac{-8}5\bigr)=\bigl(\frac25\bigr)=-1$ although $5\equiv1\pmod4$, so neither is
defined modulo $4$; their values at odd $n>0$ are $\bigl(\frac2n\bigr)^3=\bigl(\frac2n\bigr)$ and $\bigl(\frac{-1}n\bigr)\bigl(\frac2n\bigr)$, read off from \Cref{prop:rings:jacobi}(ii). For $\chi_S$ apply (v) to the primitive
characters $\chi_{q^*}$ of the distinct primes $q\in S$ (\Cref{prop:rings:qstar}); a product of characters vanishes where one factor does.
\end{proof}

\begin{corollary}\label{cor:rings:fundamental-primitive}
A fundamental discriminant $D\ne1$ is, uniquely, a product of pairwise coprime prime discriminants ($-4$, $8$, $-8$, $q^*$), and $\chi_D$ is the
product of the corresponding characters. It is primitive of conductor $\abs D$.
\end{corollary}

\begin{proof}
If $D\equiv1\pmod4$ is squarefree, $D=\pm\prod q_i$ with distinct odd primes; $\prod q_i^*=\pm\prod q_i$ is $\equiv1\pmod4$, and of the two numbers $\pm\prod q_i$ only
one is $\equiv1\pmod4$, so $D=\prod q_i^*$. If $D=4m$ with $m\equiv3\pmod4$ squarefree, $-m\equiv1\pmod4$ is squarefree and $D=(-4)\cdot(-m)$. If $D=4m$ with
$m=2m'$, $m'$ odd squarefree, then $D=8m'$ if $m'\equiv1\pmod4$ and $D=(-8)(-m')$ if $m'\equiv3\pmod4$. The odd prime discriminants occurring are the
$q^*$ for the odd $q\mid D$, and the even one is then determined. By \Cref{prop:rings:kronecker}(i) (multiplicativity in the upper entry,
all entries nonzero), $\chi_D$ is the product of the characters of the factors; these are primitive of conductors $4$, $8$, $8$, $q$
(\Cref{prop:rings:conductor}(vi)), pairwise coprime, so $\chi_D$ is primitive of conductor $\abs D$ by \Cref{prop:rings:conductor}(v). (For the $39$
fundamental discriminants with $\abs D\le60$ the conductor was also found by brute force to be $\abs D$.)
\end{proof}

\begin{theorem}[Gauss sums]\label{thm:rings:gauss-sums}
For a character $\chi$ modulo $n$ and $b\in\Z$ put
\[
\tau(\chi,b)=\sum_{a\bmod n}\chi(a)\zeta_n^{ab},\qquad\tau(\chi)=\tau(\chi,1).
\]
\begin{enumerate}[label=(\roman*)]
\item If $\gcd(b,n)=1$, $\tau(\chi,b)=\bar\chi(b)\tau(\chi)$. Always $\overline{\tau(\chi)}=\chi(-1)\tau(\bar\chi)$.
\item If $\chi$ is primitive, $\tau(\chi,b)=\bar\chi(b)\tau(\chi)$ for all $b$ (both sides vanish when $\gcd(b,n)>1$), $\abs{\tau(\chi)}^2=n$, and $\tau(\chi)\tau(\bar\chi)=\chi(-1)n$.
\item $\tau\bigl(\bigl(\frac\cdot q\bigr)\bigr)^2=q^*$ for odd primes $q$; $\tau(\chi_{-4})=i-(-i)=2i$; $\tau(\chi_{-3})=\zeta_3-\zeta_3^2=1+2\zeta_3=\sqrt{-3}$.
\item If $\chi$ modulo $n$ is induced from the primitive $\chi^{\mathrm{prim}}$ modulo $f$ and $\gcd(n/f,f)=1$, then
$\tau(\chi)=\mu(n/f)\chi^{\mathrm{prim}}(n/f)\tau(\chi^{\mathrm{prim}})$. For instance $\chi_{-3}$ modulo $6$ has $\tau=(-1)(-1)\sqrt{-3}=\sqrt{-3}$, and the principal character modulo $n$
has $\tau(\chi_0)=\mu(n)$, the sum of the primitive $n$-th roots of unity.
\end{enumerate}
\end{theorem}

\begin{proof}
(i) Substitute $c=ab$, which runs over $(\Z/n)^\times$ with $a$: $\chi(a)=\chi(c)\bar\chi(b)$. Then $\overline{\tau(\chi)}=\sum\bar\chi(a)\zeta_n^{-a}=\tau(\bar\chi,-1)=\chi(-1)\tau(\bar\chi)$.
(ii) Let $\gcd(b,n)=g>1$, $n=gn'$, $b=gb'$; then $\zeta_n^{ab}=\zeta_{n'}^{ab'}$ depends only on $a\bmod n'$. As $\chi$ is not defined modulo $n'$, there is $c$ prime to
$n$ with $c\equiv1\pmod{n'}$ and $\chi(c)\ne1$ (\Cref{prop:rings:conductor}(i)). Substituting $a\mapsto ac$: $\tau(\chi,b)=\sum\chi(ac)\zeta_{n'}^{acb'}=\chi(c)\tau(\chi,b)$, so
$\tau(\chi,b)=0=\bar\chi(b)\tau(\chi)$. Then
\[
\varphi(n)\abs{\tau(\chi)}^2=\sum_{b\bmod n}\abs{\tau(\chi,b)}^2=\sum_{a,a'}\chi(a)\overline{\chi(a')}\sum_b\zeta_n^{(a-a')b}=n\sum_a\abs{\chi(a)}^2=n\varphi(n),
\]
using $\sum_b\zeta_n^{cb}=n[n\mid c]$. Finally $n=\tau(\chi)\overline{\tau(\chi)}=\chi(-1)\tau(\chi)\tau(\bar\chi)$ by (i), and $\chi(-1)^2=1$.
(iii) The first is \Cref{prop:rings:quadratic-gauss-sum}(i) with $\zeta=\zeta_q$. $\chi_{-4}$ is $1$ at $1$ and $-1$ at $3$, and $\zeta_4=i$. $\chi_{-3}=\bigl(\frac\cdot3\bigr)$ is $1$ at
$1$ and $-1$ at $2$, and $\zeta_3-\zeta_3^2=\zeta_3+1+\zeta_3=1+2\zeta_3$.
(iv) Put $m=n/f$, choose $u,w$ with $um+wf=1$, so that $\zeta_n=\zeta_n^{um+wf}=\zeta_f^u\zeta_m^w$. For $\gcd(a,n)=1$ we have $\chi(a)=\chi^{\mathrm{prim}}(a)$, and $\gcd(a,n)=1$ iff
$\gcd(a,f)=\gcd(a,m)=1$; so $\chi(a)=\chi^{\mathrm{prim}}(a)\cdot[\gcd(a,m)=1]$ for all $a$. By the Chinese remainder theorem $a\bmod n\leftrightarrow(a\bmod f,a\bmod m)$,
\[
\tau(\chi)=\sum_{x\bmod f}\chi^{\mathrm{prim}}(x)\zeta_f^{ux}\cdot\sum_{y\in(\Z/m)^\times}\zeta_m^{wy}=\tau(\chi^{\mathrm{prim}},u)\cdot c_m,
\]
where $c_m$ is the sum of the primitive $m$-th roots of unity ($w$ is prime to $m$, as $wf\equiv1$, so $y\mapsto wy$ permutes $(\Z/m)^\times$). Since $\sum_{d\mid m}c_d$ is the
sum of all $m$-th roots of unity, which is $1$ for $m=1$ and $0$ otherwise (the coefficient of $x^{m-1}$ in $x^m-1$), Moebius inversion gives $c_m=\mu(m)$.
And $um\equiv1\pmod f$, so $u$ is prime to $f$ and $\tau(\chi^{\mathrm{prim}},u)=\overline{\chi^{\mathrm{prim}}(u)}\tau(\chi^{\mathrm{prim}})=\chi^{\mathrm{prim}}(m)\tau(\chi^{\mathrm{prim}})$ by (i). For $\chi_{-3}$ modulo $6$: $f=3$,
$m=2$, $\mu(2)=-1$, $\chi_{-3}(2)=-1$. For $\chi_0$: $f=1$, $\chi^{\mathrm{prim}}=1$ with $\tau=1$, $m=n$.
\end{proof}

All of (i) and (ii) were also checked numerically for the $278$ primitive characters of modulus $3\le n<40$. Gauss sums are the finite
Fourier transforms of characters; they reappear in the root numbers of twisted $L$-functions (\Cref{ch:lfunctions}).

\begin{proposition}[the paper's characters]\label{prop:rings:paper-characters}
The characters used in the paper are the following Dirichlet characters ($q$, $\ell$ odd primes, $S$ a finite set of odd primes):
\begin{center}
\footnotesize\setlength{\tabcolsep}{3pt}
\begin{tabular}{lll>{\raggedright\arraybackslash}p{2cm}>{\raggedright\arraybackslash}p{4.6cm}}
\toprule
name & Kronecker symbol & conductor & $\chi(-1)$ & value at a prime $p$\\
\midrule
$\chi_{-4}$ & $\bigl(\frac{-4}\cdot\bigr)$ & $4$ & $-1$ & $0$ at $p=2$; $(-1)^{(p-1)/2}$, i.e. $1$ iff $p\equiv1\pmod4$\\
$\chi_{-3}$ & $\bigl(\frac{-3}\cdot\bigr)=\bigl(\frac\cdot3\bigr)$ & $3$ & $-1$ & $0$ at $p=3$; $1$ iff $p\equiv1\pmod3$ (so $\chi_{-3}(2)=-1$)\\
$\chi_{q^*}$ & $\bigl(\frac{q^*}\cdot\bigr)=\bigl(\frac\cdot q\bigr)$ & $q$ & $(-1)^{(q-1)/2}$ & $\bigl(\frac pq\bigr)$ ($0$ at $p=q$; at $p=2$ also $\bigl(\frac2q\bigr)$)\\
$\chi_S=\prod_{q\in S}\chi_{q^*}$ & $\bigl(\frac{\prod_{q\in S}q^*}\cdot\bigr)$ & $\prod_{q\in S}q$ & $\sgn\prod_{q\in S}q^*$ & $\prod_{q\in S}\bigl(\frac pq\bigr)$\\
$\chi_{\ell^*}$ & $\bigl(\frac{\ell^*}\cdot\bigr)=\bigl(\frac\cdot\ell\bigr)$ & $\ell$ & $(-1)^{(\ell-1)/2}$ & $\bigl(\frac p\ell\bigr)$\\
$\omega_\ell$ ($\F_\ell^\times$-valued) & none & $\ell$ & $-1$ & $p\bmod\ell$\\
\bottomrule
\end{tabular}
\end{center}
Moreover:
\begin{enumerate}[label=(\roman*)]
\item $\chi_{-3}(a)\equiv a\pmod3$ for every $a\in\Z$, and $\chi_{-4}(a)\equiv a\pmod4$ for every odd $a$. In particular $p\equiv\chi_{-3}(p)\pmod3$ for every prime
$p\ne3$, and $p\equiv\chi_{-4}(p)\pmod4$ for every odd prime $p$.
\item For an element $\sigma$ of order $3$ of any group and $m$ prime to $3$, $\sigma^m=\sigma^{\chi_{-3}(m)}$; for $\sigma$ of order dividing $4$ and $m$ odd,
$\sigma^m=\sigma^{\chi_{-4}(m)}$.
\item As $\{\pm1\}\subset\F_\ell^\times$-valued functions, $\bigl(\frac\cdot\ell\bigr)=\omega_\ell^{(\ell-1)/2}$, and $\bigl(\frac\cdot\ell\bigr)=\omega_\ell$ iff $\ell=3$; so $\chi_{-3}=\omega_3$.
\item For primes $q\ge5$: $3\mid q-\chi_{-4}(q)$ iff $q\equiv\pm1\pmod{12}$.
\end{enumerate}
\end{proposition}

\begin{proof}
The rows: $\chi_{-4}$ and $\chi_{\pm8}$ are discussed in \Cref{prop:rings:conductor}(vi); for odd $a>0$, $\chi_{-4}(a)=\bigl(\frac{-1}a\bigr)\bigl(\frac2a\bigr)^2=(-1)^{(a-1)/2}$ by
\Cref{prop:rings:jacobi}(ii), and $\chi_{-4}(2)=0$; parity by \Cref{prop:rings:kronecker}(iii). The rows $\chi_{-3}$, $\chi_{q^*}$, $\chi_{\ell^*}$ are \Cref{prop:rings:qstar}
($3^*=-3$) with \Cref{prop:rings:legendre}(iv); $\chi_S$ is \Cref{prop:rings:conductor}(vi), with $\chi_S=\chi_{\prod q^*}$ by \Cref{prop:rings:kronecker}(v);
$\omega_\ell$ is defined on $(\Z/\ell)^\times=\F_\ell^\times$, which it maps isomorphically to itself, so it has order $\ell-1$, conductor $\ell$ and $\omega_\ell(-1)=-1$.
(i) $\chi_{-3}=\bigl(\frac\cdot3\bigr)$ and \Cref{prop:rings:legendre}(iv); $\chi_{-4}(a)=(-1)^{(a-1)/2}\equiv a\pmod4$ for odd $a>0$, and both sides are periodic modulo $4$.
(ii) \Cref{prop:rings:mu-in-fq}(v) and (i).
(iii) Euler's criterion (\Cref{prop:rings:legendre}(i)) is $\bigl(\frac a\ell\bigr)=\omega_\ell(a)^{(\ell-1)/2}$ in $\F_\ell$. The character $\omega_\ell$ has order $\ell-1$ and $\bigl(\frac\cdot\ell\bigr)$ has
order $2$, so they are equal only if $\ell-1=2$; for $\ell=3$, $\omega_3(a)=a\bmod3=\bigl(\frac a3\bigr)$ by (i).
(iv) If $q\equiv1\pmod4$, $\chi_{-4}(q)=1$ and the condition is $q\equiv1\pmod3$, i.e. $q\equiv1\pmod{12}$; if $q\equiv3\pmod4$, $\chi_{-4}(q)=-1$ and the condition is
$q\equiv2\pmod3$, i.e. $q\equiv11\pmod{12}$.
\end{proof}

\begin{warning}[names of characters]\label{warn:rings:names}
(a) $\chi_\ell$, with a letter subscript, is the $\ell$-adic cyclotomic character, not a quadratic character. At $\ell=3$ the two readings of "$\chi_3$"
would collide: the quadratic $\chi_{-3}$ and the cyclotomic $\chi_3^{\mathrm{cyc}}$, which agree only modulo $3$ ($\bar\chi_3^{\mathrm{cyc}}=\chi_{-3}$, the Galois form of
(iii) above). The book therefore never writes $\chi_3$, $\chi_7$ or $\chi_{11}$. (b) $\omega_\ell$ and $\bigl(\frac\cdot\ell\bigr)=\chi_{\ell^*}$ are different characters for
$\ell\ge5$: the first has order $\ell-1$, the second order $2$. This is the point of the correction to Lemma lem:h0 and to the paper's
Section 1.3: the relevant condition is $\bar\rho\simeq\bar\rho\otimes\bar\chi_\ell$, not $\bar\rho\simeq\bar\rho\otimes\chi_{\ell^*}$. (c) The character of a quadratic field $\Q(\sqrt d)$ is
$\chi_D$ with $D$ the discriminant of the field: the character of $\Q(i)$ is $\chi_{-4}$ and that of $\Q(\sqrt2)$ is $\chi_8$ (Appendix B, B-C6(b), also writes
them with the radicand as index; this chapter does not) (\Cref{ch:ant}). (d) Do not confuse the Kronecker
symbol $\bigl(\frac D2\bigr)$ (lower entry $2$) with the Jacobi symbol $\bigl(\frac2D\bigr)$: $\chi_{-3}(2)=\bigl(\frac{-3}2\bigr)=-1$ because $-3\equiv5\pmod8$, consistent with $2\equiv-1\pmod3$.
\end{warning}

\begin{paperlink}\label{pl:rings:characters}
The table of Section 1 uses $\chi_{-3}$ and $\chi_{-4}$ "the quadratic characters of $\Q(\sqrt{-3})$ and $\Q(i)$", whose values at the bad prime $q$
are read from the table above. The proof of Proposition prop:forms says "$e_3(E,q)=v_3(q-\chi_{-3}(q))\ge1=w_3(E,q)$ since $q\equiv\chi_{-3}(q)\pmod3$":
this is (i). Theorem thm:twist uses $\chi_S$ of conductor $D=\prod_{q\in S}q$ and, in its proof, that the twist map
$\iota\colon h\mapsto\sum\chi_S(n)a_n(h)q^n$ kills the coefficients at the $n$ with $\gcd(n,D)>1$: \Cref{prop:rings:conductor}(vi) (here $q^n$ is the
$q$-expansion variable, a different meaning of $q$). Section sec:anomaly observes that "the residue classes of the type $\mathrm{III}$
primes of the violating curves are $q\equiv1$ and $q\equiv11\pmod{12}$": by (iv) these are exactly the classes with $e_3(E,q)>0$. Lemma lem:h0 and
Proposition prop:ns use $\bar\chi_3^{\mathrm{cyc}}=\chi_{-3}$, whose Dirichlet side is (iii). In Example ex:qual, $a_p(\text{121d})=\chi_{-11}(p)a_p(\text{11a})$ for
$p\ne11$.
\end{paperlink}

\section{Valuations and the paper's local terms}\label{sec:rings:valuations}

Every number in the paper's main statements is an $\ell$-adic valuation: $v_\ell(\degphi)$, $t_\ell(E,q)=v_\ell(-v_q(j_E))$, and the Euler terms
$e_\ell(E,q)=v_\ell\bigl(P_q(q^{-2})/(1-a_q(E)^2q^{-2})\bigr)$ of Definition def:terms, which the table of Section 1 evaluates as $v_\ell(q-\chi_{-3}(q))$,
$v_\ell(q-\chi_{-4}(q))$, $v_\ell(q^2-1)$ and $v_\ell(q-1)+v_\ell((q+1)^2-a_q(E')^2)$. This section defines $v_p$ and $\abs\cdot_p$ on $\Q$ (the completions $\Q_p$, $\Z_p$
are \Cref{ch:local}), shows that the rationals in question are $\ell$-integral, proves the algebraic identity behind the $\kod{I}_0^*$ row,
and works out when each entry of the table is positive. \Cref{ch:groups} used $v_\ell(n)$ for integers informally; the formal
definition is the following.

\begin{definition}\label{def:rings:valuation}
For a prime $p$ and $x\in\Q^\times$ write $x=p^ka/b$ with $a,b\in\Z$ prime to $p$ and $k\in\Z$; then $v_p(x):=k$. Put $v_p(0)=+\infty$. (The exponent $k$ is well
defined: if $p^ka/b=p^{k'}a'/b'$ then $p^kab'=p^{k'}a'b$ in $\Z$, and unique factorisation gives $k=k'$.)
\end{definition}

For example $v_3(\frac{18}5)=2$, $v_5(\frac{18}5)=-1$ and $v_2(\frac{18}5)=1$. The valuation measures divisibility, and it is used throughout because the
paper's quantities are rationals whose numerators and denominators matter prime by prime: $v_\ell(x)>0$ says that $\ell$ divides $x$ in
$\Z_{(\ell)}$.

\begin{proposition}\label{prop:rings:valuation}
Let $p$ be a prime and $x,y\in\Q$.
\begin{enumerate}[label=(\roman*)]
\item $v_p(xy)=v_p(x)+v_p(y)$.
\item $v_p(x+y)\ge\min(v_p(x),v_p(y))$, with equality if $v_p(x)\ne v_p(y)$.
\item $v_p(x)\ge0$ iff $x\in\Z_{(p)}$, and $v_p(x)=0$ iff $x\in\Z_{(p)}^\times$ (\Cref{prop:rings:zp-local}).
\item On $\Z$: $p^k\mid n$ iff $v_p(n)\ge k$; $v_p(\gcd(m,n))=\min(v_p(m),v_p(n))$ and $v_p(\lcm(m,n))=\max(v_p(m),v_p(n))$ (\Cref{prop:rings:integers}).
\item For ideals, $v_p(r\Z)=v_p(r)$ is well defined; the $p$-part of $n\ne0$ is $p^{v_p(n)}$; and $v_p(\ell)=1$ if $p=\ell$, $0$ otherwise.
\item (Legendre's formula) $v_p(n!)=\sum_{k\ge1}\lfloor n/p^k\rfloor$.
\end{enumerate}
\end{proposition}

\begin{proof}
(i) Multiply the representations. (ii) If $x=p^ka/b$, $y=p^{k'}a'/b'$ with $k\le k'$, then $x+y=p^k(ab'+p^{k'-k}a'b)/(bb')$, whose numerator is an integer and
denominator prime to $p$; if $k<k'$ the numerator is $\equiv ab'\not\equiv0\pmod p$. (The cases with $x$ or $y$ zero are trivial.) (iii) is
\Cref{prop:rings:zp-local}(iv). (iv) is \Cref{prop:rings:integers}(iv). (v) $r\Z=r'\Z$ iff $r'=\pm r$. (vi) $v_p(n!)=\sum_{m=1}^nv_p(m)=\sum_{m=1}^n\sum_{k\ge1}[p^k\mid m]$, and
exchanging the sums, $\#\{m\le n:p^k\mid m\}=\lfloor n/p^k\rfloor$.
\end{proof}

The inequality in (ii) can be strict when the valuations agree: $v_3(\frac13)=v_3(\frac23)=-1$, but $\frac13+\frac23=1$ has $v_3=0$. Legendre's formula gives, for
instance, $v_3(10!)=3+1=4$, and $10!=3628800=2^8\cdot3^4\cdot5^2\cdot7$.

\begin{definition}\label{def:rings:padic-abs}
The \emph{$p$-adic absolute value} on $\Q$ is $\abs x_p=p^{-v_p(x)}$, $\abs0_p=0$; $\abs x_\infty$ is the usual absolute value. An \emph{absolute value} on a field is
a function $\abs\cdot\colon K\to\R_{\ge0}$ with $\abs x=0$ iff $x=0$, $\abs{xy}=\abs x\abs y$, and $\abs{x+y}\le\abs x+\abs y$; two are \emph{equivalent} if one is a positive power of
the other.
\end{definition}

\begin{proposition}\label{prop:rings:padic-abs}
\begin{enumerate}[label=(\roman*)]
\item $\abs\cdot_p$ is an absolute value satisfying the ultrametric inequality $\abs{x+y}_p\le\max(\abs x_p,\abs y_p)$; $d_p(x,y)=\abs{x-y}_p$ is a metric on $\Q$, in
which $p^k\to0$.
\item (Product formula) $\prod_{p\le\infty}\abs x_p=1$ for $x\in\Q^\times$ (the product over all primes and $\infty$, almost all factors being $1$).
\item $\Q$ is not complete for $d_p$.
\end{enumerate}
\end{proposition}

\begin{proof}
(i) is \Cref{prop:rings:valuation}(i), (ii); the ultrametric inequality implies the triangle inequality, and $\abs{p^k}_p=p^{-k}$.
(ii) Write $x=\pm\prod_pp^{v_p(x)}$ (\Cref{prop:rings:integers}(iii) applied to numerator and denominator); then $\abs x_\infty=\prod_pp^{v_p(x)}$ and
$\prod_p\abs x_p=\prod_pp^{-v_p(x)}$, the inverse. For instance, for $x=-\frac{12}5$: $\abs x_\infty=\frac{12}5$, $\abs x_2=\frac14$, $\abs x_3=\frac13$, $\abs x_5=5$, and the product
is $\frac{12}5\cdot\frac1{12}\cdot5=1$.
(iii) Choose $m$ and $c$ as follows: for odd $p$, $m=2$ and $c=1+2p$, which is not a square (if $1+2p=s^2$ then $2p=(s-1)(s+1)$ with both factors
even, so $4\mid2p$); for $p=2$, $m=3$ and $c=3$, not a cube (a rational cube root of an integer is an integer, \Cref{thm:rings:gauss-lemma}(iv), and
$1<3<8$). In both cases $c\equiv1\pmod p$ and $p\nmid m$. Starting from $x_1=1$, define integers $x_k\equiv1\pmod p$ with $x_k^m\equiv c\pmod{p^k}$: given $x_k$, choose
$t\in\Z$ with $mx_k^{m-1}t\equiv(c-x_k^m)/p^k\pmod p$ (possible as $mx_k^{m-1}\not\equiv0\pmod p$) and put $x_{k+1}=x_k+tp^k$; the binomial expansion gives
$x_{k+1}^m\equiv x_k^m+mx_k^{m-1}tp^k\equiv c\pmod{p^{k+1}}$. Then $\abs{x_{k+1}-x_k}_p\le p^{-k}$, so $(x_k)$ is Cauchy by the ultrametric inequality. If $x_k\to x\in\Q$, then
$\abs{x_k^m-x^m}_p\le\abs{x_k-x}_p\max(1,\abs x_p)^{m-1}\to0$ and $\abs{x_k^m-c}_p\le p^{-k}\to0$, so $x^m=c$, a contradiction. (The completion $\Q_p$ contains the limit;
\Cref{ch:local}.)
\end{proof}

\begin{theorem}[Ostrowski]\label{thm:rings:ostrowski}
Every nontrivial absolute value on $\Q$ is equivalent to $\abs\cdot_\infty$ or to $\abs\cdot_p$ for exactly one prime $p$.
\end{theorem}

We do not prove this; see \cite[Chapter II, Section 3]{Neukirch}. It is not used in the paper; it explains why the places of $\Q$ are
the primes and $\infty$.

\begin{lemma}\label{lem:rings:ell-integral}
Let $q\ne\ell$ be primes.
\begin{enumerate}[label=(\roman*)]
\item $v_\ell(q^k)=0$ for all $k\in\Z$: $q$ is a unit of $\Z_{(\ell)}$ (an "$\ell$-adic unit").
\item For $P\in\Z[X]$ of degree at most $d$ and $k\ge0$, $P(q^{-k})\in\Z_{(\ell)}$, and $v_\ell(P(q^{-k}))=v_\ell\bigl(q^{kd}P(q^{-k})\bigr)$ is the valuation of the integer
$q^{kd}P(q^{-k})$; it is finite iff $P(q^{-k})\ne0$.
\item For $a\in\Z$ and either choice of sign:
\begin{gather*}
v_\ell\Bigl(\frac{q-1}q\Bigr)=v_\ell(q-1),\qquad v_\ell(1\mp q^{-1})=v_\ell(q\mp1),\\
v_\ell(1-q^{-2})=v_\ell(q^2-1)=v_\ell(q-1)+v_\ell(q+1),
\end{gather*}
\[
v_\ell\Bigl(\frac{q-1}q\cdot\frac{(q+1)^2-a^2}{q^2}\Bigr)=v_\ell(q-1)+v_\ell\bigl((q+1)^2-a^2\bigr)\quad\text{when }(q+1)^2\ne a^2.
\]
All these valuations are $\ge0$.
\item For $q=2$, $\ell=3$: $v_3(1+2\cdot2^{-2})=v_3(3/2)=1$.
\end{enumerate}
\end{lemma}

\begin{proof}
(i) $v_\ell(q)=0$ and \Cref{prop:rings:valuation}(i). (ii) $q^{kd}P(q^{-k})=\sum_ia_iq^{k(d-i)}\in\Z$ for $P=\sum_{i\le d}a_iX^i$, and $v_\ell(q^{-kd})=0$. (iii) Each rational
is an integer divided by a power of $q$; apply (i), \Cref{prop:rings:valuation}(i) and \Cref{lem:rings:lte}(v). (iv) $1+\frac12=\frac32$.
\end{proof}

\begin{proposition}[the $\kod{I}_0^*$ identity]\label{prop:rings:i0star-identity}
Let $q\ne0$ and $a$ be complex numbers and $\alpha,\beta$ the roots of $X^2-aX+q$, so $\alpha+\beta=a$ and $\alpha\beta=q$. Then
\[
\Bigl(1-\frac\alpha{\beta q}\Bigr)\Bigl(1-\frac1q\Bigr)\Bigl(1-\frac\beta{\alpha q}\Bigr)=\frac{q-1}q\cdot\frac{(q+1)^2-a^2}{q^2},
\]
and $(q+1)^2-a^2=(q+1-a)(q+1+a)$. Equivalently, $(1-qX)(1-(a^2-2q)X+q^2X^2)$ at $X=q^{-2}$ equals the right side. If $q>0$ and $a$ is real with $a^2\le4q$, then
$(q+1)^2-a^2\ge(q-1)^2$, which is $>0$ for $q\ne1$.
\end{proposition}

\begin{proof}
$\frac\alpha\beta+\frac\beta\alpha=\frac{(\alpha+\beta)^2-2\alpha\beta}{\alpha\beta}=\frac{a^2-2q}q$ and $\frac\alpha\beta\cdot\frac\beta\alpha=1$, so
\[
\Bigl(1-\frac\alpha{\beta q}\Bigr)\Bigl(1-\frac\beta{\alpha q}\Bigr)=1-\frac{a^2-2q}{q^2}+\frac1{q^2}=\frac{q^2+2q+1-a^2}{q^2}=\frac{(q+1)^2-a^2}{q^2}.
\]
Multiply by $1-\frac1q=\frac{q-1}q$. The polynomial form: $1-(a^2-2q)q^{-2}+q^2q^{-4}$ is the same middle factor, and $1-q\cdot q^{-2}=1-q^{-1}$. Finally
$(q+1)^2-a^2\ge(q+1)^2-4q=(q-1)^2$.
\end{proof}

For an elliptic curve $E'$ over $\F_q$ with $a=q+1-\#E'(\F_q)$ and $\abs a\le2\sqrt q$ (Hasse's bound, \Cref{ch:ellcurves}), the two factors $q+1-a$ and $q+1+a$
are $\#E'(\F_q)$ and the number of points of its quadratic twist over $\F_q$; so $(q+1)^2-a^2$ is a nonzero integer, and the last valuation in
\Cref{lem:rings:ell-integral}(iii) is finite.

\begin{proposition}[the paper's local terms]\label{prop:rings:local-terms}
Let $q\ge5$ be a prime, $\ell\ne q$ an odd prime, and $a\in\Z$ with $a^2\le4q$. Put
\begin{gather*}
e_{\mathrm{II}}=v_\ell(q-\chi_{-3}(q)),\qquad e_{\mathrm{III}}=v_\ell(q-\chi_{-4}(q)),\\
e_{\mathrm{I}^*}=v_\ell(q^2-1),\qquad e_{\mathrm{I}_0^*}(a)=v_\ell(q-1)+v_\ell\bigl((q+1)^2-a^2\bigr).
\end{gather*}
These are finite nonnegative integers, equal to the valuations of $1-\chi_{-3}(q)q^{-1}$, $1-\chi_{-4}(q)q^{-1}$, $1-q^{-2}$ and of the product in
\Cref{prop:rings:i0star-identity}, and:
\begin{enumerate}[label=(\roman*)]
\item $e_{\mathrm{II}}=v_\ell(q-1)$ if $q\equiv1\pmod3$ and $v_\ell(q+1)$ if $q\equiv2\pmod3$. For $\ell=3$, $e_{\mathrm{II}}\ge1$ always, and it is unbounded ($q=53$: $v_3(54)=3$;
$q=163$: $v_3(162)=4$; these are the least primes with $e_{\mathrm{II}}\ge3$, resp. $\ge4$). For $\ell\ge5$, $e_{\mathrm{II}}\ge1$ iff $q\equiv\chi_{-3}(q)\pmod\ell$.
\item For $\ell=3$, $e_{\mathrm{III}}\ge1$ iff $q\equiv\pm1\pmod{12}$; for $\ell\ge5$, iff $q\equiv\chi_{-4}(q)\pmod\ell$.
\item For $\ell=3$, $e_{\mathrm{I}^*}\ge1$ always; for $\ell\ge5$, $e_{\mathrm{I}^*}\ge1$ iff $q\equiv\pm1\pmod\ell$.
\item For $\ell=3$, $e_{\mathrm{I}_0^*}(a)\ge1$ iff $q\equiv1\pmod3$ or $3\mid a$; for $\ell\ge5$, iff $q\equiv1\pmod\ell$ or $a\equiv\pm(q+1)\pmod\ell$.
\item (Type $\kod{I}_n$.) For $n\ge1$ and $\bar u\in\F_q^\times$, the Tamagawa refinement $v_\ell(n)+\min(v_\ell(n),v_\ell(q-1),v_\ell(c(\bar u)))$ of Lemma lem:tamfactor exceeds
$t=v_\ell(n)$ iff $\ell\mid n$, $q\equiv1\pmod\ell$ and $\bar u\in(\F_q^\times)^\ell$.
\end{enumerate}
\end{proposition}

\begin{proof}
Finiteness: $q-\chi(q)\ge q-1>0$, $q^2-1>0$, and $(q+1)^2-a^2\ge(q-1)^2>0$ by \Cref{prop:rings:i0star-identity}. The identification with valuations of
rationals is \Cref{lem:rings:ell-integral}(iii). (i) By \Cref{prop:rings:paper-characters}, $\chi_{-3}(q)=1$ iff $q\equiv1\pmod3$, and $q\equiv\chi_{-3}(q)\pmod3$. The least primes:
$e_{\mathrm{II}}\ge3$ means $q\equiv\pm1\pmod{27}$, and the candidates $26,28,53$ give the prime $53$ (and $53\equiv2\pmod3$, so $e_{\mathrm{II}}=v_3(54)=3$); $e_{\mathrm{II}}\ge4$ means $q\equiv\pm1\pmod{81}$,
and $80,82,161=7\cdot23$ are not prime, while $163$ is, with $163\equiv1\pmod3$. (ii) \Cref{prop:rings:paper-characters}(iv), and the definition for $\ell\ge5$.
(iii) $q^2\equiv1\pmod3$ for $q\ne3$; for $\ell\ge5$, $\ell\mid(q-1)(q+1)$ iff $\ell$ divides a factor. (iv) $e_{\mathrm{I}_0^*}(a)\ge1$ iff $\ell\mid q-1$ or $\ell\mid(q+1-a)(q+1+a)$. For $\ell=3$
and $q\equiv2\pmod3$, $q+1\equiv0$, so the second condition is $3\mid a^2$. (v) is \Cref{prop:rings:class-c}(iv). (Parts (i) to (iii) were also checked for all
primes $5\le q<10^5$.)
\end{proof}

\begin{example}[the paper's numbers]\label{ex:rings:paper-values}
The curve data below (Kodaira types, modular degrees, Euler factors at $2$) were recomputed with PARI and Cremona's tables; the
arithmetic is that of this section.
\begin{enumerate}[label=(\alph*)]
\item $\curve{121a1}$, type $\kod{II}$ at $11$: $\chi_{-3}(11)=-1$, so $e_3=v_3(11+1)=v_3(12)=1$; $\degphi=6$ and $v_3(6)=1$: the bound is attained (Example ex:qual(b)).
\item $\curve{539c1}$, type $\kod{I}_3^*$ at $7$ (and $\kod{I}_2$ at $11$): $e_3=v_3(7^2-1)=v_3(48)=1$, $t_3=v_3(3)=1$; $\degphi=288=2^5\cdot3^2$, and $v_3=2=e_3+t_3$.
\item $\curve{980i1}$, type $\kod{I}_5^*$ at $7$, $\ell=5$: $e_5=v_5(48)=0$, $t_5=v_5(5)=1$; $\degphi=2880=2^6\cdot3^2\cdot5$. (Here $5\mid N$, so the example illustrates the
congruence, not the theorem.)
\item $\curve{726a1}$, type $\kod{III}$ at $11$: $\chi_{-4}(11)=-1$ and $e_3=v_3(12)=1$; $11\equiv11\pmod{12}$, as \Cref{prop:rings:paper-characters}(iv) requires.
\item At $q=2$ the Kodaira type does not determine $P_2$, and $e_3=v_3(P_2(\frac14))$ (the denominator $1-a_2^2/4$ is $1$, as $a_2=0$ at an additive
prime). $\curve{52a1}$ (type $\kod{IV}^*$ at $2$) and $\curve{336f1}$ (type $\kod{II}$ at $2$) have $P_2=1+2X$, so $e_3=v_3(\frac32)=1$; $\degphi=3$, resp. $24$. $\curve{1232l1}$
(type $\kod{II}^*$ at $2$) has $P_2=1+2X-4X^2-8X^3=(1-2X)(1+2X)^2$, so $P_2(\frac14)=\frac12\cdot\frac94=\frac98$ and $e_3=2$, with $\degphi=288$; and $\curve{592c1}$ (type
$\kod{II}^*$ at $2$) has $P_2=(1-2X)(1+4X^2)$, $P_2(\frac14)=\frac58$, $e_5=1$, $\degphi=80$. The two cubic factors are the polynomial of
\Cref{prop:rings:i0star-identity} with $q=2$ and $a=0$, resp. $a=-2$: $\frac12\cdot\frac{9-a^2}4$ is $\frac98$, resp. $\frac58$.
\item $\curve{539d1}$, type $\kod{I}_0^*$ at $7$; its twist by $\chi_{-7}$ is $E'=\curve{11a3}$, with good reduction at $7$ and $a_7(E')=-2$ (so $\#E'(\F_7)=10$). Then
$e_3=v_3(6)+v_3(64-4)=1+1=2$, and $\degphi=72=2^3\cdot3^2$.
\item $\curve{11a1}$ has type $\kod{I}_5$ at $11$ and $e_3=0$ there; its twist by $\chi_{-11}$ is $\curve{121d2}$, of type $\kod{I}_5^*$ at $11$, with $e_3=v_3(11^2-1)=v_3(120)=1$. So
$e_\ell(E,q)$ is not invariant under a quadratic twist ramified at $q$, as Definition def:terms (revised) says.
\end{enumerate}
\end{example}

\begin{paperlink}\label{pl:rings:terms}
Definition def:terms:
\[
\text{"}e_\ell(E,q)=v_\ell\bigl(L_q(A_f,1)^{-1}/L_q^{\nv}(A_f,1)^{-1}\bigr)=v_\ell\bigl(P_q(q^{-2})/(1-a_q(E)^2q^{-2})\bigr)\text{".}
\]
By \Cref{lem:rings:ell-integral}(ii)
both numerator and denominator lie in $\Z_{(\ell)}$, and the valuation of the quotient is the difference of their valuations. At a multiplicative prime
$a_q^2=1$ and $P_q(X)=1-X$, so the quotient is $1$ and $e_\ell=0$; at an additive prime $a_q=0$ and the denominator is $1$. The proof of Lemma
lem:euler then reads "$e_\ell=v_\ell(1-q^{-2})=v_\ell(q^2-1)$" (type $\kod{I}_n^*$) and "$e_\ell=v_\ell(1-\chi_{-3}(q)q^{-1})=v_\ell(q-\chi_{-3}(q))$" (types $\kod{II}$ to
$\kod{IV}^*$): \Cref{lem:rings:ell-integral}(iii); its displayed identity for type $\kod{I}_0^*$ is \Cref{prop:rings:i0star-identity}, and the same polynomial
is shape (ii) of Proposition prop:shapes. Remark rem:wild ("$P_2(X)=1+2X$ and $e_3=1$") is \Cref{lem:rings:ell-integral}(iv). Remark
rem:afterquant(1) and the abstract's description of the terms rest on \Cref{prop:rings:local-terms}: under the hypotheses of Theorem thm:quant
($3\nmid2N$, $E[3]$ irreducible) a single prime of type $\kod{I}_n^*$ forces $3\mid\degphi$, because $e_3=v_3(q^2-1)\ge1$ for every $q\ge5$, and so does a prime of
type $\kod{I}_0^*$ with $q\equiv1\pmod3$, because then $v_3(q-1)\ge1$. That the table is right, i.e. that these
expressions are the $e_\ell(E,q)$, is Lemma lem:euler, proved in \Cref{ch:lfunctions,ch:paper}.
\end{paperlink}

\section{Exercises}\label{sec:rings:exercises}

\begin{exercise}[a toy Hecke algebra]\label{ex:rings:toy-hecke}
Let $r\ge1$ and $R=\{(a,b)\in\Z^2:a\equiv b\pmod r\}$, a subring of $\Z\times\Z$; let $\phi(a,b)=a$, $I=\ker\phi$ and $J=\{(a,0)\in R\}$. Show that $I\cap J=0$, that
$R/(I+J)\simeq\Z/\phi(J)=\Z/r$, and that the cokernel of $R\to R/I\times R/J$ is $\Z/r$ (\Cref{thm:rings:isomorphism,prop:rings:two-ideals}). Interpret $R$ as the
Hecke algebra of two "eigenforms" whose eigenvalues are congruent modulo $r$, and $r$ as their congruence number, as in Lemma lem:crit.
\end{exercise}

\begin{exercise}\label{ex:rings:residue}
Show that a ring homomorphism $\Z/r\Z\to k$, $k$ a field of characteristic $\ell$, exists iff $\ell\mid r$. Deduce the last step of Lemma lem:crit: if
$I+J$ is contained in a maximal ideal of $\T$ of residue characteristic $\ell$, then $\ell$ divides $\#\T/(I+J)$ (\Cref{prop:rings:residue-char}).
\end{exercise}

\begin{exercise}\label{ex:rings:eisenstein-integers}
In $\Z[\zeta_3]$, check $5+\zeta_3=(2-\zeta_3)(2+\zeta_3)$ and $7=(3+\zeta_3)(2-\zeta_3)$ (note $3+\zeta_3^2=2-\zeta_3$), and compute $\gcd(7,5+\zeta_3)$ by the Euclidean algorithm
(answer: $2-\zeta_3$, of norm $7$, up to units). Verify by listing the solutions of $a^2-ab+b^2=1$ that $\Z[\zeta_3]^\times=\mu_6$.
\end{exercise}

\begin{exercise}\label{ex:rings:check-example}
Check the values in Example ex:qual of the paper: $e_3(\curve{121a1},11)=v_3(12)=1$; $e_3(\curve{539c1},7)=v_3(48)=1$ and $t_3(\curve{539c1},7)=1$; $e_5(\curve{980i1},7)=0$
and $t_5(\curve{980i1},7)=1$. Compare with $v_3(6)$, $v_3(2^5\cdot3^2)$ and $v_5(2^6\cdot3^2\cdot5)$ (\Cref{ex:rings:paper-values}).
\end{exercise}

\begin{exercise}\label{ex:rings:forms}
Check the claim in the proof of Proposition prop:forms that $q\equiv\chi_{-3}(q)\pmod3$ for every prime $q\ne3$, so that $e_3(E,q)\ge1$ at every prime $q\ge5$ of
type $\kod{II}$, $\kod{II}^*$, $\kod{IV}$, $\kod{IV}^*$. Find the least prime $q$ with $v_3(q-\chi_{-3}(q))\ge3$ and with $v_3(q-\chi_{-3}(q))\ge4$ (answers $53$ and $163$).
\end{exercise}

\begin{exercise}\label{ex:rings:type-three}
Check the statement of Section sec:anomaly that the type $\kod{III}$ primes $q\ge5$ with $e_3>0$ are those with $q\equiv1$ or $11\pmod{12}$. For $\ell=5$, show that the type
$\kod{III}$ primes with $e_5>0$ are those with $q\equiv\pm1\pmod{20}$.
\end{exercise}

\begin{exercise}\label{ex:rings:i0star}
Prove \Cref{prop:rings:i0star-identity} by expanding the product. Then check part (1) of Remark rem:afterquant: every prime $q\ge5$ of type $\kod{I}_n^*$ has $e_3\ge1$, and a prime of
type $\kod{I}_0^*$ has $e_3\ge1$ iff $q\equiv1\pmod3$ or $3\mid a_q(E')$. Compute $e_3$ for type $\kod{I}_0^*$ at $q=7$ with $a_7(E')=-2$, and at $q=5$ and at $q=11$ with $a_q(E')=0$
(answers: $2$, $2$, $2$).
\end{exercise}

\begin{exercise}\label{ex:rings:wild}
Check Remark rem:wild: with $P_2(X)=1+2X$, $e_3=v_3(P_2(1/4))=1$; with $P_2=1$, $e_3=0$. Explain why the denominator $1-a_q^2q^{-2}$ of Definition def:terms is $1$ at
an additive prime and $1-q^{-2}=P_q(q^{-2})$ at a multiplicative one, so that $e_\ell=0$ there.
\end{exercise}

\begin{exercise}\label{ex:rings:tamfactor}
For $q=7,13,19,37$ and $\ell=3$, list the cubes and the ninth powers in $\F_q^\times$, and compute $v_3(c(u))$ for $u=2,3,-1$ (answers: $0,0,1$ for $q=7$ and $13$; $0,0,2$ for
$q=19$ and $37$). Check the statement of Lemma lem:tamfactor that $\min(v_\ell(n),v_\ell(q-1),v_\ell(c))>0$ iff $q\equiv1\pmod\ell$, $\ell\mid n$ and $u$ is an $\ell$-th power.
\end{exercise}

\begin{exercise}\label{ex:rings:power-index}
Prove \Cref{thm:rings:power-index}(iii) directly from a generator of $\F_q^\times$, without quoting \Cref{ch:groups}. Show that the $\ell$-th power map of $\F_q^\times$ is
bijective iff $q\not\equiv1\pmod\ell$. Find the least $f$ with $\mu_9\subset\F_{2^f}$ (answer $6$) and show $v_3(2^{6m}-1)=2+v_3(m)$ (\Cref{lem:rings:lte}).
\end{exercise}

\begin{exercise}\label{ex:rings:h0}
Show that the only solutions of $a^2\equiv1\pmod{\ell^k}$ ($\ell$ odd) are $a\equiv\pm1$, and deduce that a homomorphism $\chi$ onto $(\Z/\ell^k)^\times$ with $\chi^2=1$ forces
$\ell^k=3$. Explain how this gives $h_\ell(E)=0$ for $\ell\ge5$ and $h_3(E)\le1$ in Lemma lem:h0 (the Galois side is \Cref{ch:galois}).
\end{exercise}

\begin{exercise}\label{ex:rings:f16}
Construct $\F_{16}=\F_2[x]/(x^4+x+1)$ (check irreducibility: no root, and $x^4+x+1\ne(x^2+x+1)^2$, the only irreducible quadratic being $x^2+x+1$). Show
that $\bar x$ generates $\F_{16}^\times$ ($\bar x^3\ne1$ and $\bar x^5=\bar x^2+\bar x\ne1$). List the subfields, and show that the Frobenius has orbits of sizes $1,1,2,4,4,4$ on
$\F_{16}$; deduce $N_4=3$ over $\F_2$.
\end{exercise}

\begin{exercise}[Eisenstein's proof of reciprocity]\label{ex:rings:eisenstein-lemma}
Let $p\ne q$ be odd primes. Prove Eisenstein's lemma $\bigl(\frac qp\bigr)=(-1)^{\sum_{a=1}^{(p-1)/2}\lfloor2qa/p\rfloor}$ and show that the exponent has the parity of
$\sum_{u=1}^{(p-1)/2}\lfloor qu/p\rfloor$. By counting the lattice points in the rectangle $0<x<p/2$, $0<y<q/2$ on either side of the diagonal, show
$\sum_{u=1}^{(p-1)/2}\lfloor qu/p\rfloor+\sum_{v=1}^{(q-1)/2}\lfloor pv/q\rfloor=\frac{p-1}2\cdot\frac{q-1}2$, and deduce \Cref{thm:rings:qr}(i). (Both identities were checked for all
odd primes below $200$.)
\end{exercise}

\begin{exercise}\label{ex:rings:kronecker-table}
Verify that $\bigl(\frac D\cdot\bigr)$ is periodic modulo $\abs D$ for $D=-3,-4,-7,5,8,12$. Tabulate $\chi_{-3},\chi_{-4},\chi_{-7},\chi_{-11},\chi_5$ at the primes below $30$. Show
that on integers prime to $6$ the function $\bigl(\frac3\cdot\bigr)$ agrees with $\chi_{12}$, of conductor $12$, while $\bigl(\frac{-3}\cdot\bigr)=\chi_{-3}$ has conductor $3$: the sign in $q^*$
matters.
\end{exercise}

\begin{exercise}\label{ex:rings:gauss-sums}
Compute $\tau(\chi_{-4})=2i$, $\tau(\chi_{-3})=\sqrt{-3}$ and $\tau\bigl(\bigl(\frac\cdot5\bigr)\bigr)=\sqrt5$ (use $\cos(2\pi/5)=(\sqrt5-1)/4$). Show that $\chi_{-3}$ regarded modulo $6$ has
Gauss sum $\sqrt{-3}$ and that the principal character modulo $n$ has $\tau(\chi_0)=\mu(n)$.
\end{exercise}

\begin{exercise}\label{ex:rings:x4}
Show that $x^4+1$ is irreducible over $\Q$ but reducible modulo every prime (\Cref{warn:rings:x4plus1}). Check the claim of Proposition prop:ns for
$\curve{11a1}$, whose mod $3$ image is $\GL_2(\F_3)$: its $3$-division polynomial $3x^4-4x^3-60x^2-237x-21$ is irreducible (\Cref{ex:rings:psi3}).
\end{exercise}

\begin{exercise}\label{ex:rings:cyclotomic-primes}
Show that if $p\nmid n$ and $p\mid\Phi_n(a)$, then $p\equiv1\pmod n$, and deduce that there are infinitely many primes $\equiv1\pmod n$
(\Cref{prop:rings:cyclotomic-mod-p}). Check with $\Phi_{12}(2)=13$ and $\Phi_9(2)=73$.
\end{exercise}

\section{Notes and references}\label{sec:rings:notes}

Sections 1 to 4, 6 and 8 follow Dummit and Foote \cite{DummitFoote}: Sections 7.1 to 7.6 (rings, ideals, quotients, fractions, the Chinese
remainder theorem), 8.1 to 8.3 (Euclidean domains, PIDs, UFDs), 9.1 to 9.5 (polynomial rings, Gauss's lemma, Eisenstein's criterion),
13.1, 13.4 to 13.6 (field extensions, splitting fields, separability, cyclotomic polynomials) and 14.3 (finite fields); Lang
\cite[Chapter IV, Section 2, and Chapter V, Section 5]{Lang} gives Gauss's lemma and finite fields in greater generality, and Atiyah and
Macdonald \cite[Chapter 1]{AtiyahMacdonald} the existence of maximal ideals and the ideal theory of products. The Gaussian and
Eisenstein integers are treated in Ireland and Rosen \cite[Chapter 1]{IrelandRosen}, together with unique factorisation in $\Z$.

Sections 5 to 10 are classical elementary number theory: Ireland and Rosen \cite{IrelandRosen}, Chapter 4 for $(\Z/n)^\times$ and primitive
roots, Chapter 5 for quadratic reciprocity and the Jacobi symbol, Chapter 6 for the quadratic Gauss sum (Section 6.4 for Gauss's
determination of its sign, \Cref{rem:rings:gauss-sign}), Chapter 7 for finite fields, Chapter 8 for Gauss sums of general characters;
Serre \cite[Chapter I]{SerreCourse} proves reciprocity by the same computation in characteristic $p$ as \Cref{thm:rings:qr} and the second
supplement with an eighth root of unity. Among the many other proofs of reciprocity, Eisenstein's lattice-point proof is
\Cref{ex:rings:eisenstein-lemma}, and the proof through the splitting of $p$ in $\Q(\zeta_q)$ is given in \Cref{ch:ant}. The lifting-the-exponent
lemma is folklore; the form used here, with its proof by the binomial theorem, is the standard one in the proof that $(\Z/p^k)^\times$
is cyclic.

The Kronecker symbol follows Cohen \cite[Section 1.4]{Cohen93}, and fundamental discriminants Cohen \cite[Section 5.1]{Cohen93}; Dirichlet characters, conductors and
primitivity follow Washington \cite[Chapter 3]{Washington}, and orthogonality is Serre \cite[Chapter VI, Section 1]{SerreCourse}. The formula
for the Gauss sum of an imprimitive character (\Cref{thm:rings:gauss-sums}(iv)) is proved here by the Chinese remainder theorem and the
Ramanujan sum $\sum_{(y,m)=1}\zeta_m^y=\mu(m)$. Dirichlet characters as characters of $G_\Q$ are \Cref{ch:ant}, twisting of modular forms by them
\Cref{ch:hecke}.

Valuations and the $p$-adic absolute value are in Serre \cite[Chapter II]{SerreCourse} and Neukirch \cite[Chapter II]{Neukirch}, where
Ostrowski's theorem (\Cref{thm:rings:ostrowski}) is proved in Section 3. The arithmetic of the paper's local terms
(\Cref{lem:rings:ell-integral,prop:rings:i0star-identity,prop:rings:local-terms}) is elementary; the identification of these expressions with the
Euler factors of the adjoint $L$-function is Lemma lem:euler of the paper, treated in \Cref{ch:lfunctions,ch:paper}. Curve data were computed
with PARI/GP \cite{PARI} and Cremona's tables \cite{Cremona97}.

The following are proved in full here although a statement with a reference would have sufficed: the existence of
maximal ideals (from Zorn's lemma), the unique factorisation of $R[x]$ over a UFD, the primitivity of $\chi_D$ for fundamental $D$
(\Cref{cor:rings:fundamental-primitive}) and the Gauss sum of an imprimitive character. Only Ostrowski's theorem and the sign of the
quadratic Gauss sum are stated without proof, and neither is used. No theorem numbers of the sources are cited. The chapter and section numbers above are given as pointers to the relevant part of
each book; they were not verified against the printed editions.
